\section{Introduction}

\subsection{Purpose of the Document}
This document report the requirements for \acs{ard} product generation starting from the state of the art analysis D1~\cite{D1}, the use case defined in D2~\cite{D2}, the definition of the algorithms provided in D4~\cite{D4}, and the performance analysis provided in D5~\cite{D5}. 

\subsection{Reference Documents}

\renewcommand{\refname}{}
\vspace{-2em}
\bibliographystyle{apalike}
\bibliography{strings, references_paco}

\clearpage

\section{Document Rationale and General Considerations}
This product requirements specification document starts from the input State of the art Analysis, and reports the specifications for the following products:
    \begin{itemize}
        \item Co-registered \acs{slc} (\acs{cslc}) and geocoded equivalent (\acs{gslc})
        \item Normalized radar backscatter (\acs{nrb})
        \item Polarimetric covariance matrix (\acs{cov}) and geocoded equivalent (\acs{gcov})
        \item Ocean radar backscatter (\acs{orb})
        \item Precise Range Azimuth Offsets (\acs{pro} \& \acs{pao}) and geocoded equivalent (\acs{gpro} \& \acs{gpao})
        \item Interferometric products (\acs{int} / \acs{cohe} / \acs{unw}) and geocoded equivalent (\acs{gint} / \acs{gcohe} / \acs{gunw})
        \item Co-registered Multilooked Complex (\acs{cmlc}) and geocoded equivalent (\acs{gmlc})
        \item Enhanced \acs{nrb} (\acs{enrb}) \& \acs{orb} (\acs{eorb})
        \item Geo-projected Ground Motion (\acs{ggm})
        \item Stripmap equivalent \acs{cslc}
    \end{itemize}

The study aims to rationalise \acs{ard} products requirements specification by identifying common characteristics and performances relevant to ROSE-L and its joint usage with Sentinel-1 and other missions. Particular attention will be paid to identifying:
    \begin{itemize}
        \item Inputs and auxiliary data needed (e.g. \acs{dem}, \acs{etad}, models of the atmosphere)
    	\item Main variables and metadata
        \item Accuracy
        \item Timeliness       
        \item Usability, Integrability and Interoperability
    \end{itemize}

The generation of some \acs{ard} products requires the use of external data or model or, alternatively, a data-driven approach can be used. In general, the utilisation of external data could streamline processing but may lead to delays in product delivery (e.g., precise orbits, weather models based on ERA5 (\acs{ecmwf} reanalysis data), or \acs{etad} require a few weeks for delivery). Conversely, data-driven processing makes product generation independent of external data but necessitates more complex processing of image stacks to estimate and remove ionospheric effects, improve co-registration accuracy, etc.

In the following sections, the various aspects described above will be analysed for the set of identified \acs{ard} products.

It is important to highlight that only the main variables and metadata and characteristics of the proposed products have been included: those are the ones useful for describing each product and supporting the evaluation process for selecting the two \acs{ard} products to be developed in Phase 2 of the project. A detailed product specification will be prepared during this second phase, in coordination with \acs{esa} and based on the guidance provided by the Working Group, also taking into account the findings of the state-of-the-art review and of all the analyses conducted for the selected products.

\section{Co-registered and geocoded SLC (CSLC \& GSLC)}
\label{sec:cslc_gslc}

\subsection{Product Description}
\label{sec:cslc_product_description}
Image co-registration refers to the process of aligning two or more images so that each pixel corresponds to the same geographic location in the target scene. This alignment can achieve subpixel accuracy and it is a critical step in generating multidimensional image stacks, enabling advanced data analysis.

For the ROSE-L mission, the co-registered products are provided in radar geometry \cite{RD41}. This product, referred to as the \acs{cslc} (Co-Registered Single-Look Complex), plays a central role in the interferometric process. It incorporates corrections for geometric distortions (e.g., topographic effects), atmospheric delays (tropospheric and ionospheric), and other influencing factors. 

The \acs{cslc} maintains the full fidelity of the original SAR data, introducing no distortions, and serves as a foundational product from which other \acs{ard} products can be derived.
The amplitude of the \acs{cslc} corresponds to \gls{beta0}, obtained by the calibration step using the Look-up-Table from the \acs{slc} annotations.
The phase information in a \acs{cslc} is flattened with respect to the reference orbit and a Digital Elevation Model (DEM).

To support analysis in map-projected (ground) geometry, \acs{slc} or \acs{cslc} products can be automatically resampled onto a common geographic grid (e.g., \acs{utm}, \acs{ups}). This transformation removes the topographic phase and mitigates artifacts related to atmospheric disturbances or orbital inaccuracies. This map-projected product is known as the \acs{gslc} (Geocoded Single-Look Complex). While \acs{gslc} simplifies user interaction, allowing direct visualisation in \acs{gis} tools and the integration of data from multiple satellites, it is inherently less rigorous and unsuitable for high-precision interferometric applications such as coherence analysis or phase unwrapping.
Like the \acs{slc}, the phase of the \acs{gslc} is flattened with respect to the reference orbit and DEM in order to eliminate topographic contributions \cite{Zebker2017, Zheng2017}. The spatial resolution of the \acs{gslc} in map coordinates is comparable to the native resolution of the original \acs{slc}.

The \acs{cslc} product consists of \acs{slc} data that has been resampled to a reference radar grid, which may be:
\begin{itemize}
	\item virtual: based on a nominal reference orbit and sensor parameters,
	\item real: using the radar grid of a specific \acs{slc} image as the reference.
\end{itemize}

Although co-registration is essential for generating most \acs{ard} products and enabling many applications, the \acs{cslc} is often considered an intermediate product in the processing chain. However, it can also serve as a final \acs{ard} product for users who lack the expertise or the resources to perform advanced co-registration and topographic phase removal, but still wish to produce more sophisticated products than a basic \acs{ard}.

The \acs{cslc} is additionally foreseen as a standard \acs{ard} product in the ongoing re-engineering of Sentinel-1 Level-1 and Level-2 processors \cite{S1_RP_portal, S1_RP_ProdSpec}.

This planned product aligns with the evolution of the Sentinel-1 processor chain, ensuring consistency across Copernicus missions as required by the re-engineering project.

The dataset is formatted as a ZARR file. This format supports the inclusion of both the processed measurement data and the associated metadata within the same file, making the product self-contained and self-describing, with no need for additional annotations or supporting files.

\subsubsection{Input and Auxiliary Data}
\label{sec:aux_data_cslc}
The following input and auxiliary data are considered for product generation:

\begin{longtblr}[
	label = tab:in_aux_cslc_products,
	entry = none,
	caption = {\acs{cslc}/\acs{gslc} input and auxiliary data},
	]{
		width = \linewidth,
		colspec = {Q[l,20]Q[l,60]Q[c,20]},
		hlines,
		hline{1-Z} = {-}{0.05em},
	}
	\textbf{Name}   & \textbf{Description}   & \textbf{Mandatory (M) / Optional (O)} \\
	\acs{slc} & L1 Single Look Complex Images and annotations & M \\
    Radar Reference Grid & Virtual or real radar reference grid & M \\
    \acs{dem} & External DEM (e.g., Copernicus DEM at different resolutions, with 30 m as the default worldwide). & M \\
    Orbits & Standard, restituted or precise orbits. & M \\
    \acs{etad} & \acs{etad}-like products for atmospheric delay and timing corrections. & O \\
    \acs{tec} & External \acs{tec} data for delay correction & O \\
\end{longtblr}

\subsubsection{Main Variables}
\label{sec:variables_cslc}
The following tables list the main variables associated with the \acs{cslc} and \acs{gslc} products, respectively.

\begin{longtblr}[
	label = tab:var_cslc_products,
	entry = none,
	caption = {\acs{cslc} main variables},
	]{
		width = \linewidth,
		colspec = {Q[l,15]Q[l,52]Q[l,20]Q[l,13]},
		hlines,
		hline{1-Z} = {-}{0.05em},
	}
	\textbf{Name} & \textbf{Description} & \textbf{Dimensions} & \textbf{Data Type} \\
	cslc & Co-registered single look complex data: 
    
    - the amplitude is in \gls{beta0} convention, obtained by the calibration step using the Look-up-Table from the \acs{slc} annotations;
    
    - the phase corresponds to the interpolated and resampled \acs{slc} flattened phase with respect to the reference orbit and the input Digital Elevation Model (DEM) & polarisation, acquisition\_date, azimuth\_time, range\_time & complex64 \\
    phase\_model & Topographic phase model, in rad, removed from the original image phase & acquisition\_date, azimuth\_time, range\_time & float32 \\
    
\end{longtblr}

\begin{longtblr}[
	label = tab:var_gslc_products,
	entry = none,
	caption = {\acs{gslc} main variables},
	]{
		width = \linewidth,
		colspec = {Q[l,15]Q[l,52]Q[l,20]Q[l,13]},
		hlines,
		hline{1-Z} = {-}{0.05em},
	}
	\textbf{Name} & \textbf{Description} & \textbf{Dimensions} & \textbf{Data Type} \\
	gslc & Projected \acs{slc} or \acs{cslc}, with module equal by default to \gls{beta0} or, optionally to \gls{gamma0}, coefficient and phase equal to the \acs{slc} or \acs{cslc} & polarisation, acquisition\_date, x, y & complex32 \\
    phase\_model & Projected topographic phase model, in rad, removed from the original image phase & acquisition\_date, azimuth\_time, range\_time & float32 \\
    local\_inc
    
    \_angle & DEM-based local Incidence angle image [deg]. This variable is present only if the \acs{gslc} is provided in the \gls{gamma0} convention and is generated from an \acs{slc}. If the \acs{gslc} is derived from a \acs{cslc}, this variable will be null, and the reference\_inc\_angle metadata will be present & x, y & float32 \\
    anf & Area normalisation factor or local contributing area, used to obtain the gamma0 coefficient. This variable is present only if the \acs{gslc} is provided in the \gls{gamma0} convention and is generated from a \acs{slc}. If the \acs{gslc} is derived from a \acs{cslc}, this variable will be null, and the reference\_anf metadata will be present & x, y & float32 \\
\end{longtblr}

Note that, for the \acs{gslc} the fields "anf" and "local\_inc\_angle" are always included: as variables for products directly derived from \acs{slc}, and as metadata for products derived from \acs{cslc}, with links to the corresponding static layer.

The table below describes the dimensions/coordinates associated with the \acs{cslc} and \acs{gslc} main variables.

\begin{longtblr}[
	label = tab:var_cslc_coordinates,
	entry = none,
	caption = {\acs{cslc} and \acs{gslc} main variables coordinates/dimensions},
	]{
		width = \linewidth,
		colspec = {Q[l,20]Q[l,62]Q[l,18]},
		hlines,
		hline{1-Z} = {-}{0.05em},
	}
	\textbf{Name} & \textbf{Description} & \textbf{Data Type} \\
	polarisation & Acquired polarisations & string \\
    acquisition\_date & Image acquisition dates, formatted as YYYY-MM-DD & string \\
    azimuth\_time & Time along satellite track direction & datetime64 \\
    range\_time & Time across satellite track direction & float64 \\
    x & UTM/UPS x (East) coordinates [m] & float64 \\
    y & UTM/UPS y (North) coordinates [m] & float64 \\
\end{longtblr}

\subsubsection{Main Metadata}
\label{sec:metadata_cslc}
The following tables list the main metadata and per-pixel metadata associated with the \acs{cslc} and \acs{gslc} products, respectively.

\begin{longtblr}[
	label = tab:meta_cslc_products,
	entry = none,
	caption = {\acs{cslc} main metadata and per-pixel metadata},
	]{
		width = \linewidth,
		colspec = {Q[l,17]Q[l,38]Q[l,20]Q[l,13]Q[c,12]},
		hlines,
		hline{1-Z} = {-}{0.05em},
	}
	\textbf{Name} & \textbf{Description} & \textbf{Dimensions} & \textbf{Data Type} & \textbf{Per-pixel Metadata}\\
	reference\_slc & Link to input reference SLC & - & string & \\
    reference\_orbit & Link to reference orbit file & - & string & \\
    orbit\_coeffs & Polynomial coefficients describing the orbit associated with each input image belonging to the \acs{cslc} stack & acquisition\_date, axis, degree & float64 & \\
    orbit\_reference\
    
    \_time & Reference time associated with the orbit polynomial coefficients for each input image belonging to the \acs{cslc} stack & acquisition\_date & float64 & \\
    orbit\_type & Type of orbital data used for each of the images belonging to the \acs{cslc} stack: ETAD, POE, POD, RES & acquisition\_date & string & \\
    orbit\_files & Orbital file path list used for each of the images belonging to the \acs{cslc} stack & acquisition\_date & string & \\
    reference\_dem & Link to reference DEM & - & string & \\
    etad\_file & ETAD file used for each image belonging to the CSLC stack. This is optional: it is included only if used; otherwise, it is null & - & string & \\ 
    etad\_azcorr & Extended Timing Annotation Data: azimuth time correction resampled onto the common grid of the reference geometry. This is optional: it is included only if used; otherwise, it is null & acquisition\_date, azimuth\_time, slant\_range\_time & float64 & \ding{51} \\
    etad\_rgcorr & Extended Timing Annotation Data: slant range time correction resampled onto the common grid of the reference geometry. This is optional: it is included only if used; otherwise, it is null & acquisition\_date, azimuth\_time, slant\_range\_time & float64 & \ding{51} \\
    tec & Link to external \acs{tec} data, if used & - & string & \\
    reference\_inc
    
    \_angle & Link to DEM-based local incidence angle map. & - & string & \\
    reference\_anf & Link to reference area normalisation factor or local contributing area. & - & string & \\
    look\_vector & 3-D components radar unit look vector, in an Earth-Centred Earth-Fixed (\acs{ecef}) coordinate system. It consists of unit vectors from antenna to surface pixel. Only for reference geometry. &  (x, y) or (azimuth\_time, range\_time) & float64 & \ding{51} \\
    phase\_demod & Phase to be applied to demodulate each image belonging to the \acs{cslc} stack & acquisition\_date, azimuth\_time, slant\_range\_time & float32 & \ding{51} \\
    sp\_side-band & Side-band used for the application of the split-spectrum technique. It is present only when the split-spectrum method is applied; otherwise, it is set to null. \\ 
    beta0 & LUT of the beta0 coefficients obtained from each input \acs{slc} annotation belonging to the \acs{cslc} stack & acquisition\_date, azimuth\_time, slant\_range\_time & float32 & \\
    gamma0 & LUT of the gamma0 coefficients obtained from each input \acs{slc} annotation belonging to the \acs{cslc} stack & acquisition\_date, azimuth\_time, slant\_range\_time & float32 & \\
    rcs & LUT of the Radar Cross Section (RCS) coefficients obtained from each input \acs{slc} annotation belonging to the \acs{cslc} stack & acquisition\_date, azimuth\_time, slant\_range\_time & float32 & \\
    noise & LUT of thermal noise correction obtained from each input \acs{slc} annotation belonging to the \acs{cslc} stack & acquisition\_date, azimuth\_time, slant\_range\_time & float32 & \\
    crg\_accuracy & Estimated co-registration accuracy in correspondence of each patch of the images belonging to the \acs{cslc} stack & acquisition\_date, azimuth\_time, slant\_range\_time & float32 & \\
    crg\_accuracy
    
    \_bias & Estimated co-registration accuracy bias, computed as the mean value of the co-registration accuracy over all image patches & acquisition\_date & float32 & \\
    crg\_accuracy
    
    \_std & Estimated co-registration accuracy standard deviation, computed as the standard deviation of the co-registration accuracy over all image patches  & acquisition\_date & float32 & \\    
    source\_data
    
    \_attributes & Dictionary containing the attributes associated with each image belonging to the \acs{cslc} stack & - & Dictionary & \\
    x\_sample & Pixel spacing in azimuth direction & - & float32 & \\
    y\_sample & Pixel spacing in slant range direction & - & float32 & \\
\end{longtblr}

\begin{longtblr}[
	label = tab:meta_gslc_products,
	entry = none,
	caption = {\acs{gslc} main metadata and per-pixel metadata},
	]{
		width = \linewidth,
		colspec = {Q[l,17]Q[l,38]Q[l,20]Q[l,13]Q[c,12]},
		hlines,
		hline{1-Z} = {-}{0.05em},
	}
	\textbf{Name} & \textbf{Description} & \textbf{Dimensions} & \textbf{Data Type} & \textbf{Per-pixel Metadata}\\
	cslc/slc & Link to input \acs{cslc} or \acs{slc} & - & string & \\
    reference\_dem & Link to reference DEM & - & string & \\
    etad\_file & Link to ETAD file. This is optional: it is included only if used; otherwise, it is null & - & string & \\ 
    reference\_inc
    
    \_angle & Link to DEM-based local incidence angle map. It is present only if the \acs{gslc} is generated from a \acs{cslc} & - & string & \\
    reference\_anf & Link to reference area normalisation factor or local contributing area. It is present only if the \acs{gslc} is generated from a \acs{cslc} & - & string & \\
    look\_vector & 3-D components radar unit look vector, in an Earth-Centred Earth-Fixed (\acs{ecef}) coordinate system. It consists of unit vectors from antenna to surface pixel. Only for reference geometry. &  (x, y) or (azimuth\_time, range\_time) & float64 & \ding{51} \\
    mask & Mask containing, for each ground cell, the following information:
    
     - 0: no distortions
    
     - 1: layover
    
     - 2: shadow
    
     - 3: undefined & x, y & byte & \ding{51} \\
    nlooks & Number of equivalent looks for the gamma flattening. Present if the \acs{gslc} amplitude is in \gls{gamma0} convention & x, y & float32 & \ding{51} \\
    geometric
    
    \_accuracy
    & Expected geometric accuracy, including eastern and northern geometric accuracy bias and standard deviation. The measures will be performed on a validation dataset and periodically updated. & - & dictionary & \\ 
    east\_geom
    
    \_bias & Expected eastern geometric bias & - & float32 & \\
    east\_geom
    
    \_std & Expected eastern geometric standard deviation & - & float32 & \\
    north\_geom
    
    \_bias & Expected northern geometric bias & - & float32 & \\
    north\_geom
    
    \_std & Measured northern geometric standard deviation & - & float32 & \\
    geometric
    
    \_accuracy
    
    \_radial
    
    \_rmse	& Expected radial RMSE of geometric accuracy & - & float32 & \\
    x\_sample & Pixel spacing in x direction & - & float32 & \\
    y\_sample & Pixel spacing in y direction & - & float32 & \\
\end{longtblr}

The table below describes the dimensions/coordinates associated with the \acs{cslc} and \acs{gslc} main metadata.

\begin{longtblr}[
	label = tab:var_cslc_metadata_coordinates,
	entry = none,
	caption = {\acs{cslc} and \acs{gslc} main metadata coordinates/dimensions},
	]{
		width = \linewidth,
		colspec = {Q[l,20]Q[l,62]Q[l,18]},
		hlines,
		hline{1-Z} = {-}{0.05em},
	}
	\textbf{Name} & \textbf{Description} & \textbf{Data Type} \\
	acquisition\_date & Image acquisition dates, formatted as YYYY-MM-DD & string \\
    azimuth\_time & Time along satellite track direction & datetime64 \\
    range\_time & Time across satellite track direction & float64 \\
    x & UTM/UPS x (East) coordinates [m] & float64 \\
    y & UTM/UPS y (North) coordinates [m] & float64 \\
    axis & x, y, z axis index & int64 \\
    degree & orbit polynomial coefficient degree & int64 \\
\end{longtblr}


\subsubsection{Reference System Tiling and Slicing}
\label{sec:cslc_reference_system_tiling}
\autoref{tab:clsc_gslc_ref_system_slicing} reports the reference system, the slicing, and grid for \acs{cslc} and \acs{gslc} products.


\begin{longtblr}[
	label=tab:clsc_gslc_ref_system_slicing,
	entry = none,
	caption = {Reference system and Slicing for \acs{cslc} and \acs{gslc}},
	]{
		width=\linewidth,
		colspec = {Q[l,15]Q[l,40]Q[l,15]Q[l,30]},
		%vline{2-5} = {-}{},
		%vline{5-5} = {-}{},
		hlines,
		hline{1-Z} = {-}{0.05em},
		rowhead=2,
	}
	\textbf{Product} & \textbf{Reference System} & \textbf{Slicing} & \textbf{Spacing}\\
	\textbf{\acs{cslc}} & Radar: azimuth/slant-range & Burst & Same of the input data \\
	\textbf{\acs{gslc}} & Cartographic: UTM/UPS & Burst & Comparable with radar spacing \\   
	
\end{longtblr}

The ideal solution to ensure interoperability between the ROSE-L and Sentinel-1 Copernicus missions would be to adopt both a common slicing approach and a shared reference radar grid.
While the former is unfeasible—since slicing is determined by burst segmentation—the latter could still be explored.
Given that both missions share the same orbit, it is possible to define a common virtual radar grid, with spacing determined by the higher-resolution sensor. This would allow data from the different sensors to be linked using the coordinates of the reference virtual grid.


\subsection{Product Requirements}
\label{sec:cslc_requirements}

\subsubsection{Accuracy}

A critical requirement for the \acs{cslc} is that it must preserve the intrinsic radiometric and geometric characteristics of the SAR imagery. To verify this, we propose the use of quantitative image quality indicators before and after co-registration.
	
The first two indicators employed are the Peak Sidelobe Ratio (\acs{pslr}) and the Integrated Sidelobe Ratio (\acs{islr}), as previously defined. These metrics are sensitive to distortions introduced by interpolation and geometric transformations. Therefore, by comparing \acs{pslr} and \acs{islr} values computed on point targets or isolated scatterers in both the original and the co-registered SAR images (\acs{cslc} product), we can assess whether the process has introduced significant degradation.
	
Ideally, the \acs{pslr} and \acs{islr} values should remain stable, within an acceptable tolerance, and significant deviations would indicate that the interpolation step or geometric transformation negatively affects image quality and could compromise downstream applications such as InSAR or coherent change detection. 

The accuracy requirements reported in the following tables are intended to be evaluated on a set of isolated scatterers or Corner Reflectors (\acs{cr}s) with a radar cross section (\acs{rcs}) greater than 35~dB.


\begin{longtblr}[
	label = tab:cslc_general_accuracy,
	entry = none,
	caption = {Accuracy Requirements for \acs{cslc} product},
]{
	width = \linewidth,
	colspec = {Q[l,30] Q[l,60] Q[c,20]},
	hlines,
	hline{1-Z} = {-}{0.05em},
	rowhead = 1,
}

\textbf{Metric} & \textbf{Description} & \textbf{Requirement} \\

PSLR degradation (dB) & Standard deviation (1-$\sigma$) of degradation of peak sidelobe level relative to mainlobe peak in both range and azimuth directions & $\leq 0.1$~dB \\

ISLR degradation (dB) & Standard deviation (1-$\sigma$) of degradation of total sidelobe energy relative to mainlobe energy in both range and azimuth directions & $\leq 0.1$~dB \\

Resolution degradation (\%) & Standard deviation (1-$\sigma$) of resolution decreasing expressed as fraction of original resolution in both range and azimuth directions & $\leq 1\%$ \\


\end{longtblr}


Regarding the co-registration errors in both azimuth and range directions, the achievable performance is strongly influenced by both the sensor's operating band (with significant differences between C-band and L-band) and the type of co-registration strategy employed. In particular, three distinct co-registration approaches can be identified:
\begin{itemize}
    \item Model-driven. In this case, accurate orbital data is assumed to be available, but no models are used to correct for atmospheric delays.
    \item Model-driven + atmospheric correction. In addition to precise orbital information, this approach assumes the availability of state-of-the-art atmospheric correction models (e.g., \acs{etad}, JHR GIM), which are used to mitigate the impact of atmospheric propagation effects.
    \item Data-driven. This approach includes fine co-registration techniques based on cross-correlation methods, as well as azimuth correction strategies such as Enhanced Spectral Diversity (\acs{esd}), which exploits the spectral diversity in the overlap areas of consecutive bursts to estimate residual azimuth co-registration errors.
\end{itemize}
To provide a clear overview of the co-registration accuracy requirements under different processing assumptions, the expected performance thresholds are summarized separately for the C-band and L-band sensors. Due to the distinct propagation characteristics and to the sensitivity to atmospheric effects associated with each frequency band, the accuracy requirements differ accordingly. The following tables report the maximum acceptable residual errors in azimuth and range directions for each band, considering various co-registration strategies, including model driven approaches (with and without atmospheric correction) and data-driven methods. 

All accuracy values are expressed as fractions of the nominal resolution in the corresponding dimension, providing a normalized reference for evaluating the effectiveness of the co-registration process. The accuracy thresholds are evaluated at 1-$\sigma$.


\begin{longtblr}[
    label = tab:cslc_coregistration_accuracy_cband,
    entry = none,
    caption = {Coregistration Accuracy Requirements for \acs{cslc} product (C-Band)},
    ]{
        width = \linewidth,
        colspec = {Q[l,30] Q[l,60] Q[c,20]},
        hlines,
        hline{1-Z} = {-}{0.05em},
        rowhead = 1,
    }
    
    \textbf{Metric} & \textbf{Description} & \textbf{Requirement} \\
    
    \SetCell[r=2]{l} Range Accuracy (\% pixel) 
    & Model-Driven (only orbit and \acs{dem}) & $\leq 5\%$\\
    & Model-Driven + atmospheric correction & $\leq 2\%$ \\

    \SetCell[r=2]{l} Azimuth Accuracy (\% pixel) 
    & Model-Driven (only orbit and \acs{dem}) & $\leq 2\%$\\
    & Data-Driven (\acs{esd}) & $\leq 0.5\%$ \\
    
\end{longtblr}

\begin{longtblr}[
    label = tab:cslc_coregistration_accuracy_lband,
    entry = none,
    caption = {Coregistration Accuracy Requirements for \acs{cslc} product (L-Band)},
    ]{
        width = \linewidth,
        colspec = {Q[l,30] Q[l,60] Q[c,20]},
        hlines,
        hline{1-Z} = {-}{0.05em},
        rowhead = 1,
    }
    
    \textbf{Metric} & \textbf{Description} & \textbf{Requirement} \\
    
    \SetCell[r=3]{l} Range Accuracy (\% pixel) 
    & Model-Driven (only orbit and \acs{dem})& $\leq 60\%$\\
    & Model-Driven + atmospheric correction & $\leq 25\%$ \\
    & Model-Driven + atmospheric correction + Data-Driven (\acs{icc}) & $\leq 10\%$ \\

    \SetCell[r=3]{l} Azimuth Accuracy (\% pixel) 
    & Model-Driven (only orbit and \acs{dem}) & $\leq 150\%$ \\
    & Model-Driven + atmospheric correction & $\leq 50\%$ \\
    & Data-Driven (\acs{icc}) & $\leq 10\%$ \\
    
\end{longtblr}
In the case of L-band SAR systems, atmospheric effects -in particular those related to ionospheric disturbances- can significantly impact the accuracy of fine co-registration. As a result, the achievable co-registration performance is expected to exhibit strong variability depending on several factors, including the geographic location of the acquisition, the time of day, and prevailing space weather conditions.

The co-registration accuracy requirements reported herein are defined with respect to worst-case scenarios, such as acquisitions conducted at high latitudes (e.g., near the polar regions) or during periods of elevated solar activity, where ionospheric perturbations are more pronounced. Nevertheless, in more favourable conditions —such as acquisitions in equatorial or mid-latitude regions during geomagnetically quiet periods— substantially improved performance is anticipated, potentially exceeding the stated requirements by up to an order of magnitude.

\begin{comment}
For main variables and per-pixel metadata (if applicable): Schematic table with accuracy depending on the possible input datasets
\begin{itemize}
    \item IRF degradation wrt SLC (tool SCT) to be evaluated on few S1 \& SAOCOM data
    \item PSLR degradation wrt SLC (tool SCT) to be evaluated on few S1 \& SAOCOM data 
    \item ... degradation wrt SLC (tool SCT) to be evaluated on few S1 \& SAOCOM data
    \item geolocation accuracy (tool SCT) in correspondence of CRS
    \item Spectral information preservation (doppler difference wrt input data)
    \item Coregistration accuracy (rg, az) (theoretical formulas, CRS?) as a function of Ancillary data / processing
    \item phase continuity between bursts (theoretical analysis) ???? as a function of Ancillary data / processing
\end{itemize}
\end{comment}

\subsubsection{Timeliness}
\label{sec:cslc_timeliness}

In principle \acs{cslc} product could be generated as soon as L1 data are available. The quality and the accuracy of the co-registration process strongly depends on the availability of precise orbital data, models of the atmosphere (troposphere and ionosphere). One could foresee multiple releases of the product, with successive corrections according to the availability of precise orbits and external models. A first release would be within a few hours after the acquisition, for time-critical applications. A first revision including precise orbital data could be provided after few days, and the final release could be delivered whenever numerical weather reanalysis data, ETAD and accurate ionospheric models are available.

 In this approach the only criticality is on the availability of precise orbital data (in general available in about one week). \autoref{fig:cslc_basic} shows the flowchart for the generation of \acs{cslc} product.

In addition, the provided timeline indicates the time required to collect the input data and complete the main processing steps. The processing time refers to a standard \acs{s1} dataset and is approximated, considering that the latency in obtaining the input data is the dominant factor.

\begin{figure}[H]
	\centering
	\begin{tikzpicture}[node distance=1cm and 2cm]
		
		% Define the style for the nodes
		% \tikzstyle{node} = [circle, draw, fill=blue!20, text centered, minimum size=1.5cm]
		
		% Draw nodes
		\node[rectin] (init) {\acs{slc}};
		\node[block, right=of init] (coreg_geo) {Geometric co-registration};
		\node[rectaux, below=of init] (dem) {\acs{dem}};
            \node[empty, below=of coreg_geo] (empty1) {};
		\node[rectaux, below=of empty1] (orbit) {Orbital data};
		\node[rectout, right=of coreg_geo] (cslc) {\acs{cslc}};
		
		% Draw edges with arrows
		\draw[->, thick] (init) -- (coreg_geo);
		\draw[->, thick] (dem) -| (coreg_geo);
		\draw[->, thick] (orbit) -- (coreg_geo);
		\draw[->, thick] (coreg_geo) -- (cslc);
        % Timeline with linked markers and adaptive length
        \def\timelineStart{init}
        \def\timelineEnd{cslc}
        \draw[thick, ->] (\timelineStart.west |- 0,-6) -- (\timelineEnd.east |- 0,-6) node[right] {Time};
        \foreach \node/\label in {init/Start, orbit/1h-3d, cslc/<1h} {
            \draw (\node.south |- 0,-6) -- (\node.south |- 0,-6.5) node[below] {\label};
            \draw[dotted, thick] (\node.south) -- (\node.south |- 0, -6.25); % Connects the node to the timeline
        }
    
	\end{tikzpicture}
	\caption{Sample workflow for generation of \acs{cslc} using geometric co-registration method.}
    \label{fig:cslc_basic}
\end{figure}

To further improve product quality, additional data, as listed in \autoref{sec:aux_data_cslc}, can be used. This approach introduces a delay in data production, which depends on the availability of external models: \acs{etad}-like, troposphere and ionosphere. \autoref{fig:cslc_ref_mdl} shows the flowchart for refinement of co-registration considering external models.

\begin{figure}[H]
	\centering
	\resizebox{\textwidth}{!}{
	\begin{tikzpicture}[node distance=1cm and 2cm]
		
		% Define the style for the nodes
		% \tikzstyle{node} = [circle, draw, fill=blue!20, text centered, minimum size=1.5cm]
		
		% Draw nodes
		\node[rectin] (init) {\acs{slc}};
		\node[block, right=of init] (coreg_geo) {Geometric Co-registration};
		\node[block, right=of coreg_geo] (coreg_ref) {Refinement};
		\node[rectaux, below=of init] (dem) {\acs{dem}};
		\node[empty, below=of coreg_geo] (empty1) {};
            \node[rectaux, below=of empty1] (orbit) {Orbital data};
		\node[empty, below=of coreg_ref] (empty2) {};
		\node[rectaux, below=of empty2] (models) {External models};
		\node[rectout, right=of coreg_ref] (cslc) {\acs{cslc}};
		
		% Draw edges with arrows
		\draw[->, thick] (init) -- (coreg_geo);
		\draw[->, thick] (coreg_geo) -- (coreg_ref);
		\draw[->, thick] (dem) -| (coreg_geo);
		\draw[->, thick] (orbit) -- (coreg_geo);
		\draw[->, thick] (models) -- (coreg_ref);
		\draw[->, thick] (coreg_ref) -- (cslc);
            % Timeline with linked markers and adaptive length
            \def\timelineStart{init}
            \def\timelineEnd{cslc}
            \draw[thick, ->] (\timelineStart.west |- 0,-6) -- (\timelineEnd.east |- 0,-6) node[right] {Time};
            \foreach \node/\label in {init/Start, orbit/1h-3d, models/\textbf{1-30d}, cslc/<1h} {
                \draw (\node.south |- 0,-6) -- (\node.south |- 0,-6.5) node[below] {\label};
                \draw[dotted, thick] (\node.south) -- (\node.south |- 0, -6.25); % Connects the node to the timeline
            }
    
	\end{tikzpicture}
	}
	\caption{Sample workflow of generation of \acs{cslc} using geometric co-registration method plus co-registration refinement with external models.}
    \label{fig:cslc_ref_mdl}
\end{figure}


Alternatively, a data-driven approach can be followed, estimating corrections directly from a stack of SAR data. In this case, the processing is more complex, generally involving more than a single image and requiring the solution of inverse problems with high computational complexity.

\autoref{fig:cslc_data_driven} shows the flowchart for the generation of \acs{cslc} product, considering data-driven co-registration method.

\begin{figure}[H]
	\centering
	\resizebox{\textwidth}{!}{
	\begin{tikzpicture}[node distance=1cm and 2cm]
		
		% Define the style for the nodes
		% \tikzstyle{node} = [circle, draw, fill=blue!20, text centered, minimum size=1.5cm]
		
		% Draw nodes
		\node[rectin] (init) {\acs{slc}};
		\node[block, right=of init] (coreg_geo) {Geometric co-registration};
		\node[block, right=of coreg_geo] (coreg_data_driven) {Data-driven Co-registration};
		\node[rectaux, below=of init] (dem) {\acs{dem}};
		\node[empty, below=of coreg_geo] (empty1) {};
            \node[rectaux, below=of empty1] (orbit) {Orbital data};
		\node[empty, below=of coreg_data_driven] (empty2) {};
		\node[rectin, below=of empty2] (ref_stack) {Reference SLC stack};
		\node[rectout, right=of coreg_data_driven] (cslc) {\acs{cslc}};
		
		% Draw edges with arrows
		\draw[->, thick] (init) -- (coreg_geo);
		\draw[->, thick] (coreg_geo) -- (coreg_data_driven);
		\draw[->, thick] (dem) -| (coreg_geo);
		\draw[->, thick] (orbit) -- (coreg_geo);
		\draw[->, thick] (ref_stack) -- (coreg_data_driven);
		\draw[->, thick] (coreg_data_driven) -- (cslc);
            % Timeline with linked markers and adaptive length
            \def\timelineStart{init}
            \def\timelineEnd{cslc}
            \draw[thick, ->] (\timelineStart.west |- 0,-6) -- (\timelineEnd.east |- 0,-6) node[right] {Time};
            \foreach \node/\label in {init/Start, orbit/1h-3d, ref_stack/<1h, cslc/<1h} {
                \draw (\node.south |- 0,-6) -- (\node.south |- 0,-6.5) node[below] {\label};
                \draw[dotted, thick] (\node.south) -- (\node.south |- 0, -6.25); % Connects the node to the timeline
            }
	\end{tikzpicture}
	}
	\caption{Sample workflow for generation of \acs{cslc} using geometric co-registration method plus co-registration refinement exploiting data-driven approach considering a \acs{cslc} stack.}
    \label{fig:cslc_data_driven}
\end{figure}


\subsection{Usability, Integrability and Interoperability}
\label{sec:usability_integrability_CSLC}
The characteristics of the \acs{cslc} and \acs{gslc} products in terms of usability, integrability and interoperability are in the following tables.

\begin{longtblr}[
  caption = {Usability, Integrability, and Interoperability of \acs{cslc} product},
  label = {tab:cslc_usability_integrability_interoperability},
]{
  colspec = {Q[2cm] Q[12cm]},
  rowhead = 1,
  width = \linewidth
}
\hline
\textbf{Aspect} & \textbf{Description} \\
\hline
Usability &
\begin{itemize}
  \item Naturally suited for time-series analysis (e.g., surface deformation, change detection, etc.)
  \item Preserves SAR geometry for reliable spectral and filtering operations
  \item Enables phase-based and amplitude-based applications like InSAR and pixel offset tracking
\end{itemize}
\\
Integrability &
\begin{itemize}
  \item Acts as a core intermediate product in SAR processing chains (e.g., interferometry, change detection)
  \item Easily combined in multi-temporal and joint analyses
\end{itemize}
\\
Interoperability &
\begin{itemize}
  \item Cross-sensor compatibility for ROSE-L and Sentinel-1 missions
  \item Works with common \acs{eo} software
\end{itemize}
\\
\hline
\end{longtblr}

\begin{longtblr}[
  caption = {Usability, Integrability, and Interoperability of \acs{gslc} product},
  label = {tab:gsclc_characteristics},
]{
  colspec = {Q[2cm] Q[12cm]},
  rowhead = 1,
  width = \linewidth
}
\hline
\textbf{Aspect} & \textbf{Description} \\
\hline
Usability &
\begin{itemize}
  \item Geocoded and projected to ground range \acs{cslc} or \acs{slc}, making it more intuitive and accessible for non-expert users
  \item Enables direct visual interpretation and integration with other map-based or \acs{gis} data
  \item When geocoded with adequate accuracy, phase data is suitable for generating interferograms directly on the ground
  \item If radiometrically corrected, the \acs{gslc} amplitude can carry the same information as the \acs{nrb} product
\end{itemize}
\\
Integrability &
\begin{itemize}
  \item Directly usable in \acs{gis} platforms, map viewers, or decision-support tools
  \item Seamlessly integrates into workflows for change detection, urban monitoring, land deformation, and disaster response
  \item If radiometric correction is applied, it can substitute or complement the \acs{nrb} in joint or thematic analyses
\end{itemize}
\\
Interoperability &
\begin{itemize}
  \item High interoperability due to standard geographic referencing and pixel alignment across scenes
  \item Compatible with multi-source datasets (e.g., optical imagery, elevation models, land cover maps)
  \item Adheres to ARD principles, enhancing integration with downstream services and platforms
  \item Facilitates cross-mission usage through harmonized ground projection
  \item Compatible with common \acs{eo} frameworks and pipelines using georeferenced data
\end{itemize}
\\
\hline
\end{longtblr}



\section{Normalized Radar Backscatter (NRB)}
\label{s:nrb}

\subsection{Product Description}
\label{sec:nrb_prod_description}

%\begin{comment}
This product exclusively contains backscatter values and undergoes Radiometric Terrain Correction (\acs{rtc}), providing \gls{gamma0} as defined in \cite{Small2011}, which mitigates the impact of different observation geometries. It is the most readily comprehensible product, enabling direct visualisation for visual interpretation, as well as deeper analysis. Indeed, analysing backscatter time series can yield information regarding the observed object's type for classification purposes (land use, land cover), as well as its temporal variations, enabling applications such as change detection, security applications, and vegetation growth monitoring.

\acs{nrb} products have gained significant advancements in recent years, driven by both technological innovation and increasing demand for high-quality, precise data.

The Normalized Radar Backscatter (\acs{nrb}) product is foreseen and provided in several projects and it is the main \acs{ard} product in the portfolio of different operative and programmed missions.

The \acs{nrb} can be obtained through two possible alternative processing flows:
   
   1. Starting from \acs{slc} \acs{s1} data with the same Data Take ID, subject to the co-registration procedure (\acs{cslc}); \\
   2. Starting from \acs{slc} \acs{s1} data with the same Data Take ID, without applying the co-registration procedure.

The dataset is formatted as a ZARR file. This format supports the inclusion of both the processed measurement data and the associated metadata within the same file, making the product self-contained and self-describing, with no need for additional annotations or supporting files.

\subsubsection{Input and Auxiliary Data}
The following input and auxiliary data are considered for product generation:

\begin{longtblr}[
    label = tab:in_aux_nrb_products,
    entry = none,
    caption = {\acs{nrb} input and auxiliary data},
    ]{
        width = \linewidth,
        colspec = {Q[l,20]Q[l,60]Q[c,20]},
        hlines,
        hline{1-Z} = {-}{0.05em},
    }
    \textbf{Name}   & \textbf{Description}   & \textbf{Mandatory (M) / Optional (O)} \\
    \acs{slc} or \acs{cslc} & L1 Single Look Complex image or co-registered L1 Single Look Complex image and annotations & M \\
    \acs{dem} & External DEM (e.g., Copernicus DEM at different resolutions, with 30 m as the default worldwide). & M \\
    Orbits & \acs{slc} Standard, restituted or precise orbits & M \\
    \acs{etad} & \acs{etad}-like products for atmospheric delay and timing corrections & O \\
    \acs{tec} & External \acs{tec} data for delay correction & O \\
\end{longtblr}

\subsubsection{Main Variables}
\label{sec:variables_nrb}
The following tables list the main variables associated with the \acs{nrb} product.

\begin{longtblr}[
    label = tab:var_nrb_products,
    entry = none,
    caption = {\acs{nrb} main variables},
    ]{
        width = \linewidth,
        colspec = {Q[l,25]Q[l,52]Q[l,20]Q[l,13]},
        hlines,
        hline{1-Z} = {-}{0.05em},
    }
    \textbf{Name} & \textbf{Description} & \textbf{Dimensions} & \textbf{Data Type} \\
    gamma0 & Radiometric Terrain Correction \gls{gamma0} backscatter linear scaling image & polarisation, acquisition\_date, x, y & float32 \\
    local\_inc\_angle & DEM-based local Incidence angle image [deg]. This variable is present only if the \acs{nrb} is generated from \acs{slc}s. If the \acs{nrb} is derived from \acs{cslc}s, this variable will be null, and the reference\_inc\_angle metadata will be present & x, y & float32 \\
    anf & Area normalization factor or local contributing area, used to obtain the gamma0 coefficient. This variable is present only if the \acs{nrb} is generated from \acs{slc}s. If the \acs{nrb} is derived from \acs{cslc}, this variable will be null, and the reference\_anf metadata will be present & x, y & float32 \\
\end{longtblr}

Note that the fields "anf" and "local\_inc\_angle" are always included: as variables for products directly derived from \acs{slc}, and as metadata for products derived from \acs{cslc}, with links to the corresponding static layer. Either way, these variables are 2D layers, defined for each projected cell.

Furthermore, the knowledge of both \gls{gamma0}, local incidence angle ($\theta_{inc}$) and the area normalization factor (\acs{anf}) layers allows the computation of both \gls{sigma0} and \gls{beta0} with simple mathematical operations, executed on a virtual layer. Specifically these relations hold:

\begin{equation}
\begin{cases}
\sigma^0 = \gamma^0 \cdot \cos(\theta_{inc}) \\
\beta^0 = \sigma^0 \cdot \text{ANF}
\end{cases}
\end{equation}

The table below describes the dimensions/coordinates associated with the \acs{nrb} main variables.

\begin{longtblr}[
    label = tab:var_nrb_coordinates,
    entry = none,
    caption = {\acs{nrb} main variables coordinates/dimensions},
    ]{
        width = \linewidth,
        colspec = {Q[l,20]Q[l,62]Q[l,18]},
        hlines,
        hline{1-Z} = {-}{0.05em},
    }
    \textbf{Name} & \textbf{Description} & \textbf{Data Type} \\
    polarisation & Acquired polarisations & string \\
    acquisition\_date & Image acquisition dates, formatted as YYYY-MM-DD & string \\
    x & UTM/UPS x (East) coordinates [m] & float64 \\
    y & UTM/UPS y (North) coordinates [m] & float64 \\
\end{longtblr}

\subsubsection{Main Metadata}
\label{sec:nrb_product_metadata}
The following table lists the main metadata and per-pixel metadata associated with the \acs{nrb} product.

\begin{longtblr}[
    label = tab:meta_nrb_products,
    entry = none,
    caption = {\acs{nrb} main metadata and per-pixel metadata},
    ]{
        width = \linewidth,
        colspec = {Q[l,25]Q[l,38]Q[l,10]Q[l,10]Q[c,12]},
        hlines,
        hline{1-Z} = {-}{0.05em},
    }
    \textbf{Name} & \textbf{Description} & \textbf{Dimensions} & \textbf{Data Type} & \textbf{Per-pixel Metadata}\\
    reference\_slc / reference\_cslc & List of links to input reference \acs{slc} or \acs{cslc} & - & list of string & \\
    reference\_dem & Link to reference DEM & - & string & \\
    etad\_file & ETAD file used for each image belonging to the \acs{cslc}/\acs{slc} stack. This is optional: it is included only if used; otherwise, it is null & - & string & \\ 
    tec & Link to external \acs{tec} data, if used & - & string & \\ 
    reference\_inc\_angle & List of links to reference DEM-based local incidence angle maps. It is present only if the \acs{nrb} is generated from \acs{cslc}s. & - & string & \ding{51} \\
    reference\_anf & List of links to reference area normalization factor or local contributing area. It is present only if the \acs{nrb} is generated from \acs{cslc}s. & - & string & \\
    orig\_data\_id & Per pixel metadata reporting an index describing the bursts contributing to the output mosaicked pixel & x, y & int32 & \ding{51} \\
    mask & Mask containing, for each ground cell, the following information:
    
     - 0: no distortions
    
     - 1: layover
    
     - 2: shadow
    
     - 3: undefined & x, y & byte & \ding{51} \\
    nlooks & Number of equivalent looks for the gamma flattening & x, y & float32 & \ding{51} \\
    \_geometric\_accuracy & Measured geometric accuracy, including eastern and northern geometric accuracy bias and standard deviation. The measures will be performed on a validation dataset and periodically updated. & - & dictionary & \\ 
    east\_geom\_bias & Measured eastern geometric bias & - & float32 & \\
    east\_geom\_std & Measured eastern geometric standard deviation & - & float32 & \\
    north\_geom\_bias & Measured northern geometric bias & - & float32 & \\
    north\_geom\_std & Measured northern geometric standard deviation & - & float32 & \\
    geometric\_accuracy
    \_radial\_rmse& Measured radial RMSE of geometric accuracy & - & float32 & \\
    x\_sample & Pixel spacing in x direction & - & float32 & \\
    y\_sample & Pixel spacing in y direction & - & float32 & \\
\end{longtblr}

The table below describes the dimensions/coordinates associated with the \acs{nrb} main metadata.

\begin{longtblr}[
    label = tab:var_nrb_metadata_coordinates,
    entry = none,
    caption = {\acs{nrb} main metadata coordinates/dimensions},
    ]{
        width = \linewidth,
        colspec = {Q[l,20]Q[l,62]Q[l,18]},
        hlines,
        hline{1-Z} = {-}{0.05em},
    }
    \textbf{Name} & \textbf{Description} & \textbf{Data Type} \\
    x & UTM/UPS x (East) coordinates [m] & float64 \\
    y & UTM/UPS y (North) coordinates [m] & float64 \\
\end{longtblr}


\subsubsection{Reference System Tiling and Slicing}
\label{sec:nrb_reference_system_tiling}
\autoref{tab:nrb_ref_system_slicing} reports the reference system, the slicing, and grid for \acs{nrb} products.

\begin{longtblr}[
    label=tab:nrb_ref_system_slicing,
    entry = none,
    caption = {Reference system and Slicing for \acs{nrb}},
    ]{
        width=\linewidth,
        colspec = {Q[l,15]Q[l,35]Q[l,35]Q[l,15]},
        %vline{2-5} = {-}{},
        %vline{5-5} = {-}{},
        hlines,
        hline{1-Z} = {-}{0.05em},
        rowhead=2
    }
    \textbf{Product} & \textbf{Reference System} & \textbf{Slicing} & \textbf{Spacing}\\
    \textbf{\acs{nrb}} (default) & Cartographic: UTM/UPS & US-\acs{mgrs} 100 km x 100 km & 10 m\\
    \textbf{\acs{nrb}} (optional) & Cartographic: UTM/UPS & US-\acs{mgrs} 100 km x 100 km & 30 m\\
\end{longtblr}

\subsection{Product Requirements}
\label{sec:impacts_nrb}

The main quality parameters for evaluating an \acs{nrb} product are radiometric and geometric accuracies. The former is determined by the ability to compensate for local slope due to topography and obtain an estimate of the radar cross-section that is independent of the observation geometry. The objective is to reconstruct the intrinsic characteristics of the observed scene. Additionally, proper filtering could be applied to reduce speckle noise and other sources of error while maintaining the desired output resolution. The latter refers to the capability to accurately relocate points from the radar to a geographic reference system, compensating for distortion effects due to the acquisition mode and artifacts related to the propagation of the signal through the atmosphere.

\subsubsection{Accuracy}

Geolocation accuracy is inherited from the CSLC product. All geolocation accuracies values are expressed as fractions of the nominal resolution in the corresponding dimension, providing a normalized measures. All the accuracy thresholds are evaluated at 1-$\sigma$.

\begin{longtblr}[
	label = tab:nrb_accuracy_cband,
	entry = none,
	caption = {Accuracy Requirements for \acs{nrb} product (C-Band)},
	]{
		width = \linewidth,
		colspec = {Q[l,30] Q[l,60] Q[c,20]},
		hlines,
		hline{1-Z} = {-}{0.05em},
		rowhead = 1,
	}
	
	\textbf{Metric} & \textbf{Description} & \textbf{Requirement} \\
	
	\SetCell[r=2]{l} Range Accuracy (\%) 
	& Model-Driven (only orbit and \acs{dem}) & $\leq 5\%$\\
	& Model-Driven + atmospheric correction & $\leq 2\%$ \\
	
	\SetCell[r=2]{l} Azimuth Accuracy (\%) 
	& Model-Driven (only orbit and \acs{dem}) & $\leq 2\%$\\
	& Data-Driven (\acs{esd}) & $\leq 0.5\%$ \\
	
	\acs{rcs} difference (dB) & Stability of the \acs{rcs} comparing different acquisition geometries over the same area considering a resolution of 30m & $\leq 5$~dB \\
	
\end{longtblr}

\begin{longtblr}[
	label = tab:nrb_accuracy_lband,
	entry = none,
	caption = {Accuracy Requirements for \acs{nrb} product (L-Band)},
	]{
		width = \linewidth,
		colspec = {Q[l,30] Q[l,60] Q[c,20]},
		hlines,
		hline{1-Z} = {-}{0.05em},
		rowhead = 1,
	}
	
	\textbf{Metric} & \textbf{Description} & \textbf{Requirement} \\
	
	\SetCell[r=3]{l} Range Accuracy (\%) 
	& Model-Driven (only orbit and \acs{dem}) & $\leq 60\%$\\
	& Model-Driven + atmospheric correction & $\leq 25\%$ \\
	& Model-Driven + atmospheric correction + Data-Driven (\acs{icc}) & $\leq 10\%$ \\
	
	\SetCell[r=3]{l} Azimuth Accuracy (\%) 
	& Model-Driven (only orbit and \acs{dem}) & $\leq 150\%$ \\
    & Model-Driven + atmospheric correction & $\leq 50\%$ \\
	& Data-Driven (\acs{icc}) & $\leq 10\%$ \\
	
	\acs{rcs} difference (dB) & Stability of the \acs{rcs} comparing different acquisition geometries over the same area considering a resolution of 30m & $\leq 5$~dB \\
	
\end{longtblr}

\subsubsection{Timeliness}
\label{sec:nrb_timeliness}

The \acs{nrb} product can be generated as soon as L1 data is available for each \acs{slc} image independently or starting from \acs{cslc} products. Besides choosing between \acs{slc} or \acs{cslc} products, the main factor impacting the timeline is the availability of external data or reference images, which are required to refine geolocation and mitigate atmospheric effects, particularly ionospheric effects. \autoref{fig:nrb_single} shows the flowchart of generating the \acs{nrb} product using a single \acs{slc} image.

In addition, the timeline provided indicates the time required to collect the input data and complete the main processing steps. The processing time refers to a standard \acs{s1} dataset and is approximated considering that the latency in obtaining the input data is the dominant factor.

\begin{figure}[H]
		\centering
		\resizebox{\textwidth}{!}{
		\begin{tikzpicture}[node distance=0.5cm and 2cm]
%			
%			% Define the style for the nodes
%			% \tikzstyle{node} = [circle, draw, fill=blue!20, text centered, minimum size=1.5cm]
%			
%			% Draw nodes
			\node[rectin] (init) {\acs{slc}};
			\node[block, right=of init] (g2r) {Geo-to-Radar};
			\node[block, right=of g2r] (rtc) {\acs{rtc}};
			\node[block, right=of rtc] (geocoding) {Geocoding};
			\node[rectaux, below=of init] (dem) {\acs{dem}};
			\node[empty, below=of g2r] (empty0) {};
			%\node[empty, below=of empty0] (empty1) {};
			\node[rectaux, below=of empty0] (orbit) {Orbital data};
			\node[rectauxd, below=of orbit] (models) {External models};
%			\node[rectind, below=of models] (ref_stack) {Reference SLC stack};
			\node[rectout, right=of geocoding] (nrb) {\acs{nrb}};
			
%			% Draw edges with arrows
			\draw[->, thick] (init) -- (g2r);
			\draw[->, thick] (g2r) -- (rtc);
			\draw[->, thick] (rtc) -- (geocoding);
			\draw[->, thick] (dem) -| (g2r);
			\draw[->, thick] (orbit) -- (g2r);
			\draw[dotted, thick] (models) -- (orbit);
%			\draw[->, thick] (ref_stack) -| (g2r);
			\draw[->, thick] (geocoding) -- (nrb);
%			% Timeline with linked markers and adaptive length
			\def\timelineStart{init}
			\def\timelineEnd{nrb}
			\draw[thick, ->] (\timelineStart.west |- 0,-6) -- (\timelineEnd.east |- 0,-6) node[right] {Time};
			\foreach \node/\label in {dem/Start, models/\textbf{1-30d}, nrb/<1h} {
				\draw (\node.south |- 0,-6) -- (\node.south |- 0,-6.5) node[below] {\label};
				\draw[dotted, thick] (\node.south) -- (\node.south |- 0, -6.25); % Connects the node to the timeline
			}
				
		\end{tikzpicture}
		}
		\caption{Sample workflow for generation of \acs{nrb} considering each \acs{slc} image separately.}
		\label{fig:nrb_single}
\end{figure}
	
In the framework of this project, the \acs{cslc} procedure includes geometric correctness, even in case of parameter compensation that does not introduce differential errors between reference and secondary images. Taking into account this characteristic, if the \acs{nrb} is generated from the \acs{cslc} then all the calibrations have already been performed during the \acs{cslc} processing, and the reference geometry is computed only once for the reference orbit. Therefore, the schema simplifies as reported in \autoref{fig:nrb_from_cslc}, but the timeline does not change depending on the information needed to perform the Geo-to-Radar conversion or to generate the \acs{cslc} product.
\begin{figure}[H]
		\centering
		\resizebox{\textwidth}{!}{
		\begin{tikzpicture}[node distance=0.5cm and 2cm]
			
			% Define the style for the nodes
			% \tikzstyle{node} = [circle, draw, fill=blue!20, text centered, minimum size=1.5cm]
			
			% Draw nodes
			\node[rectin] (init) {\acs{slc}};
			\node[block, right=of init] (coreg) {Co-registration};
			\node[rectout, right=of coreg] (cslc) {\acs{cslc}};
			\node[block, right=of cslc] (rtc) {\acs{rtc}};
			\node[block, right=of rtc] (geocoding) {Geocoding};
			\node[rectout, right=of geocoding] (nrb) {\acs{nrb}};
			
			\node[rectaux, below=of init] (dem) {\acs{dem}};
			\node[empty, below=of coreg] (empty0) {};
			%\node[empty, below=of empty0] (empty1) {};
			\node[empty, below=of empty0] (empty2) {};
			\node[rectaux, below=of empty2] (orbit) {Orbital data};
			\node[rectauxd, below=of orbit] (models) {External models};
			\node[rectind, below=of dem] (ref_stack) {Reference SLC stack};
			\node[rectout, right=of geocoding] (nrb) {\acs{nrb}};
			
			% Draw edges with arrows
			\draw[->, thick] (init) -- (coreg);
			\draw[->, thick] (coreg) -- (cslc);
			\draw[->, thick] (cslc) -- (rtc);
			\draw[->, thick] (rtc) -- (geocoding);
			\draw[->, thick] (geocoding) -- (nrb);
			
			% Draw edges with arrows
			\draw[->, thick] (dem) -| (coreg);
			\draw[->, thick] (orbit) -- (coreg);
			\draw[dotted, thick] (models) -- (orbit);
			\draw[->, thick] (ref_stack) -| (coreg);
			\draw[->, thick] (geocoding) -- (nrb);
			%			% Timeline with linked markers and adaptive length
			\def\timelineStart{init}
			\def\timelineEnd{nrb}
			\draw[thick, ->] (\timelineStart.west |- 0,-8) -- (\timelineEnd.east |- 0,-8) node[right] {Time};
			\foreach \node/\label in {ref_stack/Start, models/\textbf{1-30d}, nrb/<1h} {
				\draw (\node.south |- 0,-8) -- (\node.south |- 0,-8.5) node[below] {\label};
				\draw[dotted, thick] (\node.south) -- (\node.south |- 0, -8.25); % Connects the node to the timeline
			}
			
		\end{tikzpicture}
		}
		\caption{Sample workflow for generation of \acs{nrb} starting from \acs{cslc} products.}
		\label{fig:nrb_from_cslc}
\end{figure}
%(required? look at applications? where? ROSE-L MRD? NISAR?)


\subsection {Usability, Integrability and Interoperability}
\label{sec:nrb_usability_integrability}
The characteristics of the \acs{nrb} products in terms of usability, integrability and interoperability are in the following table.

\begin{longtblr}[
  caption = {Usability, Integrability, and Interoperability of \acs{nrb} product},
  label = {tab:nrb_usability_integrability_interoperability},
]{
  colspec = {Q[2cm] Q[12cm]},
  rowhead = 1,
  width = \linewidth
}
\hline
\textbf{Aspect} & \textbf{Description} \\
\hline
Usability &
\begin{itemize}
  \item Radiometric calibration eliminates the need for users to correct backscatter manually.
  \item \gls{gamma0} convention enables direct comparison across sensors, acquisition times, and geometries.
  \item Orthorectified and georeferenced product makes data easier to use in mapping applications without terrain distortions.
  \item Provided in standard formats like ZARR ensures wide compatibility.
  \item The product can directly be used in classification and/or segmentation procedures.
  \item The product can be provided by using multi-polarisation: availability of VV, VH (or other) polarisations improves surface characterization.
  \item Note: Usability depends on the spatial and temporal resolution, as well as whether speckle filtering has been applied. A 10 m resolution is suitable for most applications that require radiometrically terrain-corrected backscatter maps.
\end{itemize}
\\
Integrability &
\begin{itemize}
  \item GIS Compatibility: Easily used within \acs{gis} platforms, due to geographic referencing and data format.
  \item Multi-Sensor Fusion: Can easily be integrated with optical or other SAR projected products for land cover, surface change, or hazard mapping.
  \item Environmental Monitoring: It supports use in models or indicators for forest cover, crop health or hydrology.
\end{itemize}
\\
Interoperability &
\begin{itemize}
  \item Cross-Mission Compatibility: Normalized \gls{gamma0} allows for comparison between satellites.
  \item Interdisciplinary Applications: Usable in hydrology, agriculture, forestry, urban studies, etc., enhancing cross-domain use.
  \item Compatible with common \acs{eo} frameworks and pipelines using georeferenced data
  \item High interoperability due to standard geographic referencing and pixel alignment across scenes
  \item Adheres to ARD principles, enhancing integration with downstream services and platforms
  \item Facilitates cross-mission usage through harmonized ground projection
\end{itemize}
\\
\hline
\end{longtblr}




\section{Polarimetric products (COV \& GCOV)}
\label{sec:pol_products}

\subsection{Product Description}
\label{sec:description_pol}

\acs{sar} polarimetry entails the simultaneous acquisition of a set of \acs{slc} images using different combinations of transmitted and received polarisations, and each combination is denoted as \textit{polarimetric channel}. Depending on the sensor and acquisition mode, there may be only one polarimetric channel (single-pol mode), two channels (dual-pol mode), or four channels (quad-pol mode). Using as inputs the \acs{slc} images at the polarimetric channels available, the polarimetric product is formed as a covariance matrix, as described in D1~\cite{D1}. Therefore, this \acs{ard} product contains the entries of the polarimetric covariance matrix, either in radar coordinates (\acs{cov} product) or in cartographical coordinates (\acs{gcov} product). In order to preserve the polarimetric information, it is necessary to apply the same processing to generate each one of the matrix entries. All the possible forms of the covariance matrix are detailed in D1~\cite{D1}, and not repeated here to save space.

A convenient way to explain the generation and the resulting format of the \acs{cov} and \acs{gcov} products is by re-using what is explained in \autoref{s:nrb} for the \acs{nrb} product. All matrix entries are processed as if each one of them constitutes a \acs{nrb} product. For instance, they
usually undergo Radiometric Terrain Correction (\acs{rtc}), providing Gamma-Nought as defined in \cite{Small2011}, which mitigates the impact of differing observation geometries. 

In fact, the diagonal entries of the covariance matrix correspond to the backscattering coefficient of each individual polarimetric channel, so the diagonal entries in a \acs{gcov} product coincide with the \acs{nrb} products of all channels. For instance, taking the VV-VH dual-pol mode as an example, typical of Sentinel-1 acquisitions over land, the diagonal elements  correspond to the \acs{nrb} of the VV and VH channels. As for the off-diagonal entries, they are formed by applying the same processing of the \acs{nrb} generation but to the product VV$\cdot$VH$^*$, where $*$ denotes the complex conjugate, instead of to the intensity of the VV and VH channels. In summary, \acs{cov} and \acs{gcov} products for the dual-pol case ($C_2$ matrix) contain only three layers: the two real-valued intensities (backscattering coefficients) of the two channels, and one complex value corresponding to the correlation between the two channels.

It must be pointed out that the \acs{nrb} product may be regarded either as a \textit{subset} or as a \textit{particular case} of the \acs{gcov} product, as it is fully contained in it. Such a point of view is followed by NASA in the \acs{gcov} product defined for NISAR~\cite{NISAR_GCOV_ProdSpec}, in which the \acs{nrb} product is just a particular case included in the \acs{gcov} specification.

The \acs{gcov} product shares all the properties of \acs{nrb} products, hence  enabling direct visualisation in geographical information systems for interpretation and analysis, but enriches the information content of \acs{nrb} products with the extra channels. Polarimetry outperforms backscatter in most final applications, as demonstrated in literature~\cite{Hajnsek2021}. 

Regarding the \acs{cov} product, i.e. expressed in radar coordinates, it may be preferred by some users which need to apply some specific processing steps in radar geometry to avoid distortions in the information content, like speckle filtering with a multi-look or boxcar, before geocoding the resulting matrix. In such a case, the \acs{cov} product can be provided together with a look-up table correspondence for geocoding, to be applied later by the user to convert the product to cartographical coordinates. 

As in the case of \acs{nrb}, the \acs{cov} and \acs{gcov} products can be obtained through two possible alternative processing flows:
\begin{enumerate}
    \item Starting from \acs{slc} data with the same Data Take ID, subject to the co-registration procedure (\acs{cslc});
    \item Starting from \acs{slc} data with the same Data Take ID, without applying the co-registration procedure.
\end{enumerate}

Due to the specific nature of the polarimetric \acs{ard} products, some considerations and clarifications are detailed next. 
%In addition, as for other products, existing SW tools and computational complexity aspects are commented in \autoref{ss:algoritms_pol_SW} and \autoref{ss:algoritms_pol_comput}, respectively.

\paragraph*{Reference channels}

Some processing associated with the co-registration and geocoding steps may be carried out using the actual data from the images. For instance, co-registration based on enhanced spectral diversity (ESD) exploits areas imaged by overlapping bursts. Similarly, point targets can be used to refine azimuth times of a stack of images with respect to a reference image. Regarding geocoding, a synthetic or simulated image generated from a \acs{dem} can be cross-correlated with the \acs{sar} image to improve the geocoding accuracy.
In all these cases, the computations (e.g. cross-correlation)  must be performed using only ONE of the polarimetric channels, whereas the result (e.g. interpolation of the secondary images at the computed positions) is applied to ALL the polarimetric channels.

For these cases, the selected channel is usually a co-polar one, i.e. HH or VV, because it provides a higher backscatter level and hence a higher \acs{snr} than the cross-polar channels. Consequently, the reference channel must be indicated in the metadata.

\paragraph*{Polarimetric symmetrization}

Every image acquired in quad-polarisation mode consists of four \acs{slc} images measured at the following channels: HH, HV, VH, and VV. Due to reciprocity, and after an adequate calibration, the two cross-polar channels (HV and VH) should be equal and are only different due to the presence of noise. For this reason, for the definition of the covariance matrix, only three channels are considered, i.e. HH, VV and one cross-polar channel (HV or VH). 

A common practice, also considered in the definition of the \acs{gcov}-NISAR product~\cite{NISAR_GCOV_ProdSpec}, consists in averaging the two cross-polar channels and using the resulting average in the construction of the covariance matrix in the \acs{ard} product. This operation is called symmetrization in polarimetric literature.

This step is not commented by \acs{ceos}, but the \acs{gcov}-NISAR definition in~\cite{NISAR_GCOV_ProdSpec} indicates that the metadata should include:
\begin{itemize}
\item A flag to indicate whether the symmetrization has been applied or not.
\item A reference to the symmetrization algorithm employed, if applicable.
\end{itemize}

\paragraph*{Order of spatial filtering and geocoding}
\label{p:order_pol}

In order to minimise the bias in the \acs{gcov} products, speckle filtering and/or spatial averaging of the covariance matrix should be applied in radar coordinates (slant-range and cross-range). Afterwards, with a resolution degraded by this filtering, geocoding should be performed using either nearest-neighbour resampling or the approach proposed in~\cite{Shiroma2022}. Specifically, nearest-neighbour resampling ensures that the averaged covariance
matrix elements in slant range and in geocoded ground projection are exactly the same (obviously assuming similar posting in both coordinate projections). If the approach in~\cite{Shiroma2022} is used, the degraded resolution makes it equivalent to nearest-neighbour resampling except in the borders between pixels, but no significant artifact is expected.

In the definition of the \acs{gcov}-NISAR product~\cite{NISAR_GCOV_ProdSpec}, the estimation of the output covariance matrix and its projection onto geographical coordinates is carried out simultaneously by means of the area-based interpolation associated with the \acs{rtc} proposed by Shiroma et al. in~\cite{Shiroma2022}. Therefore, in that case there is no degradation of the polarimetric information, since the averaging is carried out using the entries of the covariance matrix of the original pixels in radar coordinates.

\paragraph*{Polarimetric speckle filtering}

Speckle filter is required for extracting useful quantitative information from the \acs{sar} data, especially over natural areas. It is usually carried out using a spatial filtering of the image, which entails a degradation of the spatial resolution. To mitigate this issue, spatially-adaptive speckle filters are employed, which try to simultaneously preserve details and average homogeneous distributed areas. For this purpose, at each pixel to be filtered, a group of pixels with similar scattering properties is selected in an adaptive fashion. In this context, spatially-adaptive speckle filtering can clearly benefit from the polarimetric information of the images, expressed in covariance matrix form, since it allows for identifying similar and different land covers and elements in the scene. 

Many speckle filtering algorithms have been proposed which incorporate the polarimetric information in the selection of the group of homogeneous pixels, like the polarimetric intensity-driven adaptive-neighbourhood (IDAN) filter~\cite{Vasile2006}, the non-local filter~\cite{Deledalle2015}, tensor decomposition-based filter~\cite{Luo2023}, etc.
In all cases, the obtained performance is better than the one obtained with filters which do not exploit this information or with filters with fixed kernels (e.g. boxcar).

Speckle filters should be applied in radar coordinates (i.e. before geocoding). In this sense, the simultaneous area-based averaging and geocoding considered in the \acs{gcov}-NISAR product~\cite{NISAR_GCOV_ProdSpec} disables the possibility to rigorously apply these adaptive filters to the (already geocoded) product provided as a \acs{gcov}. It is for this reason that the \acs{cov} product is required by some users.


\paragraph*{Considerations about the Equivalent Number of Looks}

The use of spatially-adaptive speckle filters and/or the area-based interpolation associated with the \acs{rtc} defined by Shiroma et al. in~\cite{Shiroma2022} yields a variable equivalent number of looks (ENL) over the image. As a result, the layer with the values of ENL has to be available in the per-pixel metadata of the \acs{ard} products.

The ENL value is later used to help interpret or exploit some derived products. For instance, the maximum likelihood change detector for polarimetric images (also valid for single-pol backscatter data) is based on the generalised likelihood ratio test computed by using covariance matrices~\cite{Conradsen2003, Nielsen2015}. In this test, many parameters (statistical significance, false alarm probability, thresholds, etc.) depend on the ENL of the two pixels being compared, so this information must be provided together with the \acs{ard} product (as a metadata layer).

Another aspect to highlight regarding ENL is the minimum required value to properly employ some techniques. The multi-dimensional nature of the polarimetric information, exemplified by the matrix form of the data (i.e. the covariance matrix) makes it necessary to have a large number of looks to measure some polarimetric features. For instance, the entropy, alpha and anisotropy products~\cite{Cloude1997, Lee_Pottier2009} derived from the eigenvector decomposition of the covariance matrix (after transformation to coherency matrix) are widely employed in scientific studies. The estimation of these products is known to be biased when the number of looks is low, requiring large values (50+) for sufficiently reducing that bias~\cite{LopezMartinez2005, Lee2008}.

As a consequence, the ENL resulting from the area-based interpolation associated with the \acs{rtc} defined by Shiroma et al. in~\cite{Shiroma2022}, which is used in the \acs{gcov}-NISAR product~\cite{NISAR_GCOV_ProdSpec}, is likely to be smaller than the one required for some applications. The same will happen when the speckle filter applied before the geocoding of the \acs{pol} product defined in \acs{ceos}~\cite{CEOS-ARD_PFS_SAR} employs a small value of looks (e.g. a 3x3 or 5x5 boxcar). In both cases, the usually adopted solution consists in adding an extra step of spatial averaging to the previously-filtered covariance matrices and before the computation of the derived products. This extra filtering step, if applied after geocoding, will produce some artifacts as it was commented in \autoref{p:order_pol}.

The severity of the impact of these artifacts will depend on the specific case study and on the final application, but it must be taken into account.

\paragraph*{Faraday rotation}

Faraday rotation affects the distribution of backscatter intensity across the different polarimetric channels. Its effects can be compensated if quad-pol data are available. If not available, it adds up to the error budget and reduces the performance. Therefore,  single-, dual- and compact-polarimetric configurations become potentially problematic. Fortunately, in many occasions the rotation angle is small enough to not distort severely the polarimetric information. 


\subsubsection{Input and Auxiliary Data}
\label{sec:aux_data_pol}

The input and auxiliary data necessary to generate the polarimetric products coincide with the \acs{nrb} product, as detailed in \autoref{tab:in_aux_pol_products}. 

\begin{longtblr}[
	label = tab:in_aux_pol_products,
	entry = none,
	caption = {\acs{cov}/\acs{gcov} input and auxiliary data},
	]{
		width = \linewidth,
		colspec = {Q[l,20]Q[l,60]Q[c,20]},
		hlines,
		hline{1-Z} = {-}{0.05em},
	}
	\textbf{Name}   & \textbf{Description}   & \textbf{Mandatory (M) / Optional (O)} \\
	\acs{slc} or \acs{cslc} & L1 Single Look Complex image or co-registered L1 Single Look Complex image & M \\
    \acs{slc} annotation & L1 Single Look Complex image annotations or reference geometry annotations & M \\
    \acs{dem} & External DEM (e.g., Copernicus DEM at different resolutions, with 30 m as the default worldwide). & M \\
    Orbits & \acs{slc} Standard, restituted or precise orbits & M \\
    \acs{etad} & \acs{etad}-like products for atmospheric delay correction & O \\
    \acs{tec} & External \acs{tec} data for delay correction & O \\
\end{longtblr}


\subsubsection{Main Variables}
\label{sec:variables_pol}

\autoref{tab:var_pol_products} lists the main variables associated with the \acs{cov}/\acs{gcov} products.

\begin{longtblr}[
	label = tab:var_pol_products,
	entry = none,
	caption = {\acs{cov}/\acs{gcov} main variables},
	]{
		width = \linewidth,
		colspec = {Q[l,15]Q[l,52]Q[l,20]Q[l,13]},
		hlines,
		hline{1-Z} = {-}{0.05em},
	}
	\textbf{Name} & \textbf{Description} & \textbf{Dimensions} & \textbf{Data Type} \\
	gamma0 & Radiometric Terrain Corrected Gamma-Naught image (backscatter in linear scale) of each channel & polarization, acquisition\_date, x, y & float32 \\

    correlation & Radiometric Terrain Corrected cross-correlation image (in linear scale) between pairs of channels & polarization pair, acquisition\_date, x, y & complex64 \\
        
    local\_inc
    
    \_angle & DEM-based local Incidence angle image [deg]. This variable is present only if the \acs{cov}/\acs{gcov} is generated from \acs{slc}s. If it is derived from \acs{cslc}s, this variable will be null, and the reference\_inc\_angle metadata will be present & x, y & float32 \\
    
    anf & Optional variable: area normalisation factor or local contributing area, used to obtain the gamma0 coefficient. This variable is present only if the \acs{cov}/\acs{gcov} is generated from \acs{slc}s. If it is derived from \acs{cslc}s, this variable will be null, and the reference\_anf metadata will be present & x, y & float32 \\
\end{longtblr}

The dimensions and coordinates associated with the \acs{gcov} main variables are described in \autoref{tab:var_gcov_coordinates}, whereas for \acs{cov} they are listed in \autoref{tab:var_cov_coordinates}. 

\begin{longtblr}[
	label = tab:var_gcov_coordinates,
	entry = none,
	caption = {\acs{gcov} main variables coordinates/dimensions},
	]{
		width = \linewidth,
		colspec = {Q[l,20]Q[l,62]Q[l,18]},
		hlines,
		hline{1-Z} = {-}{0.05em},
	}
	\textbf{Name} & \textbf{Description} & \textbf{Data Type} \\
	polarization & Polarimetric layers (entries of the matrix) & string \\
    acquisition\_date & Image acquisition dates, formatted as YYYY-MM-DD & string \\
    x & UTM/UPS x (East) coordinates [m] & float64 \\
    y & UTM/UPS y (North) coordinates [m] & float64 \\
\end{longtblr}

\begin{longtblr}[
	label = tab:var_cov_coordinates,
	entry = none,
	caption = {\acs{cov} main variables coordinates/dimensions},
	]{
		width = \linewidth,
		colspec = {Q[l,20]Q[l,62]Q[l,18]},
		hlines,
		hline{1-Z} = {-}{0.05em},
	}
	\textbf{Name} & \textbf{Description} & \textbf{Data Type} \\
	polarization & Polarimetric layers (entries of the matrix) & string \\
    acquisition\_date & Image acquisition dates, formatted as YYYY-MM-DD & string \\
    azimuth\_time & Time along satellite track direction & datetime64 \\
    range\_time & Time across satellite track direction & float64 \\
\end{longtblr}

\subsubsection{Main Metadata}
\label{sec:metadata_pol}

As it has been commented above, both \acs{cov} and \acs{gcov} products share most characteristics with the \acs{nrb} products described in~\autoref{s:nrb}, either because they are considered as an extension of the \acs{nrb} (as in the \acs{pol} product defined by \acs{ceos}) or because \acs{nrb} is a subset of \acs{gcov} (in the \acs{gcov} definition by NASA for NISAR). Consequently, all the product metadata described in~\autoref{sec:nrb_product_metadata} apply also to \acs{cov} and \acs{gcov}. 

However, for completeness and for making the search for information easier, all the product metadata are included in~\autoref{tab:meta_gcov_products}. Key comments about specific metadata which need special attention in the polarimetric \acs{ard} products are detailed. 

\autoref{tab:meta_gcov_products} lists the main metadata and per-pixel metadata associated with the \acs{gcov} product.

\begin{longtblr}[
	label = tab:meta_gcov_products,
	entry = none,
	caption = {\acs{gcov} main metadata and per-pixel metadata},
	]{
		width = \linewidth,
		colspec = {Q[l,17]Q[l,38]Q[l,20]Q[l,13]Q[c,12]},
		hlines,
		hline{1-Z} = {-}{0.05em},
	}
	\textbf{Name} & \textbf{Description} & \textbf{Dimensions} & \textbf{Data Type} & \textbf{Per-pixel Metadata}\\
	reference\_slc / reference\_cslc & List of links to input reference \acs{slc} or \acs{cslc} & - & list of string & \\
    reference\_dem & Link to reference DEM & - & string & \\
    etad\_file & ETAD file used for each image belonging to the \acs{cslc}/\acs{slc} stack. This is optional: it is included only if used; otherwise, it is null & - & string & \\ 
    tec & Link to external \acs{tec} data, if used & - & string & \\ 
    reference\_inc
    
    \_angle & Optional attribute: list of link to reference DEM-based local incidence angle maps. It is present only if the \acs{gcov} is generated from \acs{cslc}s. & - & string & \\
    reference\_anf & Optional attribute: list of links to reference area normalisation factor or local contributing area. It is present only if the \acs{gcov} is generated from \acs{cslc}s. & - & string & \\
    orig\_data\_id & Per pixel metadata reporting an index describing the bursts contributing to the output mosaicked pixel & x, y & int32 & \ding{51} \\
    mask & Mask containing, for each ground cell, the following information:
    
     - 0: no distortions
    
     - 1: layover
    
     - 2: shadow
    
     - 3: undefined & x, y & byte & \ding{51} \\
    nlooks & Number of equivalent looks for the gamma flattening & x, y & float32 & \ding{51} \\

    filtering & If speckle filtering has been applied, metadata should include: 
    
- Reference to algorithm

- Input filtering parameters

- Type

- Window size

- Other filter parameters & - & string & \\ 
    
    reference\_inc

    pol\_calib
    
    \_matrix 
    & Polarimetric distortion matrices with the channel imbalance and the cross-talk applied for the polarimetric calibration. (Optional) & - & complex64  &   \\

    mean\_Faraday
    
    \_rotation\_angle 
    
    &  Estimated mean Faraday rotation angle, obtained from the polarimetric data and/or from models, with link to the method or publication used to derive the estimate. (Optional) & - & float32  &  \\

   ionosphere
   
   \_mask 
   &  Flag indicating whether the image is “significantly impacted” by the ionosphere (0 – false, 1 – true). Significant impact would imply that the ionospheric impact on the backscatter exceeds the radiometric calibration requirement for the image. (Optional) & - & byte  &  \\
    
    geometric
    
    \_accuracy
    & Measured geometric accuracy, including eastern and northern geometric accuracy bias and standard deviation (only for projected products) & - & dictionary & \\ 
    east\_geom
    
    \_bias & Measured eastern geometric bias (only for projected products) & - & float32 & \\
    east\_geom\_std & Measured eastern geometric standard deviation (only for projected products) & - & float32 & \\
    north\_geom
    
    \_bias & Measured northern geometric bias (only for projected products) & - & float32 & \\
    north\_geom\_std & Measured northern geometric standard deviation (only for projected products) & - & float32 & \\
    geometric
    
    \_accuracy
    
    \_radial
    
    \_rmse	& Measured radial RMSE of geometric accuracy (only for projected products) & - & float32 & \\
\end{longtblr}

And \autoref{tab:var_pol_metadata_coordinates} describes the dimensions/coordinates associated with the \acs{cov} and \acs{gcov} main metadata.

\begin{longtblr}[
	label = tab:var_pol_metadata_coordinates,
	entry = none,
	caption = {\acs{cov} and \acs{gcov} main metadata coordinates/dimensions},
	]{
		width = \linewidth,
		colspec = {Q[l,20]Q[l,62]Q[l,18]},
		hlines,
		hline{1-Z} = {-}{0.05em},
	}
	\textbf{Name} & \textbf{Description} & \textbf{Data Type} \\
	x & UTM/UPS x (East) coordinates [m] & float64 \\
    y & UTM/UPS y (North) coordinates [m] & float64 \\
        azimuth\_time & Time along satellite track direction & datetime64 \\
    range\_time & Time across satellite track direction & float64 \\
\end{longtblr}

In the case of the \acs{cov} product, it contains the same metadata as \acs{gcov} because the products are ready for geocoding in the same way as the \acs{gcov} product. The only additional metadata is a look-up table with the correspondence between each output pixel in cartographical coordinates and the pixels in radar coordinates.


\subsubsection{Reference System Tiling and Slicing}
\label{sec:gcov_reference_system_tiling}

\autoref{tab:pol_ref_system_slicing} reports the reference system, the slicing, and grid for \acs{cov} and \acs{gcov} products.

\begin{longtblr}[
	label=tab:pol_ref_system_slicing,
	entry = none,
	caption = {Reference system and Slicing for \acs{gcov}},
	]{
		width=\linewidth,
		colspec = {Q[l,15]Q[l,35]Q[l,35]Q[l,15]},
		%vline{2-5} = {-}{},
		%vline{5-5} = {-}{},
		hlines,
		hline{1-Z} = {-}{0.05em},
        rowhead=2
	}
		\textbf{Product} & \textbf{Reference System} & \textbf{Slicing} & \textbf{Spacing}\\
	\textbf{\acs{cov}} & Radar: azimuth/slant-range & Burst & Same of the input data \\
	\textbf{\acs{gcov}} & Cartographic: UTM/UPS & 100 km x 100 km & 10 m\\   
\end{longtblr}

In \autoref{tab:pol_ref_system_slicing}, the slicing for the \acs{cov} product is assigned as a burst as it is derived from \acs{slc} or \acs{cslc} products. However, the correspondence data necessary to convert it to \acs{gcov}, e.g. a look-up table with the association of radar coordinates corresponding to each cartographical coordinate, should be included as a metadata layer. After this conversion, the resulting product is really an \acs{ard} with total integrability and interoperability, like the \acs{gcov} product.

\subsection{Product Requirements}
\label{sec:impacts_pol}

The main requirements for a correct generation of \acs{cov} and \acs{gcov} products coincide with the ones stated for \acs{nrb} in \autoref{sec:impacts_nrb}, and the processing steps required to generate them coincide also with all the steps carried out to obtain \acs{nrb} products. 

All the processing steps must be applied in the same way (same algorithm and same parameters) to each individual entry of the polarimetric covariance matrix. It must be noted that the off-diagonal elements of the polarimetric matrix are complex numbers (not real numbers as the diagonal elements),  but they must be processed in the same way as the diagonal elements.



\subsubsection{Accuracy}

\begin{longtblr}[
	label = tab:pol_accuracy_cband,
	entry = none,
	caption = {Accuracy Requirements for \acs{cov} and \acs{gcov} products (C-Band)},
	]{
		width = \linewidth,
		colspec = {Q[l,30] Q[l,60] Q[c,20]},
		hlines,
		hline{1-Z} = {-}{0.05em},
		rowhead = 1,
	}
	
	\textbf{Metric} & \textbf{Description} & \textbf{Requirement} \\
	
	\SetCell[r=2]{l} Range Accuracy (\%) 
	& Model-Driven (only orbit and \acs{dem}) & $\leq 5\%$\\
	& Model-Driven + atmospheric correction & $\leq 2\%$ \\
	
	\SetCell[r=2]{l} Azimuth Accuracy (\%) 
	& Model-Driven (only orbit and \acs{dem}) & $\leq 2\%$\\
	& Data-Driven (\acs{esd}) & $\leq 0.5\%$ \\
	
	\acs{rcs} difference (dB) & Stability of the \acs{rcs} comparing different acquisition geometries over the same area considering a resolution of 30m & $\leq 5$~dB \\
	
\end{longtblr}

\begin{longtblr}[
	label = tab:pol_accuracy_lband,
	entry = none,
	caption = {Accuracy Requirements for \acs{cov} and \acs{gcov} products (L-Band)},
	]{
		width = \linewidth,
		colspec = {Q[l,30] Q[l,60] Q[c,20]},
		hlines,
		hline{1-Z} = {-}{0.05em},
		rowhead = 1,
	}
	
	\textbf{Metric} & \textbf{Description} & \textbf{Requirement} \\
	
	\SetCell[r=3]{l} Range Accuracy (\%) 
	& Model-Driven (only orbit and \acs{dem}) & $\leq 60\%$\\
	& Model-Driven + atmospheric correction & $\leq 25\%$ \\
	& Model-Driven + atmospheric correction + Data-Driven (\acs{icc}) & $\leq 10\%$ \\
	
	\SetCell[r=3]{l} Azimuth Accuracy (\%) 
	& Model-Driven (only orbit and \acs{dem}) & $\leq 150\%$ \\
    & Model-Driven + atmospheric correction & $\leq 50\%$ \\
	& Data-Driven (\acs{icc}) & $\leq 10\%$ \\
	
	\acs{rcs} difference (dB) & Stability of the \acs{rcs} comparing different acquisition geometries over the same area considering a resolution of 30m & $\leq 5$~dB \\
	
\end{longtblr}


\subsubsection{Timeliness}

The \acs{cov} and \acs{gcov} products can be generated as soon as L1 data is available for each \acs{slc} image independently or starting from \acs{cslc} products. The generation of covariance matrix for the provision of the polarimetric products does not impact significantly the timeline. 

Besides choosing between \acs{slc} or \acs{cslc} products, the main factor impacting the timeline is the availability of external data or reference images, which are required to refine geolocation and mitigate atmospheric effects, particularly ionospheric effects. 


\autoref{fig:pol_single} shows the flowchart of generating the \acs{cov} and \acs{gcov} products using a single \acs{slc} image. The associated timeline was provided in~\autoref{fig:nrb_single}, indicating the time required to collect the input data and complete the main processing steps. 

\begin{figure}[H]
	\centering
	\resizebox{\textwidth}{!}{
		\begin{tikzpicture}[node distance=1cm and 2cm]
			
			% Define the style for the nodes
			% \tikzstyle{node} = [circle, draw, fill=blue!20, text centered, minimum size=1.5cm]
			
			% Draw nodes
			\node[rectin] (init) {\acs{slc}};
			\node[block, right=of init] (g2r) {Geo-to-Radar};
            \node[block, right=of g2r] (rtc) {\acs{rtc}};
			\node[block, right=of rtc] (covm) {Covariance Matrix};
			\node[block, right=of covm] (geocoding) {Geocoding};
			\node[rectaux, below=of init] (dem) {\acs{dem}};
			\node[rectaux, below=of dem] (orbit) {Orbital data};
			\node[rectauxd, below=of orbit] (models) {External models};
			\node[rectind, below=of models] (ref_stack) {Reference SLC stack};
			\node[rectout, right=of geocoding] (gcov) {\acs{gcov}};
			\node[rectout, below=of covm] (cov) {\acs{cov}};
			
			% Draw edges with arrows
			\draw[->, thick] (init) -- (g2r);
			\draw[->, thick] (g2r) -- (rtc);
			\draw[->, thick] (rtc) --(covm);
			\draw[->, thick] (covm) -- (geocoding);
			\draw[->, thick] (covm) -- (cov);
			\draw[->, thick] (dem) -| (g2r);
			\draw[->, thick] (orbit) -| (g2r);
			\draw[->, thick] (models) -| (g2r);
			\draw[->, thick] (ref_stack) -| (g2r);
			\draw[->, thick] (geocoding) -- (gcov);
		\end{tikzpicture}
	}
	\caption{Sample workflow for generation of \acs{cov} and \acs{gcov} considering each \acs{slc} image separately. The step denoted as "Covariance Matrix" should include a speckle filtering, but is declared as optional. In the \acs{gcov} product defined by NISAR the speckle filtering and the geocoding are carried out simultaneously with an area-adaptive algorithm. The associated timeline can be consulted in the corresponding figure for \acs{nrb} (\autoref{fig:nrb_single}).}
	\label{fig:pol_single}
\end{figure}


In case the \acs{cov} or \acs{gcov} products are generated from the \acs{cslc}, all the calibrations have already been performed during the \acs{cslc} processing, and the reference geometry is computed only once for the reference orbit. Therefore, the schema simplifies as reported in \autoref{fig:pol_from_cslc}, but the timeline does not change depending on the information needed to perform Geo-to-Radar conversion or to generate the \acs{cslc} product.
\begin{figure}[H]
	\centering
	\resizebox{\textwidth}{!}{
		\begin{tikzpicture}[node distance=0.5cm and 2cm]
			
			% Define the style for the nodes
			% \tikzstyle{node} = [circle, draw, fill=blue!20, text centered, minimum size=1.5cm]
			
			% Draw nodes
			\node[rectin] (init) {\acs{slc}};
			\node[block, right=of init] (coreg) {Co-registration};
			\node[rectout, right=of coreg] (cslc) {\acs{cslc}};
            \node[block, right=of cslc] (rtc) {\acs{rtc}};
			\node[block, right=of rtc] (covm) {Covariance Matrix};
			\node[block, right=of covm] (geocoding) {Geocoding};
			\node[rectout, right=of geocoding] (gcov) {\acs{gcov}};
			\node[rectout, below=of covm] (cov) {\acs{cov}};
			
			\node[rectaux, below=of init] (dem) {\acs{dem}};
			\node[rectaux, below=of dem] (orbit) {Orbital data};
			\node[rectauxd, below=of orbit] (models) {External models};
			\node[rectind, below=of models] (ref_stack) {Reference SLC stack};
			
			% Draw edges with arrows
			\draw[->, thick] (init) -- (coreg);
			\draw[->, thick] (coreg) -- (cslc);
			\draw[->, thick] (cslc) -- (rtc);
			\draw[->, thick] (rtc) -- (covm);
			\draw[->, thick] (covm) -- (cov);
			\draw[->, thick] (covm) -- (geocoding);
			\draw[->, thick] (geocoding) -- (gcov);
			
			% Draw edges with arrows
			\draw[->, thick] (dem) -| (coreg);
			\draw[->, thick] (orbit) -| (coreg);
			\draw[->, thick] (models) -| (coreg);
			\draw[->, thick] (ref_stack) -| (coreg);
			
		\end{tikzpicture}
	}
	\caption{Sample workflow for generation of \acs{cov} and \acs{gcov} starting from \acs{cslc} products. The step denoted as "Covariance Matrix" should include a speckle filtering, but is declared as optional. In the \acs{gcov} product defined by NISAR the speckle filtering and the geocoding are carried out simultaneously with an area-adaptive algorithm. The associated timeline can be consulted in the corresponding figure for \acs{nrb} (\autoref{fig:nrb_from_cslc}).}
	\label{fig:pol_from_cslc}
\end{figure}

\subsection{Usability, Integrability and Interoperability}
\label{sec:usability_pol}

The characteristics of the \acs{gcov} products in terms of usability, integrability and interoperability are coincident, and extended, with respect to the \acs{nrb} products, which were detailed in \autoref{tab:nrb_usability_integrability_interoperability}. This is so because the first obvious way to use polarimetric products is by considering only the diagonal values of the covariance matrix. This strategy is equivalent to working with a \acs{nrb} product with several layers of intensity data (one layer per polarimetric channel). As such, all the comments about the usage of \acs{nrb} products (see \autoref{sec:nrb_usability_integrability}) apply here. 

However, it is when the matrix form is considered as a joint feature or entity when polarimetry really excels in many application domains. Therefore, the specific usability, integrability and interoperability of these matrix-form products is described next in \autoref{tab:gcov_usability_integrability_interoperability}. 

\begin{longtblr}[
  caption = {Usability, Integrability, and Interoperability of \acs{gcov} product},
  label = {tab:gcov_usability_integrability_interoperability},
]{
  colspec = {Q[2cm] Q[12cm]},
  rowhead = 1,
  width = \linewidth
}
\hline
\textbf{Aspect} & \textbf{Description} \\
\hline
Usability &
\begin{itemize}
  \item Radiometric calibration to $\gamma^0$ convention enables direct comparison across sensors, acquisition times, and geometries.
  \item Orthorectified and georeferenced product makes data ready to use in thematic mapping applications (classification and change detection). 
  \item Polarimetric features: Besides the direct use of the entries of the covariance matrix, many polarimetric features can be derived from them. The most clear example is the outputs of polarimetric decompositions~\cite{Lee_Pottier2009}. Starting from \acs{gcov} products, these features can be readily computed.
  \item Machine Learning: \acs{gcov} products, as well as decomposition outputs, can be directly ingested into machine learning models for classification, regression, or segmentation tasks without the need for extra processing. The input features to these models are usually expressed as layers or bands of data, which in this case would be formed by the entries of the covariance matrix and the outputs of decompositions. In the case the machine learning algorithm is prepared only for intensity data, there exist simple expressions which allow bijectively converting the complex matrix into a set of intensities, and vice versa~\cite{Hoekman2003, Mestre2024}.
\end{itemize}
\\

Integrability &
\begin{itemize}
  \item GIS Compatibility: Easily used within \acs{gis} platforms, due to geographic referencing and data format.
  \item Multi-Sensor Fusion: Integrates well with optical or other SAR projected products for land cover, surface change, or hazard mapping.
  \item Environmental Monitoring: Supports use in models or indicators for forest cover, crop health, or hydrology. 
  \item Users can overlay the matrix elements contained in these products with vector layers, other remote sensing data (optical, infrared), or ground truth data for comprehensive spatial analysis.
\end{itemize}
\\

Interoperability &
\begin{itemize}
  \item Cross-Mission Compatibility: Normalized $\gamma^0$ used in \acs{gcov} allows for comparison between satellites.
  \item Interdisciplinary Applications: Usable in hydrology, agriculture, forestry, urban studies, etc., enhancing cross-domain use.
  \item Compatible with common \acs{eo} frameworks and pipelines using georeferenced data.
  \item High interoperability due to standard geographic referencing and pixel alignment across scenes.
  \item Adheres to ARD principles, enhancing integration with downstream services and platforms
  \item Facilitates cross-mission usage through harmonized ground projection.
\end{itemize}
\\
\hline
\end{longtblr}






\section{Ocean products (ORB)}
\label{sec:orb_product}

\subsection{Product Description}
\label{sec:orb_prod_description}
The \acs{orb} product, either as defined in \cite{CEOS-ARD_PFS_ORB} or in the variants outlined in the proposal for this study, can be relatively straightforwardly derived from the \acs{nrb}. This derivation provides both the necessary full resolution and spatially averaged \gls{sigma0} values on a map-projected grid. 

\acs{orb} comprises multiple layers of \gls{sigma0}, each corresponds to a specific layer of preprocessing. These preprocessing steps include denoise, outlier removal, downsampling and surface roughness normalisation. The raw (full resolution) and intermediate products are also saved, including the noise-equivalent \gls{sigma0} (NESZ), a bright target detection mask (full resolution), and the normalization factor used for surface roughness correction. Comprehensive annotations detailing the processing steps applied to each layer ensures full traceability and allows users to reverse or reproduce the processing chain as needed.

A full resolution product is useful for the detection of ships or other anthropogenic targets, which could be done starting from the NRB product. Providing a \gls{sigma0} product is a logical choice when treating the data as an image of a locally homogeneous distributed surface scatterer. It also is worth noting that a full-resolution \acs{orb} product is, in general, not a good starting point for the generation of enhanced ocean products. In many cases, the generation of these products requires SLCs (or even raw data), for example for sub-look generation to enable the computation of intensity cross-spectra, or for the optimal suppression of azimuth ambiguities in coastal areas by optimal selection of the azimuth-wavenumber (also known as Doppler frequency) weighting. 

For applications where the NRCS is relevant (e.g. surface stress or wind, oil-spill detection), a low-speckle spatially averaged product including polarimetric information, and a compensation of the known part of the angle of incidence dependence is preferable to a full-resolution product which first needs to be properly despeckled, noting that speckle filtering requires a level of expertise that makes the product deviate from the ARD rationale. Moreover, given the limitation of the effective resolution by the azimuth cut-off, which is largely wavelength independent, a spatially averaged product to a generally achievable spatial resolution facilitates the creation of intercomparable products by different sensors.

% The \acs{ceos} \acs{ard} Ocean Radar Backscatter (ORB) product is a full-resolution intensity product with the necessary corrections applied to allow a nominal interpretation as a Normalized Radar Cross Section (NRCS, $\sigma_0$). The \acs{ceos} specification of the ORB product, closely related to the NRB product, includes geoid correction, to convert to $\sigma_0$ but it does not include any despeckling. This choice is questionable given that the azimuth resolution of a SAR image is limited not by the radar characteristics or the processed Doppler bandwidth, as over land, but by the coherence time of the surface as observed by the radar. This results in the so-called azimuth cut-off,
% \begin{equation}
%         \lambda_c \approx \frac{R\cdot \sigma_v}{v_\mathrm{orb}},
%     \end{equation}
% with $R$ the slant-range to the surface and $v_\mathrm{orb}$ the orbital velocity of the satellite. The result is that the effective azimuth resolution of the intensity image is limited to \qtyrange{50}{400}{\meter}, depending on the sea state, which, in turn, is mostly wind-driven. Most notably, this effective resolution is largely wavelength independent, which would, in principle, facilitate the creation of intercomparable products by different sensors, if these are first averaged to a generally achievable spatial resolution.

% A full resolution product is still useful for the detection of ships or other anthropogenic targets, for which the $\sigma_0$ convention is not particularly useful. Ship detection could be done starting from the NRB product. Providing a $\sigma_0$ product is a logical choice when treating the data as an image of a locally homogeneous distributed surface scatterer. %Given the limitation of the effective resolution by the azimuth cut-off, it is surprising that the ORB product is defined as a full nominal resolution product. 
% For applications where the NRCS is relevant (e.g. surface stress or wind, oil-spill detection), a low-speckle spatially averaged product including polarimetric information, and a compensation of the known part of the angle of incidence dependence is preferable to a full-resolution product which first needs to be properly despeckled, noting that speckle filtering requires a level of expertise that makes the product deviate from the ARD rationale.

% It also is worth noting that a full-resolution \acs{orb} product is, in general, not a good starting point for the generation of enhanced ocean products. In many cases, the generation of these products requires SLCs (or even raw data), for example for sub-look generation to enable the computation of intensity cross-spectra, or for the optimal suppression of azimuth ambiguities in coastal areas by optimal selection of the azimuth-wavenumber (AKA Doppler frequency) weighting.


% The \acs{orb} product, either as defined in \cite{CEOS-ARD_PFS_ORB} or in the variants discussed in the proposal for this study, can be relatively straightforwardly derived from the \acs{nrb}.

% While many references can be found to prototype implementations of \acs{nrb} products for both Sentinel-1 and ROSE-L, there is limited evidence of similar efforts for \acs{orb}, although \acs{orb}-type products are likely generated as intermediate products by many stakeholders in their analysis of ocean data.

% Lower resolution ocean \textit{roughness} products derived from Sentinel-1 data are, in contrast, provided (for example) through the through Oceandatalab's Ocean Virtual Lab portal.

% \begin{longtblr}[
%   label = tab:orb_products,
%   entry = none,
%   caption = {Data layers of the \acs{orb} },
% ]{
%   width = \linewidth,
%   colspec = {Q[c,20]Q[c,30]Q[l,40]},
%   hlines,
%   hline{1-Z} = {-}{0.05em},
% }
% \textbf{\acs{ard} Product} & \textbf{Description} & \textbf{Notes}\\
% $\sigma_{0,1}$ & Full resolution (single look) radar intensity, in $\sigma_0$ convention & This is the default \acs{orb} product defined in  \cite{CEOS-ARD_PFS_ORB} \\
% $\sigma_{0,L}$ & Spatially-averaged, LxL m$^2$ resolution radar intensity, in $\sigma_0$ convention. & The possibility of multilooked product is implicit in the \acs{ceos} \acs{orb} definition, as it defines metadata regarding speckle filtering, etc. \\
% $\sigma_{0,L}/\sigma_{0,\mathrm{mod}}$ & (Mean) wind normalized backscatter measurements & This product is not required, but it defined in \cite{CEOS-ARD_PFS_ORB} and highly useful as it removes the geometry dependence of the backscatter intensity, allowing a direct interpretation of atmospheric phenomena. The complementary metadata also allow a low-effort wind-retrieval. //
% 1 & $C_{3mod}$(1,1) & HH backscatter & Real  \\  
% 2 & $C_{3mod}$(1,2) & HH·HV* correlation & Complex  \\
% 3 & $C_{3mod}$(1,3) & HH·VV* correlation & Complex \\
% 4 & $C_{3mod}$(2,2) & HV backscatter & Real \\
% 5 & $C_{3mod}$(2,3) & HV·VV* correlation & Complex \\
% 6 & $C_{3mod}$(3,3) & VV backscatter & Real \\
%\end{longtblr}

\subsubsection{Input and Auxiliary Data}
\label{sec:aux_data_ocean}
% To generate the product the following auxiliary data are necessary:
% \begin{itemize}
%     \item external Geoid
%     \item \acs{slc} precise orbits
%     \item \acs{etad}-like product for atmospheric delay and timing corrections.
%     \item \acs{slc} ancillary data
% \end{itemize}

To generate the \acs{orb} product coincide with the \acs{nrb} product, the following input and auxiliary data are necessary: 
\begin{longtblr}[
	label = tab:in_aux_orb_products,
	entry = none,
	caption = {\acs{orb} input and auxiliary data},
	]{
		width = \linewidth,
		colspec = {Q[l,20]Q[l,60]Q[c,20]},
		hlines,
		hline{1-Z} = {-}{0.05em},
	}
	\textbf{Name}   & \textbf{Description}   & \textbf{Mandatory (M) / Optional (O)} \\
	\acs{slc} or \acs{cslc} & L1 Single Look Complex image or co-registered L1 Single Look Complex image and annotations & M \\
    \acs{dem} & External DEM (e.g., Copernicus DEM at different resolutions, with 30 m as the default worldwide). & O \\
    Orbits & \acs{slc} Standard, restituted or precise orbits & M \\
    \acs{etad} & \acs{etad}-like products for atmospheric delay and timing corrections & O \\
    \acs{tec} & External \acs{tec} data for delay correction & O \\
    Geoid & External geoid & M \\
    \acs{era5} & ECMWF \acs{era5} wind data for angle of incidence compensation & M \\
\end{longtblr}

\subsubsection{Main Variables}
\label{sec:variables_orb}
The main variables of the \acs{orb} product are listed as follows:
\begin{longtblr}[
	label = tab:var_orb_products,
	entry = none,
	caption = {\acs{orb} main variables},
	]{
		width = \linewidth,
		colspec = {Q[l,15]Q[l,52]Q[l,20]Q[l,13]},
		hlines,
		hline{1-Z} = {-}{0.05em},
	}
	\textbf{Name} & \textbf{Description} & \textbf{Dimensions} & \textbf{Data Type} \\
    $\sigma_{0,\mathrm{raw}}$ & Full resolution (single look) radar intensity, in \gls{sigma0} convention & polarization, acquisition\_date, x, y & float32 \\
    $\sigma_{0,\mathrm{nesz}}$ & Spatially-averaged, LxL m$^2$ resolution noise equivalent sigma zero, in \gls{sigma0} convention & polarization, acquisition\_date, x, y & float32 \\
    $\sigma_{0,\mathrm{denoise}}$ & Spatially-averaged, LxL m$^2$ resolution denoized radar intensity, in \gls{sigma0} convention & polarization, acquisition\_date, x, y & float32 \\
    mask & Full resolution (single look) bright target mask & polarization, acquisition\_date, x, y & bool \\
    $\sigma_{0,L}$ & Spatially-averaged, LxL m$^2$ resolution denoized and outlier removed radar intensity, in \gls{sigma0} convention & polarization, acquisition\_date, x, y & float32 \\
    $\sigma_{0,L}/\sigma_{0,\mathrm{mod}}$ & Spatially-averaged, LxL m$^2$ resolution (mean) wind normalized backscatter measurements & polarization, acquisition\_date, x, y & float32 \\
\end{longtblr}
The table below describes the dimensions/coordinates associated with the \acs{orb} main variables.
\begin{longtblr}[
	label = tab:var_orb_coordinates,
	entry = none,
	caption = {\acs{orb} main variables coordinates/dimensions},
	]{
		width = \linewidth,
		colspec = {Q[l,20]Q[l,62]Q[l,18]},
		hlines,
		hline{1-Z} = {-}{0.05em},
	}
	\textbf{Name} & \textbf{Description} & \textbf{Data Type} \\
	polarization & Acquired polarizaions & string \\
    acquisition\_date & Image acquisition dates, formatted as YYYY-MM-DD & string \\
    x & UTM/UPS x (East) coordinates [m] & float64 \\
    y & UTM/UPS y (North) coordinates [m] & float64 \\
\end{longtblr}
\subsubsection{Main Metadata}
\label{sec:orb_product_metadata}
The following table lists the main metadata and per-pixel metadata associated with the \acs{orb} product.

\begin{longtblr}[
	label = tab:meta_orb_products,
	entry = none,
	caption = {\acs{orb} main metadata and per-pixel metadata},
	]{
		width = \linewidth,
		colspec = {Q[l,17]Q[l,38]Q[l,20]Q[l,13]Q[c,12]},
		hlines,
		hline{1-Z} = {-}{0.05em},
	}
	\textbf{Name} & \textbf{Description} & \textbf{Dimensions} & \textbf{Data Type} & \textbf{Per-pixel Metadata}\\
	reference\_slc / reference\_cslc & List of
    links to input reference \acs{slc} or \acs{cslc} & - & list of string & \\
    reference\_geoid & Link to reference GEOID. Only for \acs{eorb} & - & string & \\
    tec & Link to external \acs{tec} data, if used. & - & string & \\ 
    era5 & Link to external \acs{era5} data, if used.  & - & string & \\
        x\_sample & Pixel spacing in x direction & - & float32 &\\
        y\_sample & Pixel spacing in y direction & - & float32 &\\
        ground\_heading & Azimuth direction & - & float32 & \\
        look\_direction & Planar angle between	north and each range direction & - & float32 &\\
        inc\_angle & Angle of incidence & x, y & float32 & \ding{51} \\
        nlooks & Number of equivalent looks & x, y & float32 & \ding{51} \\
        azimuth\_time & Azimuth time & x, y & datetime64 & \ding{51} \\
        slant\_range\_time & Slant range time & x, y & datetime64 & \ding{51} \\
\end{longtblr}

The table below describes the dimensions/coordinates associated with the \acs{orb} main metadata.

\begin{longtblr}[
	label = tab:var_orb_metadata_coordinates,
	entry = none,
	caption = {\acs{orb} main metadata coordinates/dimensions},
	]{
		width = \linewidth,
		colspec = {Q[l,20]Q[l,62]Q[l,18]},
		hlines,
		hline{1-Z} = {-}{0.05em},
	}
	\textbf{Name} & \textbf{Description} & \textbf{Data Type} \\
	x & UTM/UPS x (East) coordinates [m] & float64 \\
    y & UTM/UPS y (North) coordinates [m] & float64 \\
\end{longtblr}

% The indicated main metadata for these products derive from the analysis of \acs{ceos} Product Format Specification \cite{CEOS-ARD_PFS_SAR}.% and of the metadata available information for the products list in \autoref{tab:orb_products}.

% Most of the listed products are declared compliant with CEOS requirements whereby it is assumed that their metadata include at least the threshold requirements of \cite{CEOS-ARD_PFS_ORB} and \cite{CEOS-ARD_PFS_SAR}.

% In \autoref{tab:metadata_orb_products}, ``All" means that the correspondent metadata are present in all the products in \autoref{tab:orb_products} and in CEOS recommendations. %Note that NISAR \acs{rtc} is included in the NISAR \acs{gcov} product. 
% This table can be understood as a \acs{orb}-specific revision of \autoref{tab:metadata_nrb_products}.

% \begin{longtblr}[
%   label = tab:metadata_orb_products,
%   entry = none,
%   caption = {ORB product main metadata from \acs{sota} analysis},
% ]{
%   width = \linewidth,
%   colspec = {Q[l,15]Q[l,42]Q[l,18]Q[l,25]},
%   hlines,
%   hline{1-Z} = {-}{0.05em},
%   rowhead=1,
% }
% \textbf{Metadata name}   & \textbf{Description}   & \textbf{Reference Product} & \textbf{Notes} \\
% Product Type & Product type name & All & - \\
% Source data attributes & Subsection describing \acs{sar} acquisitions used to generate the product. Main attributes for each source are:

% -Source Data Access

% -Instrument

% - Acquisition Time

% - Acquisition Parameters (Band, Centre frequency, Observation mode, Polarization(s), Antenna pointing, Beam ID)

% - Orbit Information

% - Processing Parameters

% - Image Attributes (geometry, azimuth and range pixel spacing and resolution, incidence angles, range distance from the sensor to the surface)

% & All & - \\
% Product Data Access & Processing parameters details of the product (processing facility, processing date, software version, \acs{url} or \acs{doi} from where product can be retrieved.
%  & All & - \\
% Polarization(s) & Indication of the polarization. & All & - \\
% Product Sample Spacing & - Pixel (column) spacing

% -	Line (row) spacing
%  & All & - \\
% Product Equivalent Number of Looks & Equivalent Number of Looks (\acs{enl}) &  & Not required by \acs{ceos} (listed as ”goal”). 
%  \\
% Product Resolution & Average spatial resolution of the product along columns and rows
%  & - & Not required by \acs{ceos} (listed as ”goal”).\\

% Range of validity & Variable range of validity &  & The analysed documentation does not explicitly mention the presence of this field for all the other indicated products. \\
% Product Filtering & Flag indicating if speckle or noise filter has been applied. Metadata should include reference to algorithm as \acs{doi} or \acs{url} and input filtering parameters. & All & As threshold. Not relevant in \acs{ceos} \acs{orb} specification, but necessary if we consider, as strongly recommended, lower resolution speckle-filtered products. \\
% Product Bounding Box & Two opposite corners of the product file (bounding box, including any zero-fill values) are identified, expressed in the coordinate reference system.
% Four corners of the product file are recommended for scenes crossing the Antimeridian, or the North or the South Pole. 
%  & All & - \\
% Product Geographical Extent & The geometry of the \acs{sar} image footprint expressed in WGS84, in a standardised format (e.g., WKT Polygon). & All & - \\
% Product Image Size & 
% - Number of lines

% - Number of pixels per line

% - File header size (if applicable)

% - Number of no-data border pixels (if appl.)
%  & All & - \\
% Backscatter Measurements & "Geoid-corrected"  Sigma-Nought ($\sigma^0$) convention, for each polarization.
% Metadata should include:

% - Measurement Type: Sigma-Nought (NRCS)

% - Backscatter Expression Convention: Linear amplitude or linear power

% - Polarization: HH, HV, VV, VH

% - Data Format: GeoTIFF, HDF5, NetCDF, etc.

% - Data Type: Integer, Float, etc.

% - Bits per Sample

% - Byte Order

% & CEOS \acs{orb} & Partially present in all the listed products. \\
% Scaling Conversion & Equation to convert linear amplitude/power to logarithmic decibel scale, including, if applicable, the associated calibration (dB offset) factor, and/or the equation used to convert compressed data (int8/int16/float16) to float32. &  CEOS \acs{orb} & - \\

% GMF & Geophysical model function used to relate model winds to reference NRCS, $\sigma_{0,\mathrm{mod}}$. & & This could be provided as a mathematical model with or as a reference and URL to a GMF description and implementation. \\

% FR method & Method used to estimate FR. & &  Under some assumptions, FR can be estimated from data, or, alternatively, from ionospheric models. \\
% %
% % \acs{rtc} Algorithm & Reference to \acs{rtc} algorithm and sources of used auxiliary data. & All & - \\
% % Product Pixel Coordinate Convention & Coordinate referring to the Centre, or the Upper Left Corner or the Lower Left Corner of a pixel & All & - \\
% Product Coordinate Reference System & Parameters required to geolocate data in that map projection, expressed in a standardised format (e.g., WKT). 

% Indicate \acs{epsg} code, if defined for the \acs{crs}
%  & All & - \\
% Reference Orbit & Absolute orbit number of reference \acs{slc}. It includes the baseline for the source data relative to the reference orbit used (for approximate use only).
%  & All & The generation of \acs{orb}s is on single \acs{slc}. \\
% Geometric Correction Algorithms & Information on the algorithm used for geometric correction & - & Not required by \acs{ceos} \\
% % Digital Elevation Model & - Reference to Digital Elevation Model used for co-registration step, ortho-rectification and \acs{rtc} procedure.

% % - Reference to Earth Gravitational Model (\acs{egm}) used for geometric correction.
% %  & All & All steps, requiring it, use the same \acs{dem}.\\
% Radiometric Accuracy & Uncertainty on  $\sigma^0$. & - & Not required by \acs{ceos} (listed as ”goal”). \\
% Geometric Accuracy & Geolocation accuracy estimates, based on SLC assessment (\acs{slc} \acs{ale}) and/or measured directly in the final projected final product (\acs{ard} \acs{ale}). &  & Not required by \acs{ceos} (listed as ”goal”). 

% %For QA4EO RTC, the geometric accuracy does not include the relative RMSE and reference.
% \\
% Gridding Convention & All necessary information of used grid and used convention & All & - \\
% Processing information & Information on the processing: common band filtering, spatial averaging method, \acs{rfi} correction, co-registration algorithm, etc.. & All & Some of this information could be already available through source-data
% pointers. \\
% Units & Units of the variables of measurements. & All & Not found in \acs{ceos}.
% %\centering
% \end{longtblr}

% \autoref{tab:per-pixel_metadata_orb_products} indicates the main product per-pixel metadata. The names and definitions of the majority of the per-pixel metadata are exacted from \acs{ceos} Product Format Specification \cite{CEOS-ARD_PFS_SAR}.% and of the per-pixel metadata available information for the products list in \autoref{tab:orb_products}.

% Most of the listed products are declared compliant with CEOS requirements whereby it is assumed that their per-pixel metadata includes at least the threshold requirements of \cite{CEOS-ARD_PFS_ORB} and \cite{CEOS-ARD_PFS_SAR}.

% In \autoref{tab:per-pixel_metadata_orb_products}, "All" means that the correspondent per-pixel metadata is present in all the products in \autoref{tab:orb_products} and in CEOS recommendations. %Note that NISAR \acs{rtc} is included in the NISAR \acs{gcov} product.

% \begin{longtblr}[
%   label = tab:per-pixel_metadata_orb_products,
%   entry = none,
%   caption = {\acs{orb} product main per-pixel metadata from \acs{sota} analysis},
% ]{
%   width = \linewidth,
%   colspec = {Q[l,15]Q[l,33]Q[l,18]Q[l,34]},
%   hlines,
%   hline{1-Z} = {-}{0.05em},
% }
% \textbf{Metadata name}   & \textbf{Description}   & \textbf{Reference Product} & \textbf{Notes} \\
% Data mask & Mask image indicating:

% - Valid data

% - Invalid data

% - No data & All & - \\
% Equivalent Number of Looks Image & Equivalent Number of Looks (\acs{enl}) &  & Not required by \acs{ceos}, but important for characterization of radiometric quality for optionally spatially-averaged products. \\
% Scattering Area & geoid-based scattering area image &  & Not required by \acs{ceos} (indicated as a goal requirement). It depends on the used algorithm, e.g. if number of independent input pixels is constant this areas is going to vary for each output pixel. \\
% Local incidence angle & DEM based local incidence angle image & All &  \\
% Ellipsoidal incidence angle & Ellipsoidal incidence angle image & %NISAR RTC (GCOV), QA4EO RTC & Not required by \acs{ceos} (indicated as goal requirement). The analysed ROSE-L documentation does not explicitly mention the presence of this field.
% \\
% % Per-pixel \acs{dem} and per-pixel GEOID & DEM and/or GEOID using for the processing & CEOS \acs{nrb} & This can be useful for a variety of operations (undoing the geometric correction, the atmospheric correction, etc.). Probably it is not necessary report the DEM/GEOID for each pixel of the product: it could be sufficient the information from "Digital Elevation Model" metadata. \\
% Noise Power Image & Estimated noise equivalent $\sigma^0$  & CEOS \acs{orb} & - \\
% % Gamma-to-Sigma Ratio Image & Ratio between the integrated area in the Gamma projection and the integrated area in the Sigma projection (on the ground) & All & - \\
% Radiometric Accuracy & Estimated uncertainty (e.g. $1\sigma$) of the intensity products. & CEOS \acs{orb} & Not required by CEOS but important for some applications, for example for uncertainty propagation to surface winds estimates \\
% $U_{10s,m}$ & Stress-equivalent model 10-m wind used to generate reference intensity, $\sigma_{0,\mathrm{mod}}$  & - & -\\
% Faraday Rotation & Estimated FR per pixel & - & Not specified in CEOS \acs{orb}, but potentially relevant to correct or address leakage of co-polar signal in cross-polar channel.


% %\centering
% \end{longtblr}

\subsubsection{Reference System Tiling and Slicing}
A gridding/sampling frame, consistent with the reference system and slicing defined for the \acs{nrb}, is generally used.
\begin{longtblr}[
	label=tab:orb_ref_system_slicing,
	entry = none,
	caption = {Reference system and Slicing for \acs{orb}},
	]{
		width=\linewidth,
		colspec = {Q[l,15]Q[l,35]Q[l,35]Q[l,15]},
		%vline{2-5} = {-}{},
		%vline{5-5} = {-}{},
		hlines,
		hline{1-Z} = {-}{0.05em},
        rowhead=2
	}
	\textbf{Product} & \textbf{Reference System} & \textbf{Slicing} & \textbf{Spacing}\\
	\textbf{\acs{orb}} (default) & Cartographic: UTM/UPS & US-\acs{mgrs} 100 km x 100 km & 10 m\\
    \textbf{\acs{orb}} (optional) & Cartographic: UTM/UPS & US-\acs{mgrs} 100 km x 100 km & 30 m\\	
\end{longtblr}
\subsection{Product Requirements}
\subsubsection{Accuracy}
Main impacts on \acs{orb} product generation and main requirements are similar to those indicated in the correspondent section of the \acs{nrb} product. The limitations, in this case, are less stringent for the local slope (here the geoid is used) than for the geolocation accuracy (less binding over ocean).
\begin{longtblr}[
	label = tab:orb_accuracy_cband,
	entry = none,
	caption = {Accuracy Requirements for \acs{orb} product (C-Band)},
	]{
		width = \linewidth,
		colspec = {Q[l,30] Q[l,60] Q[c,20]},
		hlines,
		hline{1-Z} = {-}{0.05em},
		rowhead = 1,
	}
	
	\textbf{Metric} & \textbf{Description} & \textbf{Requirement} \\
	
	\SetCell[r=2]{l} Range Accuracy (\%) 
	& Model-Driven (only orbit and \acs{dem}) & $\leq 5\%$\\
	& Model-Driven + atmospheric correction & $\leq 2\%$ \\
	
	\SetCell[r=2]{l} Azimuth Accuracy (\%) 
	& Model-Driven (only orbit and \acs{dem}) & $\leq 2\%$\\
	& Data-Driven (\acs{esd}) & $\leq 0.5\%$ \\
	
	\acs{rcs} difference (dB) & Stability of the \acs{rcs} comparing different acquisition geometries over the same area considering a resolution of 30m & $\leq 5$~dB \\
	
\end{longtblr}

\begin{longtblr}[
	label = tab:orb_accuracy_lband,
	entry = none,
	caption = {Accuracy Requirements for \acs{orb} product (L-Band)},
	]{
		width = \linewidth,
		colspec = {Q[l,30] Q[l,60] Q[c,20]},
		hlines,
		hline{1-Z} = {-}{0.05em},
		rowhead = 1,
	}
	
	\textbf{Metric} & \textbf{Description} & \textbf{Requirement} \\
	
	\SetCell[r=3]{l} Range Accuracy (\%) 
	& Model-Driven (only orbit and \acs{dem}) & $\leq 60\%$\\
	& Model-Driven + atmospheric correction & $\leq 25\%$ \\
	& Model-Driven + atmospheric correction + Data-Driven (\acs{icc}) & $\leq 10\%$ \\
	
	\SetCell[r=3]{l} Azimuth Accuracy (\%) 
	& Model-Driven (only orbit and \acs{dem}) & $\leq 150\%$ \\
    & Model-Driven + atmospheric correction & $\leq 50\%$ \\
	& Data-Driven (\acs{icc}) & $\leq 10\%$ \\
	
	\acs{rcs} difference (dB) & Stability of the \acs{rcs} comparing different acquisition geometries over the same area considering a resolution of 30m & $\leq 5$~dB \\
	
\end{longtblr}

\subsubsection{Timeliness}
For many potential use cases of \acs{orb} products, timeliness can be critical, as the information in the image is highly ephemeral. The production of the product has the same timeline of \acs{nrb} product. In case of \acs{orb} being the requirement on geolocation more relaxed, it is possible to work only with nominal geometry without need for refinement of timing. Atmospheric corrections to produce unbiased NRCS estimates and to provide ancillary surface wind estimates are needed to allow the use of the \acs{orb} product for atmospheric applications (wind estimation, Marine Atmospheric Boundary Layer (MABL) parameterization, etc). \autoref{fig:orb_single} shows the flowchart of generating the \acs{orb} product using a single \acs{slc} image.
\begin{figure}[H]
	\centering
	\resizebox{\textwidth}{!}{
		\begin{tikzpicture}[node distance=0.5cm and 2cm]
			
			% Define the style for the nodes
			% \tikzstyle{node} = [circle, draw, fill=blue!20, text centered, minimum size=1.5cm]
			
			% Draw nodes
			\node[rectin] (init) {\acs{slc}};
			\node[block, right=of init] (g2r) {Geo-to-Radar};
			\node[block, right=of g2r] (rtc) {\acs{rtc}};
			\node[block, right=of rtc] (geocoding) {Geocoding};
                \node[block, right=of geocoding] (denoise) {Denoise};
                \node[block, below=of denoise] (outlier) {Outlier removal};
                \node[block, below=of outlier] (downsample) {Downsample};
                \node[block, below=of downsample] (roughness) {Normalization};
                \node[rectout, below=of roughness] (orb) {\acs{orb}};
			\node[rectaux, below=of init] (dem) {\acs{dem}};
                \node[empty, below=of g2r] (empty0) {};
			\node[rectaux, below=of empty0] (orbit) {Orbital data};
                \node[rectaux, left=of downsample] (era5) {ERA5 wind};
			\node[rectauxd, below=of orbit] (models) {External models};
                \node[rectauxd, right=of downsample] (cmod5n) {CMOD5N};
			% \node[rectind, below=of models] (ref_stack) {Reference SLC stack};
			
			% Draw edges with arrows
			\draw[->, thick] (init) -- (g2r);
			\draw[->, thick] (g2r) -- (rtc);
			\draw[->, thick] (rtc) -- (geocoding);
			\draw[->, thick] (dem) -| (g2r);
			\draw[->, thick] (orbit) -- (g2r);
			\draw[dotted, thick] (models) -- (orbit);
			% \draw[->, thick] (ref_stack) -| (g2r);
			\draw[->, thick] (geocoding) -- (denoise);
                \draw[->, thick] (denoise) -- (outlier);
                \draw[->, thick] (outlier) -- (downsample);
                \draw[->, thick] (downsample) -- (roughness);
                \draw[->, thick] (roughness) -- (orb);
                \draw[->, thick] (era5) -- (downsample);
                \draw[->, thick] (cmod5n) -- (downsample);
                % Timeline with linked markers and adaptive length
			\def\timelineStart{init}
			\def\timelineEnd{orb}
			\draw[thick, ->] (\timelineStart.west |- 0,-8) -- (\timelineEnd.east |- 0,-8) node[right] {Time};
			\foreach \node/\label in {dem/Start, models/\textbf{1-30d}, orb/<10min} {
				\draw (\node.south |- 0,-8) -- (\node.south |- 0,-8.5) node[below] {\label};
				\draw[dotted, thick] (\node.south) -- (\node.south |- 0, -8.25); % Connects the node to the timeline
			}
			
		\end{tikzpicture}
	}
	\caption{Sample workflow for generation of \acs{orb} considering each \acs{slc} image separately.}
	\label{fig:orb_single}
\end{figure}


\subsection{Usability, Integrability and Interoperability}
\label{sec:orb_usability_integrability}
The characteristics of the \acs{orb} products in terms of usability, integrability and interoperability  are in the following table.

\begin{longtblr}[
  caption = {Usability, Integrability, and Interoperability of \acs{orb} product},
  label = {tab:orb_usability_integrability_interoperability},
]{
  colspec = {Q[2cm] Q[12cm]},
  rowhead = 1,
  width = \linewidth
}
\hline
\textbf{Aspect} & \textbf{Description} \\
\hline
Usability &
\begin{itemize}
    \item Pre-Calibrated Data: the \acs{orb} product is radiometrically calibrated, meaning users do not need to perform additional steps to correct the radar backscatter values. This simplifies the analysis for the users, including those who may not be experts in radar data processing.
    \item Normalized Backscatter: since the backscatter values are normalized (converted into \gls{sigma0} by applying a geoid correction procedure), they can be easily compared across different acquisitions, sensors, or times. 
    \item Geometrically Corrected (Georeferenced): \acs{orb} products are orthorectified and projected into standard map coordinates, reducing the need for users to handle raw slant-range data or worry about terrain-induced distortions (unless further specific geometric corrections are needed for the user's application).
    \item Scalable and Flexible Data Formats: \acs{orb} products are typically provided in common geospatial formats (such as GeoTIFF, NetCDF, etc.), which are compatible with a wide range of geospatial tools (e.g., QGIS) and programming libraries (e.g., GDAL, Python). This facilitates easy integration into existing analysis workflows.
    \item Ready for Machine Learning: \acs{orb} products, being standardized and normalized, can be directly ingested into machine learning models for classification, regression, or segmentation tasks without the need for extensive preprocessing. Their calibrated and consistent nature makes them suitable for training AI models on large datasets.
\end{itemize}
\\
Integrability &
\begin{itemize}
    \item Integration into \acs{gis} platforms: \acs{orb} products are easy to integrate into Geographic Information Systems (\acs{gis}) due to their georeferencing and common data formats. Users can overlay these products with vector layers, other remote sensing data (optical, infrared), or ground truth data for comprehensive spatial analysis.
    \item Cloud-based processing: \acs{orb} products can be easily integrated into cloud-based processing environments like Google Earth Engine (\acs{gee}), \acs{aws}, or Sentinel Hub, etc. This allows large-scale, on-demand processing without requiring local resources. Users can perform global or regional analysis by leveraging cloud computing for processing extensive datasets.
    \item Compatibility with remote sensing software: software packages like ESA SNAP (Sentinel Application Platform), ENVI, or commercial and free GIS software like ArcGIS and QGIS can directly process and analyse \acs{orb} products, making them easy to integrate into radar data analysis pipelines.
\end{itemize}
\\
Interoperability &
\begin{itemize}
  \item Cross-mission compatibility for ROSE-L and Sentinel-1 missions
  \item Variety of applications, particularly for monitoring and analysing oceanic and coastal environments: Sea Surface Wind Monitoring, Wave Height and Ocean State Estimation, Ocean Current Mapping, Oil Spill Detection, Sea Ice and Iceberg Monitoring, Coastal Erosion and Sediment Transport Analysis, Ship Detection and Maritime Surveillance.
  \item High interoperability due to standard geographic referencing and pixel alignment across scenes
\end{itemize}
\\
\hline
\end{longtblr}

% \begin{itemize}
%     \item Usability.
%         \begin{itemize}
%             \item Pre-Calibrated Data: the \acs{orb} product is radiometrically calibrated, meaning users do not need to perform additional steps to correct the radar backscatter values. This simplifies the analysis for the users, including those who may not be experts in radar data processing.
%             \item Normalized Backscatter: since the backscatter values are normalized (converted into sigma-nought by applying a geoid correction procedure), they can be easily compared across different acquisitions, sensors, or times. %This makes it highly useful for time series analysis, change detection, and long-term monitoring without worrying about differences in acquisition geometry or sensor calibration. This increases the product’s usability for land cover classification, land use analysis, and surface change monitoring.
%             \item Geometrically Corrected (Georeferenced): \acs{orb} products are orthorectified and projected into standard map coordinates, reducing the need for users to handle raw slant-range data or worry about terrain-induced distortions (unless further specific geometric corrections are needed for the user's application).
%             \item Scalable and Flexible Data Formats: \acs{orb} products are typically provided in common geospatial formats (such as GeoTIFF, NetCDF, etc.), which are compatible with a wide range of geospatial tools (e.g., QGIS) and programming libraries (e.g., GDAL, Python). This facilitates easy integration into existing analysis workflows.
%             \item Ready for Machine Learning: \acs{orb} products, being standardized and normalized, can be directly ingested into machine learning models for classification, regression, or segmentation tasks without the need for extensive preprocessing. Their calibrated and consistent nature makes them suitable for training AI models on large datasets.
%         \end{itemize}
%     \item Integrability into different systems and workflows.
%         \begin{itemize}
%             \item Integration into \acs{gis} platforms: \acs{orb} products are easy to integrate into Geographic Information Systems (\acs{gis}) due to their georeferencing and common data formats. Users can overlay these products with vector layers, other remote sensing data (optical, infrared), or ground truth data for comprehensive spatial analysis.
%             % \item Interoperability with other \acs{ard} products: the \acs{nrb} could be integrated with other \acs{ard} products (such as optical data, LiDAR, or \acs{sar} products) for multi-sensor analysis. For instance, combining \acs{nrb} data with optical satellite imagery can enhance land cover classification or environmental monitoring efforts. In addition, they can complement \acs{sar} interferometric products (e.g., interferograms, coherence maps) by providing context on surface properties and changes in scattering characteristics. This is valuable in areas like infrastructure monitoring, natural hazard detection, and land subsidence studies.
%             \item Cloud-based processing: \acs{orb} products can be easily integrated into cloud-based processing environments like Google Earth Engine (\acs{gee}), \acs{aws}, or Sentinel Hub, etc.. This allows large-scale, on-demand processing without requiring local resources. Users can perform global or regional analysis by leveraging cloud computing for processing extensive datasets.
%             \item Compatibility with remote sensing software: software packages like ESA SNAP (Sentinel Application Platform), ENVI, or commercial and fre GIS software like ArcGIS and QGIS can directly process and analyse \acs{orb} products, making them easy to integrate into radar data analysis pipelines.
%             % \item Environmental and Agricultural Monitoring Systems: \acs{nrb} products can be integrated into environmental monitoring platforms to assess parameters like soil moisture, forest cover, and agricultural productivity. By coupling radar data with models or indices, they offer insights into ecosystem health or crop conditions even under cloud cover or adverse weather conditions.
%         \end{itemize}
% \end{itemize}

% % The usability of NRB products is highly dependent on the spatial and temporal resolution. For high-frequency monitoring or fine-scale analysis, users must ensure that the \acs{nrb} product offers the necessary resolution. For example, low-resolution \acs{nrb} products might be unsuitable for fine-scale monitoring.

% % \acs{nrb} products often include multiple polarizations (e.g., VV, VH), which provide complementary information about surface properties. The ability to use different polarizations enhances integrability into diverse applications like vegetation analysis (VH polarization is sensitive to biomass) or water body detection (VV polarization).

% % In case of integration with other products, it could be necessary and easier to define a common data format, reference system, and geographic grid.

% % In some of the analysed \acs{nrb} products the speckle filter is not applied. For specific analyses like land cover classification or vegetation monitoring, speckle noise mitigation may be required. This is an additional preprocessing step that one may need to consider.

% % Finally, radar backscatter can be influenced by surface roughness, moisture content, and other environmental factors. These must be accounted for when interpreting the data, especially in applications where surface characteristics (e.g., soil moisture, vegetation growth) vary over time.

% The Ocean Radar Backscatter (\acs{orb}) product has a variety of applications, particularly for monitoring and analysing oceanic and coastal environments. Key applications include:
% \begin{itemize}
%     \item Sea Surface Wind Monitoring,
%     \item Wave Height and Ocean State Estimation,
%     \item Ocean Current Mapping,
%     \item Oil Spill Detection,
%     \item Sea Ice and Iceberg Monitoring,
%     \item Coastal Erosion and Sediment Transport Analysis,
%     \item Ship Detection and Maritime Surveillance.
% \end{itemize}

% % Normalized Radar Backscatter products are versatile for a wide range of applications, from environmental monitoring to infrastructure management, disaster response, and resource exploration. Below are some of the key applications of \acs{orb} products:
% % \begin{itemize}
% %     \item Land Cover Classification.
% %         \begin{itemize}
% %             \item Vegetation Monitoring: \acs{orb} products can help in identifying different types of vegetation and monitoring their health, structure, and density, as different land covers (e.g., forests, agricultural fields, grasslands) have distinct radar backscattering signatures.
% %             \item Urban vs. Rural Classification: Radar backscatter can distinguish between built-up areas and natural land covers.
% %         \end{itemize}
% %     \item Agriculture.
% %         \begin{itemize}
% %             \item Crop Monitoring and Yield Prediction: \acs{orb} can be used to monitor crop growth stages and predict yields by identifying differences in crop structures.
% %             \item Soil Moisture Estimation: SAR backscatter is sensitive to soil moisture content, making \acs{orb} products useful for assessing soil moisture in agricultural fields, which is crucial for water management and irrigation practices.
% %         \end{itemize}
% %     \item Flood Detection: Radar backscatter is affected by water, enabling \acs{orb} products to detect flooded areas
% %         \begin{itemize}
% %             \item in agricultural lands after heavy rains or during irrigation events.
% %             \item in urban zones
% %             \item for tracking water levels in real-time, even during cloudy or rainy conditions
% %             \item for creating flood risk maps
% %         \end{itemize}
% %     \item Change Detection
% %         \begin{itemize}
% %             \item for disaster management (earthquake and post-disaster damage assessment, forest fire and burned area mapping)
% %             \item for deforestation monitoring
% %             \item for urbanization monitoring
% %             \item for volcanic monitoring (e.g., appearance of new lava flows)
% %             \item for monitoring the polar ice caps
% %         \end{itemize}
% %     \item Forest Biomass Estimation \& Carbon Stock Estimation for climate change monitoring
% %     \item Water Resource Management.
% %         \begin{itemize}
% %             \item Wetland Mapping: The ability of \acs{orb} to differentiate between water and land makes it useful for mapping and monitoring wetlands and other water bodies.
% %             \item River and Lake Monitoring: \acs{orb} products can help monitor changes in river flows and lake levels, essential for hydrology studies and water resource management.
% %         \end{itemize}
% %     \item Military and Security Applications: \acs{orb} products can detect changes in terrain, infrastructure, border and coast or activity in a region, even under poor visibility, making them valuable for military and security operations.
% %     \item Environmental Monitoring.
% %         \begin{itemize}
% %             \item Coastal Erosion Monitoring: \acs{orb} products are useful in tracking shoreline changes and coastal erosion, which is important for coastal management and protection.
% %             \item Desertification and Land Degradation: to monitor areas undergoing desertification or other forms of land degradation by detecting changes in surface properties.
% %             \item Monitoring of Environmental Protection Areas: \acs{orb} products can be used to track illegal activities, such as deforestation or construction, in protected areas.
% %         \end{itemize}
% % \end{itemize}


\section{Precise Range \ Azimuth Offsets (PRO / PAO \& GPRO / GPAO )}

\subsection{Product Description}
\label{sec:pro_pao_prod_description}
The pixel offset products contain information on the shifts in range and azimuth between a reference and secondary \acs{cslc}s (pair of \acs{cslc}). Offsets can be estimated by exploiting phase and/or amplitude information. Whenever possible, techniques based on the analysis of the phase signal provide information with better accuracy. 

The phase information can be used to retrieve, both slant range and azimuth offsets, in areas where two different looks of the same scene are available. In the case of Sentinel-1 data, this operation is only feasible within the limited areas where bursts overlap, due to the characteristics of its acquisition mode. However, in TOPSAR mode, the use of phase modulation improves sensitivity, making it more effective for this kind of analysis. It is important to ensure that the interferometric phase remains continuous between adjacent bursts. Additionally, the residual phase that remains after standard processing can be used to detect and highlight local movements in the azimuth direction, providing valuable insights into dynamic surface phenomena. The range offsets, can be derived by standard differential interferometry: in this case, \acs{pro} and \acs{int} coincide, as well as \acs{gpro} and \acs{gunw}.

In case of using amplitude information, the range and azimuth local residual offsets can be estimated via cross-correlation procedure.

The selection criteria of the pairs of images, from which the products will be generated, will consider:
\begin{itemize}
    \item on demand request: the products could be required for a specific event (e.g., earthquake, fast landslides, fast glacier motion, etc.). In this case, the selected pairs of images, including images acquired in ascending and descending passes, before and after the event, and temporally closer to it.
    \item systematic production in correspondence of occurrence of a specific event (included in a predefined list of events, e.g. earthquake, fast landslides, glacier collapse, etc.). In this case, the selection criteria is the same of the on demand request.
    \item systematic production: the products could be generated for pairs of consecutive images acquired in interferometric mode over predefined geographic areas.
\end{itemize}

The fact that \acs{pro} \& \acs{pao} are derived from pairs of \acs{cslc}s, defined on the reference radar grid common to both Sentinel-1 and ROSE-L missions, implies that the output radar grid of these products will also be the same for Sentinel-1 and ROSE-L. This ensures the joint use of products generated from data acquired by the two missions.

\subsubsection{Input and Auxiliary Data}
\label{sec:pro_pao_aux_data}
The following input and auxiliary data are considered for product generation:

\begin{longtblr}[
	label = tab:in_aux_pro_pao_products,
	entry = none,
	caption = {"Precise Range \& Azimuth Offsets" input and auxiliary data},
	]{
		width = \linewidth,
		colspec = {Q[l,20]Q[l,60]Q[c,20]},
		hlines,
		hline{1-Z} = {-}{0.05em},
	}
	\textbf{Name}   & \textbf{Description}   & \textbf{Mandatory (M) / Optional (O)} \\
	\acs{cslc}s & Pair of co-registered L1 Single Look Complex image (\acs{cslc}) and annotations. The \acs{cslc} must be corrected by using \acs{etad}. & M \\
    \acs{dem} & External DEM (e.g., Copernicus DEM at different resolutions, with 30 m as the default worldwide). & M only for \acs{gpao} and \acs{gpro} \\
    Precise Orbits & \acs{cslc} precise orbits & M only for \acs{gpao} and \acs{gpro} \\
\end{longtblr}

\acs{etad} products are not considered as inputs for \acs{pro}/\acs{pao} \& \acs{gpro}/\acs{gpao}. All corrections using \acs{etad} are applied at the \acs{cslc} level.

\subsubsection{Main Variables}
The following tables list the main variables associated with the "Precise Range \& Azimuth Offsets" products.

\begin{longtblr}[
	label = tab:var_px_off_products,
	entry = none,
	caption = {"Precise Range \& Azimuth Offsets" main variables},
	]{
		width = \linewidth,
		colspec = {Q[l,15]Q[l,32]Q[l,40]Q[l,13]},
		hlines,
		hline{1-Z} = {-}{0.05em},
	}
	\textbf{Name} & \textbf{Description} & \textbf{Dimensions} & \textbf{Data Type} \\
    ${\Delta_{rg}\: (m)}$ & Offsets measured along slant range direction (in meters) & polarization, acquisition\_date, azimuth\_time or y (for radar or projected geometry, respectively), slat\_range\_time or x (for radar or projected geometry, respectively) & float32 \\
    ${\Delta_{az}\: (m)}$ & Offsets measured along azimuth direction (in meters) & polarization, acquisition\_date, azimuth\_time or y (for radar or projected geometry, respectively), slat\_range\_time or x (for radar or projected geometry, respectively) & float32 \\
\end{longtblr}

The table below describes the dimensions/coordinates associated with the "Precise Range \& Azimuth Offsets" main variables.

\begin{longtblr}[
	label = tab:var_px_off_coordinates,
	entry = none,
	caption = {"Precise Range \& Azimuth Offsets" main variables coordinates/dimensions},
	]{
		width = \linewidth,
		colspec = {Q[l,20]Q[l,62]Q[l,18]},
		hlines,
		hline{1-Z} = {-}{0.05em},
	}
	\textbf{Name} & \textbf{Description} & \textbf{Data Type} \\
	polarization & Acquired polarizations & string \\
    acquisition\_date & Image acquisition dates, formatted as YYYY-MM-DD & string \\
    azimuth\_time & Time along satellite track direction & datetime64 \\
    range\_time & Time across satellite track direction & float64 \\
    x & UTM/UPS x (East) coordinates [m] & float64 \\
    y & UTM/UPS y (North) coordinates [m] & float64 \\
\end{longtblr}

\subsubsection{Main Metadata}
\label{sec:metadata_pix_off}
The following table lists the main metadata and per-pixel metadata associated with the "Precise Range \ Azimuth Offsets" product.

\begin{longtblr}[
	label = tab:meta_px_off_products,
	entry = none,
	caption = {"Precise Range \ Azimuth Offsets" main metadata and per-pixel metadata},
	]{
		width = \linewidth,
		colspec = {Q[l,17]Q[l,38]Q[l,20]Q[l,13]Q[c,12]},
		hlines,
		hline{1-Z} = {-}{0.05em},
	}
	\textbf{Name} & \textbf{Description} & \textbf{Dimensions} & \textbf{Data Type} & \textbf{Per-pixel Metadata}\\
    geometric
    
    \_accuracy
    
    \_type	& Type of measured geometric accuracy (only for projected products)	& - & string & \\
    \_geometric
    
    \_accuracy
    & Measured geometric accuracy, including eastern and northern geometric accuracy bias and standard deviation (only for projected products). The measures will be performed on a validation dataset and periodically updated. & - & dictionary & \\ 
    east\_geom
    
    \_bias & Measured eastern geometric bias & - & float32 & \\
    east\_geom\_std & Measured eastern geometric standard deviation (only for projected products) & - & float32 & \\
    north\_geom
    
    \_bias & Measured northern geometric bias & - & float32 & \\
    north\_geom
    
    \_std & Measured northern geometric standard deviation (only for projected products) & - & float32 & \\
    geometric
    
    \_accuracy
    
    \_radial
    
    \_rmse	& Measured radial RMSE of geometric accuracy (only for projected products) & - & float32 & \\
     mask & Mask containing, for each ground cell for projected product or pixel for product in radar geometry, respectively, the following information:
    
     - 0: Invalid data
    
     - 1: Valid data

     - 2: No data
    
     & x and y & byte & \ding{51} \\
     correlation
     
     \_peak & Peak of the normalized correlation coefficient. Present only if the products are derived by using the images' amplitude & x and y for projected products,
     
     range\_time and azimuth\_time for products in radar geometry & float32 & \ding{51} \\
     coherence & coherence map. Present only if the products are derived by using the images' phase & x and y for projected products,
     
     range\_time and azimuth\_time for products in radar geometry & float32 & \ding{51} \\
     x\_sample & Pixel spacing in x direction & - & float32 & \\
     y\_sample & Pixel spacing in y direction & - & float32 & \\
\end{longtblr}

The table below describes the dimensions/coordinates associated with the "Precise Range \ Azimuth Offsets" main metadata.

\begin{longtblr}[
	label = tab:var_pix_off_metadata_coordinates,
	entry = none,
	caption = {"Precise Range \ Azimuth Offsets" main metadata coordinates/dimensions},
	]{
		width = \linewidth,
		colspec = {Q[l,20]Q[l,62]Q[l,18]},
		hlines,
		hline{1-Z} = {-}{0.05em},
	}
	\textbf{Name} & \textbf{Description} & \textbf{Data Type} \\
	azimuth\_time & Time along satellite track direction & datetime64 \\
    range\_time & Time across satellite track direction & float64 \\
    x & UTM/UPS x (East) coordinates [m] & float64 \\
    y & UTM/UPS y (North) coordinates [m] & float64 \\
\end{longtblr}

\subsubsection{Reference System Tiling and Slicing}
\label{sec:propao_reference_system_tiling}
\autoref{tab:px_off_ref_system_slicing} reports the reference system, the slicing, and grid for "Precise Range \ Azimuth Offsets" products.

\begin{longtblr}[
	label=tab:px_off_ref_system_slicing,
	entry = none,
	caption = {Reference system and Slicing for "Precise Range \ Azimuth Offsets"},
	]{
		width=\linewidth,
		colspec = {Q[l,15]Q[l,40]Q[l,15]Q[l,30]},
		%vline{2-5} = {-}{},
		%vline{5-5} = {-}{},
		hlines,
		hline{1-Z} = {-}{0.05em},
		rowhead=2,
	}
	\textbf{Product} & \textbf{Reference System} & \textbf{Slicing} & \textbf{Spacing}\\
	\textbf{\acs{pro} / \acs{pao}} & Radar: azimuth/slant-range & Burst & ~50 m \\
	\textbf{\acs{gpro} / \acs{gpao}} & Cartographic: UTM/UPS & Burst-projected & 50 m \\   
	
\end{longtblr}

It is planned to use the same reference radar grid and the same slicing for the ROSE-L and Sentinel-1 missions, to ensure interoperability between the products generated with the data of both missions.

\subsection{Product Requirements}

\subsubsection{Accuracy}
\label{pro_pao_accuracy}

The accuracy of the \acs{pro} / \acs{pao} products is critically influenced by the chosen estimation technique, the sensor characteristics, and the underlying scene properties. For the purpose of this analysis, the remaining sources of error stem primarily from thermal noise, spatial decorrelation, local terrain-induced geometric distortions, and processing algorithm limitations, regardless of the presence of some residual ionospheric effects after the corrections carried out previously over the SLC/CSLC products.

Two main approaches are commonly employed to estimate the offsets:
\begin{itemize}
    \item Interferometric phase-based methods. When the SAR scene preserves high coherence between acquisitions, interferometric techniques can be used to derive precise displacement measurements. In particular, the phase difference provides a direct measurement of displacement along the radar line-of-sight. This phase, once unwrapped and converted through known radar geometry, yields an accurate estimate of the range displacement. However, azimuthal displacement cannot be directly measured from a single interferogram. In specific configurations, such as burst overlap regions in Sentinel-1's Terrain Observation with Progressive Scans (TOPS) mode, a two-look geometry can be exploited. By generating interferograms from overlapping subswaths or looks and forming an "interferogram of interferograms", it becomes possible to retrieve the azimuth component of the displacement, albeit with lower sensitivity and increased noise.
    \item Amplitude-based pixel offset techniques. In low-coherence or geometrically complex scenes, pixel offset techniques based on amplitude cross-correlation are often preferred. These methods estimate the local displacement by identifying the maximum correlation between image patches across the two SAR acquisitions. Pixel offset is less sensitive to decorrelation but typically achieves lower precision compared to interferometric methods. Moreover, the final resolution and accuracy are influenced by the window size used for matching, the oversampling strategy, the image signal-to-noise and signal-to-clutter ratio.
\end{itemize}
Overall, the accuracy of PRO/PAO products varies depending on the selected method, the radar mode, and scene content. Interferometric methods generally offer higher precision (especially in range), but require high coherence and are limited in their ability to extract azimuth displacement unless specific multi-look geometries are available. Pixel offset is complementary to SAR interferometry and is applicable in case of larger displacements, where the interferometric phase become incoherent, at the expense of somewhat lower resolution and accuracy.

The following table summarizes the accuracy requirements for \acs{pro} and \acs{pao} products, considering the interferometric phase-based methods. These requirements are computed separately for ROSE-L and Sentinel-1 satellites, since they depend on satellite properties like wavelength and spectral diversity, and are intended as evaluated on pixels with coherence greater than 0.7, and performing an average over cell of 200m x 200m on the measurements in azimuth direction.
\begin{longtblr}[
  label = tab:pro_pao_accuracy_phase_s1,
  entry = none,
  caption = {Accuracy Requirements for \acs{pro}/\acs{pao} Products using Phase-based methods in case of \acs{s1} data},
]{
  width = \linewidth,
  colspec = {Q[l,35] Q[l,55] Q[c,20]},
  hlines,
  hline{1-Z} = {-}{0.05em},
  rowhead = 1,
}

\textbf{Metric} & \textbf{Description} & \textbf{Requirement} \\

$\sigma_{\Delta^{ph}_{rg}}$ & Range offset accuracy for \acs{s1} mission & $\leq 0.005$ m \\

$\sigma_{\Delta^{ph}_{az}}$ & Azimuth offset accuracy for \acs{s1} mission & $\leq 0.04$ m \\

\end{longtblr}


\begin{longtblr}[
  label = tab:pro_pao_accuracy_phase_rosel,
  entry = none,
  caption = {Accuracy Requirements for \acs{pro}/\acs{pao} Products using Phase-based methods in case of \acs{rosel} data},
]{
  width = \linewidth,
  colspec = {Q[l,35] Q[l,55] Q[c,20]},
  hlines,
  hline{1-Z} = {-}{0.05em},
  rowhead = 1,
}

\textbf{Metric} & \textbf{Description} & \textbf{Requirement} \\

$\sigma_{\Delta^{ph}_{rg}}$ & Range offset accuracy for \acs{rosel} mission & $\leq 0.02$ m \\

$\sigma_{\Delta^{ph}_{az}}$ & Azimuth offset accuracy for \acs{rosel} mission & $\leq 1.2$ m \\

\end{longtblr}


The following table summarizes the accuracy requirements for \acs{pro} and \acs{pao} products, considering the amplitude-based pixel offset technique. These requirements are intended to be evaluated using patches with of 200m x 200m, and a coherence value greater than 0.5.

\begin{longtblr}[
  label = tab:pro_pao_accuracy_pixel_offset,
  entry = none,
  caption = {Accuracy Requirements for \acs{pro}/\acs{pao} Products using Pixel Offset in case of \acs{s1} and \acs{rosel} data},
]{
  width = \linewidth,
  colspec = {Q[l,35] Q[l,55] Q[c,20]},
  hlines,
  hline{1-Z} = {-}{0.05em},
  rowhead = 1,
}

\textbf{Metric} & \textbf{Description} & \textbf{Requirement} \\

$\sigma_{\Delta^{amp}_{rg}}$ & Range offset accuracy assuming $\sim$ 5 m range resolution & $\leq 0.5$ m \\

$\sigma_{\Delta^{amp}_{az}}$ & Azimuth offset accuracy assuming $\sim$ 20.0 m azimuth resolution & $\leq 2$ m \\

\end{longtblr}

\subsubsection{Timeliness}
\acs{pro} and \acs{pao} are generated from interferometric products (phase-based approach) or by-products of the data-driven co-registration (pixel offset approach). Displacement or varying atmospheric conditions between acquisitions have the same effect on the SAR image: a misplacement of the object. The data-driven co-registration process can estimate the offsets between images but cannot separate atmospheric effects from the actual displacement of objects on the ground. Therefore, it is necessary to filter out artifacts, particularly ionospheric effects. Model-based or data-driven approaches can be followed for \acs{cslc} and interferometric products generation. 

Depending on the target application, the model-based approach can be adopted for non-time-critical use cases, while the data-driven approach can be selected for products requiring faster delivery times.

In the following, the different possible workflows are reported, depending on the \acs{pro} and \acs{pao} techniques, as well as the exploitation of external models or data-driven approaches.
In addition, the timeline provided indicates the time required to collect the input data and complete the main processing steps. The processing time refers to a standard \acs{s1} dataset and is approximated, considering that the latency in obtaining the input data is the dominant factor.

The first workflow refers to a phase-based technique, including interferometric processing, and the use of external models for compensating atmospheric and processing artifacts. The timeliness of this workflow is strongly affected by the availability of external data.

\begin{figure}[H]
	\resizebox{\textwidth}{!}{
		\begin{tikzpicture}[node distance=1cm and 2cm]
			
			% Define the style for the nodes
			% \tikzstyle{node} = [circle, draw, fill=blue!20, text centered, minimum size=1.5cm]
			
			% Draw nodes
			\node[rectin] (init) {\acs{slc}};
			\node[block, right=of init] (coreg_geo) {Geometric co-registration};
			\node[block, right=of coreg_geo] (coreg_ref) {Refinement};
			\node[rectaux, below=of init] (dem) {\acs{dem}};
			\node[empty, below=of coreg_geo] (empty1) {};
			\node[rectaux, below=of empty1] (orbit) {Orbital data};
			\node[empty, below=of coreg_ref] (empty2) {};            
			\node[rectaux, below=of empty2] (models) {External models};
			\node[rectout, right=of coreg_ref] (cslc) {\acs{cslc}};
			\node[block, right=of cslc] (interf) {Interf};
			\node[rectout, right=of interf] (phase) {PHASE};
			\node[block, right=of phase] (propaoproc) {PHDIFF};
			\node[rectout, right=of propaoproc] (pro_pao) {\acs{pro} and \acs{pao}};
			
			%\node[block, right=of coreg_ref] (interf) {PHASE};
			%\node[rectout, right=of interf] (pro_pao) {\acs{pro} & \acs{pao}};
			
			% Draw edges with arrows
			\draw[->, thick] (init) -- (coreg_geo);
			\draw[->, thick] (coreg_geo) -- (coreg_ref);
			\draw[->, thick] (dem) -| (coreg_geo);
			\draw[->, thick] (orbit) -- (coreg_geo);
			\draw[->, thick] (models) -- (coreg_ref); % Changed -| to --
			\draw[->, thick] (coreg_ref) -- (cslc);
			\draw[->, thick] (cslc) -- (interf);
			\draw[->, thick] (interf) -- (phase);
			\draw[->, thick] (phase) -- (propaoproc);
			\draw[->, thick] (propaoproc) -- (pro_pao);
			
			% Timeline with linked markers and adaptive length
			\def\timelineStart{init}
			\def\timelineEnd{pro_pao}
			%\pgfnodealias{start}{\timelineStart}
			%\pgfnodealias{end}{\timelineEnd}
			\draw[thick, ->] (\timelineStart.west |- 0,-6) -- (\timelineEnd.east |- 0,-6) node[right] {Time};
			
			\foreach \node/\label in {init/Start, orbit/1h-3d, models/\textbf{1-30d}, cslc/<1h, phase/<1h, pro_pao/<1h} {
				\draw (\node.south |- 0,-6) -- (\node.south |- 0,-6.5) node[below] {\label};
				\draw[dotted, thick] (\node.south) -- (\node.south |- 0, -6.25); % Connects the node to the timeline
			}
		\end{tikzpicture}
	}	
	\caption{Sample workflow for generation of \acs{pro} and \acs{pao} products exploiting interferometric processing and co-registration refinement with external models.}
	\label{fig:pro_pao_phase_model_driven}
\end{figure}

Alternatively, a data-driven approach can be followed, estimating corrections directly from a stack of SAR data. In this case, the processing is more complex, generally involving more than a single image and requiring the solution of inverse problems with high computational complexity.
\begin{figure}[H] % Changed H to h!
	\resizebox{\textwidth}{!}{
		\begin{tikzpicture}[node distance=1cm and 2cm]			
			% Draw nodes
			\node[rectin] (init) {\acs{slc}};
			\node[block, right=of init] (coreg_geo) {Geometric co-registration};
			\node[block, right=of coreg_geo] (coreg_data_driven) {Data-driven co-registration};
			\node[rectaux, below=of init] (dem) {\acs{dem}};
			\node[empty, below=of coreg_geo] (empty1) {};
			\node[rectaux, below=of empty1] (orbit) {Orbital data};
                \node[empty, below=of coreg_data_driven] (empty2) {};  
			\node[rectin, below=of empty2] (ref_stack) {Reference SLC stack};
			\node[rectout, right=of coreg_data_driven] (cslc) {\acs{cslc}};
			\node[block, right=of cslc] (interf) {Interf};
			\node[rectout, right=of interf] (phase) {PHASE};
			\node[block, right=of phase] (propaoproc) {PHDIFF};
			\node[rectout, right=of propaoproc] (pro_pao) {\acs{pro} and \acs{pao}};
			% Draw edges with arrows
			\draw[->, thick] (init) -- (coreg_geo);
			\draw[->, thick] (coreg_geo) -- (coreg_data_driven);
			\draw[->, thick] (dem) -| (coreg_geo);
			\draw[->, thick] (orbit) -- (coreg_geo);
			\draw[->, thick] (ref_stack) -- (coreg_data_driven);
			\draw[->, thick] (coreg_data_driven) -- (cslc);
			\draw[->, thick] (cslc) -- (interf);
			\draw[->, thick] (interf) -- (phase);
			\draw[->, thick] (phase) -- (propaoproc);
			\draw[->, thick] (propaoproc) -- (pro_pao);			
			% Timeline with linked markers and adaptive length
			\def\timelineStart{init}
			\def\timelineEnd{pro_pao}
			\draw[thick, ->] (\timelineStart.west |- 0,-6) -- (\timelineEnd.east |- 0,-6) node[right] {Time};
			\foreach \node/\label in {init/Start, orbit/1h-3d, ref_stack/<1h, cslc/<1h, phase/<1h, pro_pao/<1h} {
				\draw (\node.south |- 0,-6) -- (\node.south |- 0,-6.5) node[below] {\label};
				\draw[dotted, thick] (\node.south) -- (\node.south |- 0, -6.25); % Connects the node to the timeline
			}
		\end{tikzpicture}	
	}
	\caption{Sample workflow for generation of \acs{pro} and \acs{pao} products exploiting interferometric processing and co-registration data-driven.}
	\label{fig:pro_pao_phase_data_driven}
\end{figure}

In case of pixel-offset technique a dense cross-correlation is applied to the \acs{cslc} product in order to extract the residual local offset. 

\begin{figure}[H] % Changed H to h!
	\resizebox{\textwidth}{!}{
		\begin{tikzpicture}[node distance=1cm and 2cm]			
			% Draw nodes
			\node[rectin] (init) {\acs{slc}};
			\node[block, right=of init] (coreg_geo) {Geometric co-registration};
			\node[block, right=of coreg_geo] (coreg_data_driven) {Data-driven co-registration};
			\node[rectaux, below=of init] (dem) {\acs{dem}};
                \node[empty, below=of coreg_geo] (empty1) {};
			\node[rectaux, below=of empty1] (orbit) {Orbital data};
                \node[empty, below=of coreg_data_driven] (empty2) {};
			\node[rectin, below=of empty2] (ref_stack) {Reference SLC stack};
			\node[rectout, right=of coreg_data_driven] (cslc) {\acs{cslc}};
			\node[block, right=of cslc] (crosscorr) {Dense correlation};
			\node[rectout, right=of crosscorr] (pro_pao) {\acs{pro} and \acs{pao}};
			% Draw edges with arrows
			\draw[->, thick] (init) -- (coreg_geo);
			\draw[->, thick] (coreg_geo) -- (coreg_data_driven);
			\draw[->, thick] (dem) -| (coreg_geo);
			\draw[->, thick] (orbit) -- (coreg_geo);
			\draw[->, thick] (ref_stack) -- (coreg_data_driven);
			\draw[->, thick] (coreg_data_driven) -- (cslc);
			\draw[->, thick] (cslc) -- (crosscorr);
			\draw[->, thick] (crosscorr) -- (pro_pao);			
			% Timeline with linked markers and adaptive length
			\def\timelineStart{init}
			\def\timelineEnd{pro_pao}
			\draw[thick, ->] (\timelineStart.west |- 0,-6) -- (\timelineEnd.east |- 0,-6) node[right] {Time};
			\foreach \node/\label in {init/Start, orbit/1h-3d, ref_stack/<1h, cslc/<1h, pro_pao/<1h} {
				\draw (\node.south |- 0,-6) -- (\node.south |- 0,-6.5) node[below] {\label};
				\draw[dotted, thick] (\node.south) -- (\node.south |- 0, -6.25); % Connects the node to the timeline
			}
		\end{tikzpicture}	
	}
	\caption{Sample workflow for generation of \acs{pro} and \acs{pao} products exploiting pixel offset processing and co-registration data-driven.}
	\label{fig:pro_pao_amp_data_driven}
\end{figure}

\subsection{Usability, Integrability and Interoperability}
\label{sec:usability_integrability_pix_off}
The characteristics of the "Precise Range \& Azimuth Offsets" products in terms of usability, integrability and interoperability are in the following table.

\begin{longtblr}[
  caption = {Usability, Integrability, and Interoperability of "Pixel Offsets" Products},
  label = {tab:pixel_offsets_characteristics},
]{
  colspec = {Q[2cm] Q[12cm]},
  rowhead = 1,
  width = \linewidth
}
\hline
\textbf{Aspect} & \textbf{Description} \\
\hline
Usability &
\begin{itemize}
  \item Displacement Monitoring: Ideal for detecting large displacements such as landslides, glacier flow, earthquakes, or mining subsidence, in both range and azimuth directions.
  \item Complementary to InSAR: Effective where InSAR fails (e.g., due to decorrelation or large deformation). Allows capturing horizontal (east/north) and vertical motion in case of using double geometry.
  \item Global to Local Monitoring: Suitable for regional or global-scale observations (e.g., glaciers, tectonics) and local events (e.g., fast landslides).
  \item Broad Application Domain: Useful in natural hazard analysis, glaciology, volcanology, earthquake monitoring, and fine co-registration tasks
\end{itemize}
\\
Integrability &
\begin{itemize}
  \item Integration with InSAR: Combines with interferometric measurements to derive full 3D deformation fields when ascending/descending geometries are used. This integration allows a more comprehensive knowledge of the analyzed phenomena.
  \item Multi-Sensor Compatibility: Can be derived from various sensors and integrated in joint analyses.
  \item Input for Geophysical Models: Supports models of surface dynamics (e.g., glacier flow, landslides, earthquake rupture), improving simulation accuracy.
\end{itemize}
\\
Interoperability &
\begin{itemize}
  \item Sensor-Independent Formats: Pixel offsets in projected geometry can be standardized across sensors for interoperability in displacement workflows.
  \item Geolocation Harmonization: When geocoded, products align spatially with other \acs{eo} datasets (e.g., NRB, optical, DEMs) to support multi-source analysis.
  \item Software Ecosystem Support: Usable in major \acs{eo} software platforms, supporting scientific and operational pipelines.
\end{itemize}
\\
\hline
\end{longtblr}


\section{Interferometric products (INT / COHE / UNW \& GINT / GCOHE / GUNW)}

\subsection{Product Description}
The Interferogram (\acs{int}) product represents the wrapped phase difference (in radians) between two co-registered SAR acquisitions and is used to analyze surface deformation and coherence. It is generated from co-registered \acs{slc}s or \acs{cslc}s, with corrections applied for co-registration errors (e.g., Enhanced Spectral Diversity) and, when necessary, ionospheric effects (e.g., via split-spectrum technique).

The \acs{int} product includes:
\begin{itemize}
    \item Wrapped phase information
    \item Normalized interferometric coherence (\acs{cohe})
    \item Auxiliary layers such as DEM, baseline lookup tables, and optional phase correction maps
\end{itemize}

The product is available in both:
\begin{itemize}
    \item \textit{Radar geometry} (burst-wise or swath-merged)
    \item \textit{Ground-projected format} (\acs{gint} and \acs{gcohe}) in projected coordinates
\end{itemize}

The resolution may vary depending on the application, from full-resolution (e.g., urban ground deformation) to coarser, multilooked resolutions for regional-scale analyses.

The Unwrapped Interferogram (\acs{unw}) product contains the unwrapped phase (in radians), where the 2$\pi$ ambiguities present in the wrapped interferogram have been resolved. This product enables direct estimation of surface displacement, provided that the phase unwrapping process is successful and robust.

The \acs{unw} product includes:
\begin{itemize}
    \item Unwrapped phase
    \item Coherence map (\acs{cohe})
    \item Geometry masks and connected components
    \item Auxiliary layers, baseline lookup tables, and optional phase correction maps
\end{itemize}

\acs{unw} can be delivered in:
\begin{itemize}
    \item \textit{Radar geometry} (as \acs{unw})
    \item \textit{Ground-projected format} (\acs{gunw} and \acs{gcohe}) in projected coordinates
\end{itemize}

Because phase unwrapping is a known ill-posed problem and no fully reliable algorithm exists that universally guarantees high-quality results, the dissemination of \acs{unw} products must be carefully assessed. Their generation is typically limited to areas where sufficient coherence and signal stability allow for accurate unwrapping.

The interferometric products are derived from pairs of SAR acquisitions. These products are planned to be generated specifically from pairs of consecutive or time-close acquisitions.

\acs{int}, \acs{cohe}, and \acs{unw} will be generated using \acs{cslc}s and are defined on a common reference radar grid shared by both the Sentinel-1 and ROSE-L missions. The output radar grid of these products will likewise be aligned across both missions, ensuring interoperability and joint usability of products generated from data acquired by either satellite.

\subsubsection{Input and Auxiliary Data}
The following input and auxiliary data are considered for product generation:

\begin{longtblr}[
    label = tab:in_aux_int_products,
    entry = none,
    caption = {Interferometric products' input and auxiliary data},
    ]{
    width = \linewidth,
    colspec = {Q[l,20]Q[l,60]Q[c,20]},
    hlines,
    hline{1-Z} = {-}{0.05em},
    }
    \textbf{Name}   & \textbf{Description}   & \textbf{Mandatory (M) /       Optional (O)} \\
    Reference \acs{cslc} & Reference L1 Co-registered Single Look
    Complex Image and annotations & M \\
    Secondary \acs{cslc} & Secondary L1 Co-registered Single Look
    Complex Image and annotations & M \\
    \acs{dem} & External DEM (e.g., Copernicus DEM at different resolutions, with 30 m as the default worldwide). & M \\
    Precise Orbits & \acs{cslc} precise orbits. & M \\
    \acs{etad} & \acs{etad}-like products for atmospheric delay and timing corrections. & O \\
    \acs{tec} & External \acs{tec} data for delay correction & O \\
\end{longtblr}

\subsubsection{Main Variables}
The following tables list the main variables associated with the interferometric product.

\begin{longtblr}[
    label = tab:var_int_products,
    entry = none,
    caption = {Interferometric product main variables},
    ]{
        width = \linewidth,
        colspec = {Q[l,25]Q[l,52]Q[l,20]Q[l,13]},
        hlines,
        hline{1-Z} = {-}{0.05em},
    }
    \textbf{Name} & \textbf{Description} & \textbf{Dimensions} & \textbf{Data Type} \\
    interferogram & Wrapped interferometric phase between Secondary and Reference \acs{cslc}s, corrected for the topographic contribution and, optionally, for ionospheric and atmospheric effects. Present in \acs{int} and \acs{unw} and their projected version & polarization and 
    
    - range\_time, azimuth\_time for radar geometry products
    
    - x, y for projected products & float32 \\
    unwrapped 
    
    \_interferogram & Unwrapped interferometric phase between the Secondary and Reference \acs{cslc}s, corrected for the topographic contribution and, optionally, for ionospheric and atmospheric effects. Present in \acs{unw} and its projected version & polarization and 
    
    - range\_time, azimuth\_time for radar geometry products
    
    - x, y for projected products & float32 \\
    coherence & The coherence of between the reference and secondary \acs{cslc}s forming the interferometric pair & polarization and 
    
    - range\_time, azimuth\_time for radar geometry products
    
    - x, y for projected products & float32 \\
\end{longtblr}

The table below describes the dimensions/coordinates associated with the interferometric product main variables.

\begin{longtblr}[
    label = tab:var_int_coordinates,
    entry = none,
    caption = {Interferometric product main variables coordinates/dimensions},
    ]{
        width = \linewidth,
        colspec = {Q[l,20]Q[l,62]Q[l,18]},
        hlines,
        hline{1-Z} = {-}{0.05em},
    }
    \textbf{Name} & \textbf{Description} & \textbf{Data Type} \\
    polarization & Acquired polarizaions & string \\
    azimuth\_time & Time along satellite track direction & datetime64 \\
    range\_time & Time across satellite track direction & float64 \\
    x & UTM/UPS x (East) coordinates [m] & float64 \\
    y & UTM/UPS y (North) coordinates [m] & float64 \\
\end{longtblr}

\subsubsection{Main Metadata}
The following tables list the main metadata and per-pixel metadata associated with the interferometric products.

\begin{longtblr}[
    label = tab:meta_int_products,
    entry = none,
    caption = {Interferometric products' main metadata and per-pixel metadata},
    ]{
    width = \linewidth,
    colspec = {Q[l,17]Q[l,38]Q[l,20]Q[l,13]Q[c,12]},
    hlines,
    hline{1-Z} = {-}{0.05em},
    }
    \textbf{Name} & \textbf{Description} & \textbf{Dimensions} & \textbf{Data Type} & \textbf{Per-pixel Metadata}\\
    ionospheric
    
    \_phase & The ionospheric phase subtracted from the phase, if performed &  
    - range\_time, azimuth\_time for radar geometry products
    
    - x, y for projected products &  float32 & \ding{51} \\
    tropospheric
    
    \_phase & The tropospheric phase subtracted from the phase, if performed &  
    - range\_time, azimuth\_time for radar geometry products
    
    - x, y for projected products &  float32 & \ding{51} \\
    topographic
    
    \_phase & Topographic phase estimated by using DEM & - range\_time, azimuth\_time for radar geometry products
    
    - x, y for projected products & float32 & \ding{51} \\
    reference\_cslc & Link to input reference \acs{cslc} & - & string & \\
    secondary\_cslc & Link to input secondary \acs{cslc} & - & string & \\
    reference\_orbit & Link to reference orbit file & - & string & \\
    reference\_dem & Link to reference DEM & - & string & \\
    etad\_file & ETAD file used for each image belonging to the CSLC stack. This is optional: it is included only if used; otherwise, it is null & - & string & \\ 
    etad\_azcorr & Extended Timing Annotation Data: azimuth time correction resampled onto the common grid of the reference geometry. This is optional: it is included only if used; otherwise, it is null & 
    - azimuth\_time, slant\_range\_time  for product in radar geometry 
    
    - x, y for projected products & float64 & \ding{51} \\
    etad\_rgcorr & Extended Timing Annotation Data: slant range time correction resampled onto the common grid of the reference geometry. This is optional: it is included only if used; otherwise, it is null & - azimuth\_time, slant\_range\_time  for product in radar geometry 
    
    - x, y for projected products & float64 & \ding{51} \\
    tec & Link to external \acs{tec} data, if used & - & string & \\ 
    look\_vector & 3-D components radar unit look vector, in an Earth-Centred Earth-Fixed (\acs{ecef}) coordinate system. It consists of unit vectors from antenna to surface pixel. Only for reference geometry. &  (x, y) or (azimuth\_time, range\_time) & float64 & \ding{51} \\
    % los\_vector & Director cosines of the line of sight & - azimuth\_time, slant\_range\_time  for product in radar geometry 
    
    % - x, y for projected products & float32 & \ding{51} \\
    mask & Mask containing, for pixel, the following information:
    
     - 0: no distortions
    
     - 1: layover
    
     - 2: shadow
    
     - 3: undefined & slant\_range\_time  for product in radar geometry 
    
    - x, y for projected products & byte & \ding{51} \\
     
    ml\_fact & Applied range-azimuth multilook factor & slant\_range\_time  for product in radar geometry 
    
    - x, y for projected products & float32 & \ding{51} \\
    x\_sample & Pixel spacing in azimuth or North direction for radar geometry and projected products, respectively & - & float32 & \\
    y\_sample & Pixel spacing in slant range or East direction for radar geometry and projected products, respectively & - & float32 & \\
    \_geometric
    
    \_accuracy
    & Expected geometric accuracy, including eastern and northern geometric accuracy bias and standard deviation. Only for projected products. The measures will be performed on a validation dataset and periodically updated. & - & dictionary & \\ 
    east\_geom
    
    \_bias & Expected eastern geometric bias.  Only for projected products & - & float32 & \\
    east\_geom
    
    \_std & Expected eastern geometric standard deviation.  Only for projected products & - & float32 & \\
    north\_geom
    
    \_bias & Expected northern geometric bias.  Only for projected products & - & float32 & \\
    north\_geom
    
    \_std & Measured northern geometric standard deviation.  Only for projected products & - & float32 & \\
    geometric
    
    \_accuracy
    
    \_radial
    
    \_rmse	& Expected radial RMSE of geometric accuracy. Only for projected products & - & float32 & \\
    
\end{longtblr}
    
The table below describes the dimensions/coordinates associated with the interferometric products' main metadata.

\begin{longtblr}[
    label = tab:var_int_metadata_coordinates,
    entry = none,
    caption = {Interferometric products' main metadata coordinates/dimensions},
    ]{
    width = \linewidth,
    colspec = {Q[l,20]Q[l,62]Q[l,18]},
    hlines,
    hline{1-Z} = {-}{0.05em},
    }
    \textbf{Name} & \textbf{Description} & \textbf{Data Type} \\
    polarization & Acquired polarizations & string \\
    azimuth\_time & Time along satellite track direction & datetime64 \\
    range\_time & Time across satellite track direction & float64 \\
    x & UTM/UPS x (East) coordinates [m] & float64 \\
    y & UTM/UPS y (North) coordinates [m] & float64 \\
\end{longtblr}

\subsubsection{Reference System Tiling and Slicing}
\label{sec:interf_reference_system_tiling}
\autoref{tab:int_ref_system_slicing} reports the reference system, the slicing, and grid for interferometric products.

\begin{longtblr}[
    label=tab:int_ref_system_slicing,
    entry = none,
    caption = {Reference system and Slicing for interferometric products},
    ]{
        width=\linewidth,
        colspec = {Q[l,20]Q[l,35]Q[l,30]Q[l,15]},
        %vline{2-5} = {-}{},
        %vline{5-5} = {-}{},
        hlines,
        hline{1-Z} = {-}{0.05em},
        rowhead=2
    }
    \textbf{Product} & \textbf{Reference System} & \textbf{Slicing} & \textbf{Spacing}\\
    \textbf{\acs{int}+\acs{cohe}} (default) & Radar: azimuth/slant-range & Burst & ~50 m \\
    \textbf{\acs{int}+\acs{cohe}} (optional) & Radar: azimuth/slant-range & Burst & Full radar resolution \\
    \textbf{\acs{unw}+\acs{cohe}} (default) & Radar: azimuth/slant-range & Burst & ~50 m \\
    \textbf{\acs{unw}+\acs{cohe}} (optional) & Radar: azimuth/slant-range & Burst & Full radar resolution \\ 
    \textbf{\acs{gint}+\acs{gcohe}} & Cartographic: UTM/UPS & Burst-projected & 50 m \\
    \textbf{\acs{gunw}+\acs{gcohe}} & Cartographic: UTM/UPS & Burst-projected & 50 m \\
\end{longtblr}

\subsection{Product Requirements}
\subsubsection{Accuracy}
There is no standard praxis for defining requirements for phase products. The most natural approach would be to derive them from user requirements at the application level.

Requirements on product accuracy should ideally be specified in terms of bias and variance of the estimated phase or mean velocity. This is not trivial, as it depends on the target's correlation characteristics, and user requirements typically do not provide this level of detail.

Therefore, for this set of products, it is not straightforward to define numerical performance requirements. However, several key elements can be identified that critically influence the overall quality of the derived interferometric outputs.

The expected performance of interferometric products—specifically, the interferogram and coherence map—is primarily governed by three main factors: the accuracy of the co-registration process, the quality of the orbital data and Digital Elevation Model (\acs{dem}) used for phase flattening, and the effectiveness of atmospheric correction strategies in mitigating residual propagation delays.

First, achieving sub-pixel co-registration accuracy is essential to preserve phase coherence between image pairs. As a general guideline, a co-registration error lower than 10\% of a resolution cell is required to avoid significant decorrelation. Misalignments between images lead to spatial decorrelation which directly degrades the quality of interferometric measurements.

Second, the precision of the orbital state vectors and the spatial resolution and vertical accuracy of the \acs{dem} are critical for effective phase flattening. Inaccuracies in these inputs can introduce residual phase ramps or localized artifacts, reducing the fidelity of the interferometric phase and potentially compromising subsequent geophysical analyses.

Third, atmospheric disturbances—such as ionospheric phase delays in L-band systems or tropospheric variability in C-band—can introduce spatially variable phase distortions. The ability to estimate and correct these residuals, whether through external atmospheric models or empirical, data-driven approaches, is essential for preserving the spatial consistency of both the interferometric phase and the derived coherence.

Collectively, these factors determine the quality and usability of interferometric products and must be evaluated with respect to the sensor characteristics, acquisition geometry, and the requirements of the specific application.

When considering the unwrapped phase product derived from the interferogram, additional challenges arise, making it even more difficult to define numerical quality indicators for this type of \acs{ard}. The quality and reliability of the unwrapped phase are, in addition to the previously described elements, also strongly influenced by the spatial and temporal characteristics of the ground deformation being observed.

In particular, deformation signals with abrupt spatial gradients (e.g., localized subsidence, fault ruptures) or rapid temporal evolution (e.g., coseismic events or landslides) can significantly affect the robustness of phase unwrapping algorithms. Such conditions often lead to areas with high phase gradients, where the unwrapping process becomes more susceptible to errors, resulting in  discontinuities or unwrapping ambiguities.

Common unwrapping algorithms, such as the Minimum Cost Flow (MCF) and the Statistical-Cost Network-Flow Algorithm for Phase Unwrapping (SNAPHU), rely on assumptions of spatial smoothness and phase continuity. When these assumptions are violated due to complex deformation patterns or low-coherence areas, the algorithms may fail to recover a consistent and continuous deformation signal.

For all the geocoded \acs{ard}s—namely, \acs{gint}, \acs{gcohe}, and \acs{gunw}—which correspond respectively to the interferogram phase (\acs{int}), coherence (\acs{cohe}), and unwrapped phase (\acs{unw}) \acs{ard} products in radar geometry, the quality levels shall be equivalent to those of their radar-coordinate counterparts. This equivalence shall be ensured by applying a geocoding algorithm that preserves the intrinsic phase and amplitude information, reprojecting the data onto the ground reference grid without introducing radiometric or geometric distortions.

As with the previously discussed quality elements, it is also difficult to define a strict numerical requirement for geolocation accuracy, as it depends on several external factors. In particular, the requirements applicable to geolocation accuracy are directly dependent on the quality of the orbital data and the \acs{dem} used to project ground coordinates onto the radar geometry. The accuracy achievable by the geocoding procedure is linked to the precision of this mapping process, which depends on the spatial and vertical accuracy of the \acs{dem} and the precision of the orbit information employed.

\begin{itemize}
\item phase continuity between adjacent bursts
\item fringes in correspondence of flat terrain
\end{itemize}

\subsubsection{Timeliness}
In terms of timeliness, the \acs{int} product generation is closely related to the generation of \acs{cslc} products, considering all calibration operations performed at \acs{cslc} level. In fact, the successive phases of filtering, interferogram and coherence computation, and even the most demanding process of phase unwrapping, do not significantly impact the product delivery time.

On the contrary, if corrections are not applied during \acs{cslc} product generation, additional time must be considered for the retrieval of atmospheric models or the application of data-driven techniques for ionospheric phase screen removal to provide the \acs{int} product.

%\hl{(required? look at applications? where? ROSE-L MRD? NISAR?)
%ROSE-L MRD mentions now 1h latency in case of emergencies. This requirement might be very difficult to achieve.}

\subsection{Usability, Integrability and Interoperability}
\label{sec:usability_integrability_InSAR}
The characteristics of the interferometric products in terms of usability, integrability and interoperability are in the following table.

\begin{longtblr}[
  caption = {Usability, Integrability, and Interoperability of Interferometric Products},
  label = {tab:insar-usability-integrability},
]{
  colspec = {Q[2cm] Q[12cm]},
  rowhead = 1,
  width = \linewidth
}
\hline
\textbf{Aspect} & \textbf{Description} \\
\hline
Usability &
\begin{itemize}
  \item Ideal for detecting ground motion such as fast landslides, subsidence, earthquakes, glacier deformation.
  \item Sensitive to vertical and east-west displacements.
  \item Applicable at multiple scales: Suitable for both local and regional/global deformation monitoring.
  \item \acs{cohe} and \acs{gcohe} products support change detection, feature extraction, and thematic classification.
\end{itemize} 
\\
Integrability &
\begin{itemize}
  \item Interdisciplinary use: Applicable in geoscience, civil engineering, urban planning, agriculture, and disaster risk management.
  \item Workflow extensibility: Seamlessly embedded in advanced processing chains for ground deformation estimation.
\end{itemize} 
\\
Interoperability &
\begin{itemize}
  \item Geospatial alignment: Aligned to shared grids and geographic reference systems, allowing overlay with optical, DEM, and other vector layers.
  \item Cross-mission consistency: Enables joint use of products from multiple missions (e.g., Sentinel-1 and ROSE-L) using common radar grids.
\end{itemize} 
\\
\hline
\end{longtblr}

\section{Co-registered Multilooked Complex (CMLC \& GMLC)}


\subsection{Product Description}

The multi-looked phase product \acs{cmlc} is going to be generated based on \acs{slc} products and can be later ground-projected and mosaicked for improved usability, or it may be directly generated from \acs{gslc} products.

Many of the characteristics of this product are shared with the INT and GINT products described above. All model-driven and data driven corrections (atmosphere, solid earth tides, etc.) can be applied to this product and are naturally inherited from the \acs{cslc} if the \acs{cslc} already includes them. 

The target posting is 50 m, according to the ROSE-L MRD. In the case of existing triangulated phase intermediate products, the posting might be in the order of 50-200~m depending on the sensor, whereas the number of averaged pixels can vary even within the same scene, for example when pixels are selected adaptively. More recent realizations make use of machine learning techniques to estimate the phase, so that it might not even be possible to determine the number of looks in a conventional way.

It is important for this product to include a quality layer which can be very useful for the user to guide phase unwrapping. Some quality layer is normally accompanying the triangulated phases we take inspiration from: for example \cite{Ferretti2011, Ansari2018} introduces coherence indices that describe succinctly the success of the triangulation procedure. Also in the procedure described in \cite{DeZan2022} there is a possibility to measure a pair of coherence indices. However, in this case, they vary for each point and each acquisition, so they do not describe synthetically the entire time series.

It is recommended to generate \acs{cmlc} products from \acs{cslc}s defined on a reference radar grid common to Sentinel-1 and ROSE-L missions for interoperability. In any case, the geocoded version (\acs{gmlc}) will be on the same grid for Sentinel-1 and ROSE-L.

Sub-band phases would be very useful to enable ionospheric corrections based on the split-spectrum technique. This means that we should consider generating 2-3 sets of phase time series for each dataset.

\subsubsection{Input and Auxiliary Data}
The input data for \acs{cmlc} is the \acs{cslc} product in radar geometry. The geocoded \acs{gmlc} product can be derived either by geocoding the \acs{cmlc} or directly from \acs{gslc} products. The former option allows more easily to work with a fixed number of looks which is considered preferable.

It would be advisable to incorporate ETAD-like corrections in the phase. This can happen implicitly when the corrections are already included in the \acs{cslc} input products.
If a more refined version of ETAD-like corrections is needed for high-end applications (e.g. precise strain measurement over large distances), this can be done at a later stage. To this end, it should be possible to undo the applied corrections.


\begin{longtblr}[
	label = tab:in_aux_mlc_products,
	entry = none,
	caption = {\acs{mlc} input and auxiliary data},
	]{
		width = \linewidth,
		colspec = {Q[l,20]Q[l,60]Q[c,20]},
		hlines,
		hline{1-Z} = {-}{0.05em},
	}
	\textbf{Name}   & \textbf{Description}   & \textbf{Mandatory (M) / Optional (O)} \\
	\acs{cslc} or \acs{gslc} & Co-registered (preferred) or geocoded L1 Single Look Complex image and annotations & M \\
    %\acs{dem} & External DEM (e.g., Copernicus DEM at different resolutions, with 30 m as the default worldwide). & M \\
    %Precise Orbits & \acs{slc} precise orbits & M \\
    \acs{etad} & \acs{etad}-like products for atmospheric delay and timing corrections & O \\
    \acs{tec} & External \acs{tec} data for delay correction, if not already included as an \acs{etad} layer & O \\
\end{longtblr}


\subsubsection{Main Variables}

The following tables (\ref{tab:var_cmlc_products} and \ref{tab:var_gmlc_products}) collect the main variables for the \acs{cmlc} and \acs{gmlc} products.

\begin{longtblr}[
	label = tab:var_cmlc_products,
	entry = none,
	caption = {\acs{cmlc} main variables},
	]{
		width = \linewidth,
		colspec = {Q[l,13]Q[l,42]Q[l,19]Q[l,14]Q[l,12]},
		hlines,
		hline{1-Z} = {-}{0.05em},
	}
	\textbf{Name} & \textbf{Description} & \textbf{Dimensions} & \textbf{Data Type} \\
	phase & Co-registered multi-look complex data.
       
    The phase corresponds conceptually to the interferometric phase between each acquisition and a fixed reference (e.g. the first acquisition in the stack) that one could compute directly by forming an interferogram with the reference. In practice, the \acs{mlc} phase should be less noisy. The \acs{mlc} phase can be used to generate directly multi-looked interferograms between any pair of acquisitions.
    & polarization, acquisition\_date, azimuth\_time, range\_time & complex64 (or float64) \\
    long-term coherence & The long-term coherence of each acquisition with a past reference. This information can be useful as an indication of the noisiness of interferograms computed between acquisitions that are separated by a large time interval (e.g. 1 year or more). The exact interpretation depends on the chosen processing parameters. & polarization, acquisition\_date, azimuth\_time, range\_time & float32 \\
    short-term coherence & The short-term coherence of each acquisition with the current reference (recent past). This information can be useful as an indication of the noisiness of interferograms computed between acquisitions that are separated by a short time interval (e.g. a few weeks), f.i. to guide phase unwrapping. The exact interpretation depends on the chosen processing parameters. & polarization, acquisition\_date, azimuth\_time, range\_time & float32 \\
\end{longtblr}

\begin{longtblr}[
	label = tab:var_gmlc_products,
	entry = none,
	caption = {\acs{gmlc} main variables},
	]{
		width = \linewidth,
		%colspec = {Q[l,15]Q[l,37]Q[l,20]Q[l,15]Q[l,13]},
        colspec = {Q[l,13]Q[l,42]Q[l,19]Q[l,14]Q[l,12]},
		hlines,
		hline{1-Z} = {-}{0.05em},
	}
	\textbf{Name} & \textbf{Description} & \textbf{Dimensions} & \textbf{Data Type} \\
	phase & Geocoded multi-look complex data:
       
    The phase corresponds conceptually to the interferometric phase between each acquisition and a fixed reference (e.g. the first acquisition in the stack) that one could compute directly by forming an interferogram with the reference. In practice, the \acs{mlc} phase should be less noisy. The \acs{mlc} phase can be used to generate directly multi-looked interferograms between any pair of acquisitions.
    & polarization, acquisition\_date, x, y & complex64 (or float64) \\
    long-term coherence & The long-term coherence of each acquisition with a past reference. This information can be useful as an indication of the noisiness of interferograms computed between acquisitions that are separated by a large time interval (e.g. 1 year or more). The exact interpretation depends on the chosen processing parameters. & polarization, acquisition\_date, x, y & float16 \\
    short-term coherence & The short-term coherence of each acquisition with the current reference (recent past). This information can be useful as an indication of the noisiness of interferograms computed between acquisitions that are separated by a short time interval (e.g. a few weeks), f.i. to guide phase unwrapping. The exact interpretation depends on the chosen processing parameters. & polarization, acquisition\_date, x, y & float16 \\
\end{longtblr}

The table below describes the dimensions/coordinates associated with the \acs{mlc} main variables.

\begin{longtblr}[
	label = tab:var_mlc_coordinates,
	entry = none,
	caption = {\acs{mlc} main variables coordinates/dimensions},
	]{
		width = \linewidth,
		colspec = {Q[l,20]Q[l,62]Q[l,18]},
		hlines,
		hline{1-Z} = {-}{0.05em},
	}
	\textbf{Name} & \textbf{Description} & \textbf{Data Type} \\
	polarization & Acquired polarizations & string \\
    band & Acquired band (e.g. main or auxiliary band) & string \\
    acquisition\_date & Image acquisition dates, formatted as YYYY-MM-DD & string \\
    azimuth\_time & Time along satellite track direction & datetime64 \\
    range\_time & Time across satellite track direction & float64 \\
    x & UTM/UPS x (East) coordinates [m] & float64 \\
    y & UTM/UPS y (North) coordinates [m] & float64 \\
\end{longtblr}


\subsubsection{Main Metadata}

The following table lists the main metadata and per-pixel metadata associated with the \acs{mlc} product. This product is not contemplated in \acs{ceos} documents nor in current SAR missions.  Many of the metadata are going to be identical for all the products in a stack and should not be repeated for each acquisition.

\begin{longtblr}[
	label = tab:meta_mlc_products,
	entry = none,
	caption = {\acs{cmlc}/\acs{gmlc} main metadata and per-pixel metadata},
	]{
		width = \linewidth,
		colspec = {Q[l,17]Q[l,38]Q[l,20]Q[l,13]Q[c,12]},
		hlines,
		hline{1-Z} = {-}{0.05em},
	}
	\textbf{Name} & \textbf{Description} & \textbf{Dimensions} & \textbf{Data Type} & \textbf{Per-pixel Metadata}\\
	input\_slc  & List of links to input reference \acs{slc} products & - & list of strings & \\
    reference\_dem & Link to reference DEM & - & string & \\
    phase\_demod & Phase to be applied to demodulate each image belonging to the \acs{slc} stack, computed using DEM & (x, y) or (azimuth\_time, range\_time) & float64 & \ding{51} \\ 
    Ionospheric phase & The ionospheric phase estimated with split-spectrum and subtracted from the phase. This variable allows undoing the correction for an advanced user who would want to utilize a different correction & (x, y) or (azimuth\_time, range\_time) & float32 & \ding{51} \\
    etad\_file & ETAD file used for each image belonging to the \acs{cslc}/\acs{slc} stack. This is optional: it is included only if used; otherwise, it is null & - & string & \\ 
    tec & Link to external \acs{tec} data, if used & - & string & \\ 
    look\_vector & 3-D components radar unit look vector, in an Earth-Centred Earth-Fixed (\acs{ecef}) coordinate system. It consists of unit vectors from antenna to surface pixel. Only for reference geometry. &  (x, y) or (azimuth\_time, range\_time) & float64 & \ding{51} \\
    inc\_angle & List of links to reference DEM-based local incidence angle maps. & - & string & \\

    
    %reference\_anf & Optional attribute: list of links to reference area normalization factor or local contributing area. & - & string & \\
    mask & Mask containing, for each ground cell, the following information:
    
     - 0: no distortions
    
     - 1: layover
    
     - 2: shadow
    
     - 3: undefined & (x, y) or (azimuth\_time, range\_time) & byte & \ding{51} \\
    nlooks & Number of equivalent looks for interferometric phase computation (only for reference geometry) & (x, y) or (azimuth\_time, range\_time) & float32 & \ding{51} \\
    
    processing\_
    
    parameters  & List of processing parameters & - & list of strings & \\
    units & Units of the variables of measurements.  & - & list of strings & \\
    
\end{longtblr}



\subsubsection{Reference System Tiling and Slicing}
\label{sec:mlc_reference_system_tiling}
\autoref{tab:mlc_ref_system_slicing} reports the reference system, the slicing, and grid for "Multi-Look Complex" products.


\begin{longtblr}[
	label=tab:mlc_ref_system_slicing,
	entry = none,
	caption = {Reference system and Slicing for \acs{cmlc} and \acs{gmlc}},
	]{
		width=\linewidth,
		colspec = {Q[l,15]Q[l,40]Q[l,15]Q[l,30]},
		%vline{2-5} = {-}{},
		%vline{5-5} = {-}{},
		hlines,
		hline{1-Z} = {-}{0.05em},
		rowhead=2,
	}
	\textbf{Product} & \textbf{Reference System} & \textbf{Slicing} & \textbf{Spacing}\\
	\textbf{\acs{cslc}} & Radar: azimuth/slant-range & Burst & sub-grid of original \acs{slc} radar grid, approximately 50 m \\
	\textbf{\acs{gslc}} & Cartographic: UTM/UPS & Burst & 50 m \\   
	
\end{longtblr}

\subsection{Product Requirements}

\subsubsection{Accuracy}
There is no standard praxis around the definition of requirements for multi-looked phase products. The most natural way of defining requirements would be to derive them from user requirements at the application level. 

Requirements on the product accuracy should ideally be specified in terms of bias and variance on the estimated phase or mean velocity. This is not trivial as it depends on the target correlation characteristics and user requirements do not have this level of detail.

One possibility is to define requirements with reference to a standardized coherence scenario that includes decorrelation, closure phases, and long-term coherence. 

Another possibility, which we prefer here, is to define the performance requirements relative to a full-covariance matrix estimator such as EMI or relative to an ideal performance such as the one predicted by the Cram\'er-Rao bound.

We are going to ask that the phase bias relative to a reference full-covariance algorithm should be smaller than 1 mm / year (TBC), which we set looking at the deformation precision requirement. 

The standard deviation degradation with respect to a reference full-covariance algorithm should be smaller than 3 dB (TBC).

The generation of this product requires a precise co-registration of the \acs{cslc} input products. It is important that the co-registration offsets are annotated in the \acs{cslc} and that it is possible to distinguish the geometrical/modeled part from the data-driven part (estimated typically with local cross-correlation).


\begin{longtblr}[
  label = tab:mlc_accuracy,
  entry = none,
  caption = {Accuracy Requirements for \acs{mlc} products},
]{
  width = \linewidth,
  colspec = {Q[l,35] Q[l,55] Q[c,20]},
  hlines,
  hline{1-Z} = {-}{0.05em},
  rowhead = 1,
}
\textbf{Metric} & \textbf{Description} & \textbf{Requirement} \\
$v_{\mathrm{bias}}$ & Deformation velocity bias relative to reference full-covariance method& $\leq 1$ mm/year \\
$20 \log_{10}(\sigma_\varphi/\sigma_\mathrm{ref})$ &  Phase std degradation relative to reference full-covariance method & $\leq 3$ dB \\
\end{longtblr}


\subsubsection{Timeliness}
The timeliness requirements depend on the applications which are intended to be covered by this product. If we ignore emergency requirements, assuming they can be covered by nearest-pair interferograms, the \acs{mlc} product can have a more relaxed latency requirement of e.g. 1 day or 1 week. For precise tectonic measurements, it is worth including the best reanalysis tropospheric corrections which are currently available only with a latency of a few months. %If there is access to predicted troposphere layers from ECMWF, one should quantify the degradation with respect to reanalysis datasets to determine whether they are sufficient for applications.

One could therefore foresee two releases of the product, the second being only a correction of the phase variable - not a full reprocessing. A first release would be within one to two days after the acquisition, to be able to include quality orbit information.
A second release would be as soon as final numerical weather reanalysis data are available (currently, three months for ERA-5 reanalysis).

\autoref{fig:cmlc_from_cslc} shows the flowchart for generating the \acs{cmlc} and \acs{gmlc} products from \acs{cslc} input.

\begin{figure}[H]
		\centering
		\resizebox{\textwidth}{!}{
		\begin{tikzpicture}[node distance=0.5cm and 2cm]
			
			% Define the style for the nodes
			% \tikzstyle{node} = [circle, draw, fill=blue!20, text centered, minimum size=1.5cm]
			
			% Draw nodes
			\node[rectin] (init) {\acs{cslc}};
			\node[block, right=of init] (phase_est) {Phase estim.};
			\node[rectout, right=of phase_est] (cmlc) {\acs{cmlc}};
			\node[block, right=of cmlc] (refine) {Refinement};
			\node[block, below=of cmlc] (geocode) {Geocoding};
			\node[rectout, below=of geocode] (gmlc) {\acs{gmlc}};
			\node[rectout, right=of refine] (cmlc2) {\acs{cmlc}};
			\node[rectaux, below=of refine] (models) {External models};
			\node[block, right=of gmlc] (refine2) {Refinement};
			\node[rectout, right=of refine2] (gmlc2) {\acs{gmlc}};
			
			% Draw edges with arrows
			\draw[->, thick] (init) -- (phase_est);
			\draw[->, thick] (phase_est) -- (cmlc);
			\draw[->, thick] (cmlc) -- (refine);
			\draw[->, thick] (refine) -- (cmlc2);
			\draw[->, thick] (models) -- (refine);
			\draw[->, thick] (cmlc) -- (geocode);
			\draw[->, thick] (geocode) -- (gmlc);
			\draw[->, thick] (models) -- (refine2);
			\draw[->, thick] (gmlc) -- (refine2);
			\draw[->, thick] (refine2) -- (gmlc2);
			
			% Draw edges with arrows
			%\draw[->, thick] (dem) -| (coreg);
			%\draw[->, thick] (orbit) -- (coreg);
			\draw[dotted, thick] (gmlc2) -- (cmlc2);
						
			% Timeline with linked markers and adaptive length
			\def\timelineStart{init}
			\def\timelineEnd{cmlc2}
			\draw[thick, ->] (\timelineStart.west |- 0,-6) -- (\timelineEnd.east |- 0,-6) node[right] {Time};
			\foreach \node/\label in {init/Start+(1-12h), phase_est/{<1h}, refine2/20-90d, gmlc2/<1h} {
				\draw (\node.south |- 0,-6) -- (\node.south |- 0,-6.5) node[below] {\label};
				\draw[dotted, thick] (\node.south) -- (\node.south |- 0, -6.25); % Connects the node to the timeline
			}
			
		\end{tikzpicture}
		}
		\caption{Sample timeline for the generation of \acs{cmlc} and \acs{gmlc} starting from \acs{cslc} products.}
		\label{fig:cmlc_from_cslc}
\end{figure}


\subsection {Usability, Integrability and Interoperability}

\autoref{tab:mlc_characteristics} summarizes usage characteristics of \acs{cmlc} and \acs{gmlc} products.
\begin{longtblr}[
  caption = {Usability, Integrability, and Interoperability of "Multi-looked Complex" Products},
  label = {tab:mlc_characteristics},
]{
  colspec = {Q[2cm] Q[12cm]},
  rowhead = 1,
  width = \linewidth
}
\hline
\textbf{Aspect} & \textbf{Description} \\
\hline
Usability &
\begin{itemize}
  \item Displacement Monitoring: useful for detecting large displacements such as landslides, earthquakes, mining subsidence, etc. with InSAR sensitivity. Usability to monitor ice motion needs to be demonstrated.
  \item Time series: naturally suited for time-series analysis of surface displacement.
  \item Advanced phase estimation: it is robust to deformation bias so that the user does not have to be a  specialist in phase estimation.
  \item Global to Local Monitoring: Suitable for regional or global-scale observations (e.g., tectonics) and local events (e.g., fast landslides).
  \item The \acs{mlc} wrapped phase can even be fed directly to a machine learning model trained for motion detection or parameter inversion~\cite{Biggs2022} or, possibly, phase unwrapping.
  \item Broad application domain: Useful in natural hazard analysis, volcanology, earthquake monitoring, urban subsidence, hydrology, etc.
\end{itemize}
\\
Integrability &
\begin{itemize}
  \item GIS Compatibility: Easily used within \acs{gis} platforms, due to geographic referencing and data format.
  \item Multi-Sensor Fusion: Integrates well with optical or other SAR projected products for land cover, surface change, or hazard mapping.
  \item Multi-Sensor compatibility: Can be derived from various sensors and integrated in joint analyses.
  \item Multi-sensor phase unwrapping: the \acs{gmlc} product can be used in a joint search for a phase-unwrapping solution, which can benefit from wavelength diversity.
\end{itemize}
\\
Interoperability &
\begin{itemize}
  \item Sensor-Independent Formats: Phase values projected geometry can be standardized across sensors for interoperability in displacement workflows.
  \item Geolocation Harmonization: When geocoded, products align spatially with other \acs{eo} datasets (e.g., NRB, optical, DEMs) to support multi-source analysis.
  \item Software Ecosystem Support: Usable in major \acs{eo} software platforms, supporting scientific and operational pipelines.
\end{itemize}
\\
\hline
\end{longtblr}



\section{Enhanced NRB \& ORB (ENRB \& EORB)}

\subsection{Product Description}
\label{sec:enrb_eorb_product_description}
As is well known, SAR images are subject to various artifacts and noise sources that can hinder the implementation of many high-level applications. The Enhanced \acs{nrb} and \acs{orb} products are designed to retain the information content of their corresponding standard products, while mitigating these disturbances to improve overall image quality and usability.

These "enhanced" versions of the \acs{nrb} and \acs{orb} \acs{ard} products are generated by applying specific filtering procedures to the input \acs{slc} or \acs{cslc} data:

\begin{itemize}
    \item \textit{Ghost Removal:} Ghosts are spurious image artifacts caused by SAR ambiguities in the range or azimuth directions. These artifacts can significantly degrade application performance, particularly in tasks such as object detection and recognition. Ghost removal is especially beneficial for applications that utilize the \acs{orb} product.
    \item \textit{Sidelobe Suppression:} Sidelobes in SAR imagery may introduce two critical issues. First, \textit{false alarms} can occur when strong scatterers produce sidelobes that are mistakenly interpreted as genuine targets. Second, \textit{missed detections} may arise when the sidelobes of a strong target obscure the mainlobe of a nearby weaker target. Sidelobe suppression improves the reliability of detection, particularly in scenarios requiring either the \acs{orb} or \acs{nrb} products.
\end{itemize}

With the exception of the intensity cross-spectra, the Enhanced \acs{nrb} and \acs{orb} products closely match their standard counterparts (see \autoref{sec:nrb_prod_description} and \autoref{sec:orb_product}), but with additional processing steps aimed at reducing ghosting and sidelobe interference as much as possible.

The proposed intensity cross-spectra component follows the specification of Harmony’s Intensity Cross-Spectral Matrix (ICSM) product. For each polarization, the ICSM provides both intensity power spectra and cross-spectra, enhancing the frequency-domain analysis of the SAR backscatter.

\subsubsection{Input and Auxiliary Data}
The following input and auxiliary data are considered for product generation:

\begin{longtblr}[
    label = tab:in_aux_enh_products,
    entry = none,
    caption = {Enhanced \acs{nrb} and \acs{orb} input and auxiliary data},
    ]{
        width = \linewidth,
        colspec = {Q[l,20]Q[l,60]Q[c,20]},
        hlines,
        hline{1-Z} = {-}{0.05em},
    }
    \textbf{Name}   & \textbf{Description}   & \textbf{Mandatory (M) / Optional (O)} \\
    \acs{slc} or \acs{cslc} & L1 Single Look Complex image or co-registered L1 Single Look Complex image & M \\
    \acs{slc} annotation & L1 Single Look Complex image annotations or reference geometry annotations & M \\
    \acs{dem} & External DEM (e.g., Copernicus DEM at different resolutions, with 30 m as the default worldwide). & M for \acs{nrb} \\
    Precise Orbits & \acs{slc} precise orbits & M \\
    \acs{etad} & \acs{etad}-like products for atmospheric delay correction & O \\
    \acs{tec} & External \acs{tec} data for delay correction & O \\
    Geoid & External geoid & M for \acs{orb} \\
    \acs{era5} & ECMWF \acs{era5} wind data for angle of incidence compensation & M for \acs{orb} \\
\end{longtblr}

\subsubsection{Main Variables}
\label{sec:variables_enh}
The following tables list the main variables associated with the \acs{enrb} product.

\begin{longtblr}[
    label = tab:var_enrb_products,
    entry = none,
    caption = {\acs{enrb} main variables},
    ]{
        width = \linewidth,
        colspec = {Q[l,25]Q[l,52]Q[l,20]Q[l,13]},
        hlines,
        hline{1-Z} = {-}{0.05em},
    }
    \textbf{Name} & \textbf{Description} & \textbf{Dimensions} & \textbf{Data Type} \\
    gamma0 & Radiometric Terrain Correction gamma naught backscatter linear scaling image. Only for \acs{enrb} & polarization, acquisition\_date, x, y & float32 \\
    sigma0\_raw ($\sigma_{0,\mathrm{raw}}$) & Full resolution (single look) radar intensity, in $\sigma_0$ convention. Only for \acs{eorb}  & polarization, acquisition\_date, x, y & float32 \\
    sigma0\_nesz ($\sigma_{0,\mathrm{nesz}}$) & Spatially-averaged, LxL m$^2$ resolution noise equivalent sigma zero, in $\sigma_0$ convention. Only for \acs{eorb} & polarization, acquisition\_date, x, y & float32 \\
    sigma0\_denoise ($\sigma_{0,\mathrm{denoise}}$) & Spatially-averaged, LxL m$^2$ resolution denoized radar intensity, in $\sigma_0$ convention. Only for \acs{eorb} & polarization, acquisition\_date, x, y & float32 \\
    sigma0\_L ($\sigma_{0,L}$) & Spatially-averaged, LxL m$^2$ resolution denoized and outlier removed radar intensity, in $\sigma_0$ convention. Only for \acs{eorb} & polarization, acquisition\_date, x, y & float32 \\
    sigma0\_Lnorm ($\sigma_{0,L}/\sigma_{0,\mathrm{mod}}$) & Spatially-averaged, LxL m$^2$ resolution (mean) wind normalized backscatter measurements. Only for \acs{eorb} & polarization, acquisition\_date, x, y & float32 \\
    local\_inc\_angle & DEM-based local Incidence angle image [deg], only for \acs{enrb}. This variable is present only if the \acs{enrb} is generated from \acs{slc}s. If the \acs{enrb} is derived from \acs{cslc}s, this variable will be null, and the reference\_inc\_angle metadata will be present & x, y & float32 \\
    anf & Optional variable: area normalization factor or local contributing area, used to obtain the gamma0 coefficient; only for \acs{enrb}. This variable is present only if the \acs{enrb} is generated from \acs{slc}s. If the \acs{nrb} is derived from \acs{cslc}s, this variable will be null, and the reference\_anf metadata will be present & x, y & float32 \\
    mask & Full resolution (single look) bright target mask. Only for \acs{eorb} & polarization, acquisition\_date, x, y & bool \\
\end{longtblr}

The table below describes the dimensions/coordinates associated with the \acs{enrb} and \acs{eorb} main variables.

\begin{longtblr}[
    label = tab:var_enh_coordinates,
    entry = none,
    caption = {\acs{enrb} and \acs{eorb} main variables coordinates/dimensions},
    ]{
        width = \linewidth,
        colspec = {Q[l,20]Q[l,62]Q[l,18]},
        hlines,
        hline{1-Z} = {-}{0.05em},
    }
    \textbf{Name} & \textbf{Description} & \textbf{Data Type} \\
    polarization & Acquired polarizaions & string \\
    acquisition\_date & Image acquisition dates, formatted as YYYY-MM-DD & string \\
    x & UTM/UPS x (East) coordinates [m] & float64 \\
    y & UTM/UPS y (North) coordinates [m] & float64 \\
\end{longtblr}

\subsubsection{Main Metadata}
\label{sec:enh_product_metadata}
The following table lists the main metadata and per-pixel metadata associated with the \acs{enrb} and  \acs{eorb} products.

\begin{longtblr}[
    label = tab:meta_enh_products,
    entry = none,
    caption = { \acs{enrb} and  \acs{eorb} main metadata and per-pixel metadata},
    ]{
        width = \linewidth,
        colspec = {Q[l,25]Q[l,38]Q[l,10]Q[l,10]Q[c,12]},
        hlines,
        hline{1-Z} = {-}{0.05em},
    }
    \textbf{Name} & \textbf{Description} & \textbf{Dimensions} & \textbf{Data Type} & \textbf{Per-pixel Metadata}\\
    reference\_slc / reference\_cslc & List of links to input reference \acs{slc} or \acs{cslc} & - & list of string & \\
    reference\_dem & Link to reference DEM. Only for \acs{enrb} & - & string & \\
    reference\_geoid & Link to reference GEOID. Only for \acs{eorb} & - & string & \\
    etad\_file & ETAD file used for each image belonging to the \acs{cslc}/\acs{slc} stack. This is optional: it is included only if used; otherwise, it is null & - & string & \\ 
    tec & Link to external \acs{tec} data, if used. & - & string & \\ 
    era5 & Link to external \acs{era5} data, if used.  & - & string & \\
    reference\_inc\_angle & Optional attribute: list of link to reference DEM- or GEOID-based local incidence angle maps. It is present only if the \acs{nrb} is generated from \acs{cslc}s. & - & string & \ding{51} \\
    ground\_heading & Azimuth direction. Only for \acs{eorb} & x, y & float32 & \ding{51} \\
    reference\_anf & Optional attribute: list of links to reference area normalization factor or local contributing area; only for \acs{enrb}. It is present only if the \acs{enrb} is generated from \acs{cslc}s. & - & string & \\
    orig\_data\_id & Per pixel metadata reporting an index describing the bursts contributing to the output mosaicked pixel & x, y & int32 & \ding{51} \\
    mask & Mask containing, for each ground cell, the following information:
    
     - 0: no distortions
    
     - 1: layover
    
     - 2: shadow
    
     - 3: undefined.

     Only for \acs{enrb}
     & x, y & byte & \ding{51} \\
    nlooks & Number of equivalent looks for the gamma flattening & x, y & float32 & \ding{51} \\ 
    \_geometric\_accuracy & Measured geometric accuracy, including eastern and northern geometric accuracy bias and standard deviation (only for \acs{enrb}) & - & dictionary & \\ 
    east\_geom\_bias & Measured eastern geometric bias (only for \acs{enrb}) & - & float32 & \\
    east\_geom\_std & Measured eastern geometric standard deviation (only for \acs{enrb}) & - & float32 & \\
    north\_geom\_bias & Measured northern geometric bias (only for \acs{enrb}) & - & float32 & \\
    north\_geom\_std & Measured northern geometric standard deviation (only for \acs{enrb}) & - & float32 & \\
    geometric\_accuracy
    \_radial\_rmse& Measured radial RMSE of geometric accuracy (only for \acs{enrb}) & - & float32 & \\
    x\_sample & Pixel spacing in x direction & - & float32 & \\
    y\_sample & Pixel spacing in y direction & - & float32 & \\
    ghost\_rm\_flag & Flag indicating the application of ghosts removal & - & string & \\
    ghost\_rm
    
    \_algo & Link to the algorithm used for ghosts removal & - & string & \\
    side\_lobe
    \_reduction\_flag & Flag indicating the application of side lode reduction & - & string & \\
    side\_lobe
    
    \_reduction
    
    \_algo & Link to the algorithm used for side lobe reduction & - & string & \\
    
\end{longtblr}

The table below describes the dimensions/coordinates associated with the \acs{enrb} and \acs{eorb} main metadata.

\begin{longtblr}[
    label = tab:var_enh_metadata_coordinates,
    entry = none,
    caption = {\acs{enrb} and \acs{eorb} main metadata coordinates/dimensions},
    ]{
        width = \linewidth,
        colspec = {Q[l,20]Q[l,62]Q[l,18]},
        hlines,
        hline{1-Z} = {-}{0.05em},
    }
    \textbf{Name} & \textbf{Description} & \textbf{Data Type} \\
    x & UTM/UPS x (East) coordinates [m] & float64 \\
    y & UTM/UPS y (North) coordinates [m] & float64 \\
\end{longtblr}

\subsubsection{Reference System Tiling and Slicing}
\label{sec:enh_reference_system_tiling}

\autoref{tab:enh_ref_system_slicing} reports the reference system, the slicing, and grid for \acs{enrb} and \acs{eorb} products.

\begin{longtblr}[
    label=tab:enh_ref_system_slicing,
    entry = none,
    caption = {Reference system and Slicing for \acs{enrb} and \acs{eorb}},
    ]{
        width=\linewidth,
        colspec = {Q[l,15]Q[l,35]Q[l,35]Q[l,15]},
        %vline{2-5} = {-}{},
        %vline{5-5} = {-}{},
        hlines,
        hline{1-Z} = {-}{0.05em},
        rowhead=1
    }
    \textbf{Product} & \textbf{Reference System} & \textbf{Slicing} & \textbf{Spacing}\\
    \textbf{\acs{enrb}} (default) & Cartographic: UTM/UPS & US-\acs{mgrs} 100 km x 100 km & 10 m\\
    \textbf{\acs{enrb}} (optional) & Cartographic: UTM/UPS & US-\acs{mgrs} 100 km x 100 km & 30 m\\
    \textbf{\acs{eorb}} (default) & Cartographic: UTM/UPS & US-\acs{mgrs} 100 km x 100 km & 10 m\\
    \textbf{\acs{eorb}} (optional) & Cartographic: UTM/UPS & US-\acs{mgrs} 100 km x 100 km & 30 m\\	
\end{longtblr}

\subsection{Product Requirements}
No additional criticalities are foreseen beyond those listed for standard \acs{nrb} or \acs{orb} products.

\subsubsection{Accuracy}

% La risoluzione 

As regards Ghost Removal, two approaches have been evaluated (both the analysed methods base on the asymmetry of the Doppler Spectrum of \acs{sar} ghosts):
\begin{itemize}
    \item The “asymmetric mapping and selective filtering” methods described in \cite{Avolio_et_al2014} and \cite{DiMartino_et_al2014} mitigates azimuth ambiguities through a two-step process: asymmetric Wiener filters are applied to isolate ambiguity-free spectral regions, and binary maps identify affected areas. Only those ambiguous regions are selectively replaced with their filtered counterparts, preserving the rest of the image.
    \item The method in \cite{Choi2021} exploits the spectral distinction between azimuth ambiguities and genuine targets: while true targets exhibit broad spectral support, ghost signals are narrowly concentrated near the Doppler Nyquist frequency. By forming three sub-aperture images and projecting them onto a rotated coordinate system, the method concentrates target power into a single component while spreading ghost energy across all components, enabling effective ambiguity suppression through simple recombination.
\end{itemize}
This technique is often seen as an intermediate step for further analyses, such as detection (especially in maritime environments) or classification, but some metrics can help assessing the performances. The following table summarizes the accuracy requirements for the Ghost Removal technique for \acs{enrb} and \acs{eorb} products, considering coastal environments as test site.

\begin{longtblr}[
	label = tab:enrb_ghost_accuracy,
	entry = none,
	caption = {Accuracy Requirements for \acs{enrb} / \acs{eorb} products, regarding Ghost Removal (Azimuth and Range Ambiguities)},
]{
	width = \linewidth,
	colspec = {Q[l,30] Q[l,60] Q[c,20]},
	hlines,
	hline{1-Z} = {-}{0.05em},
	rowhead = 1,
}

\textbf{Metric} & \textbf{Description} & \textbf{Requirement} \\

\textbf{ASR / RSR (dB)} & Ambiguity-to-Signal Ratio in azimuth (ASR) or range (RSR), defined as the ratio of average power in ambiguity regions (typically displaced by integer multiples of PRF or sampling interval) to the mainlobe peak power. & $\leq -28$~dB \\

\textbf{GBR Reduction (dB)} & Difference in ghost-to-background ratio before and after suppression. GBR is computed as the ratio between the average power in ambiguity-affected regions and the average power in nearby background (clutter or noise) regions. & $\geq 5$~dB \\

\end{longtblr}

As regards Sidelobe Suppression techniques, Spatially Variant Apodization (SVA) technique (and its refinements) have been considered, due to their capacity to suppress sidelobes while preserving the resolution for strong scatterers, following \cite{Fischer2006} and \cite{Liu2018}. The following table summarizes the accuracy requirements for the Sidelobe Suppression technique for \acs{enrb} and \acs{eorb} products, considering a strong pointlike scatterer in a uniform region as test site.

\begin{longtblr}[
	label = tab:enrb_sidelobe_accuracy,
	entry = none,
	caption = {Accuracy Requirements for \acs{enrb} / \acs{eorb} products, regarding Sidelobe Suppression application},
]{
	width = \linewidth,
	colspec = {Q[l,30] Q[l,60] Q[c,20]},
	hlines,
	hline{1-Z} = {-}{0.05em},
	rowhead = 1,
}

\textbf{Metric} & \textbf{Description} & \textbf{Requirement} \\

\textbf{PSLR (dB)} & Peak sidelobe level relative to mainlobe peak & $\leq -25$~dB \\

\textbf{ISLR (dB)} & Total sidelobe energy relative to mainlobe energy & $\leq -10$~dB \\

\textbf{FWHM Deviation (\%)} & Full Width at Half Maximum variation from theoretical mainlobe width & $\leq 10\%$ \\

\textbf{SCNR (dB)} & Signal-to-clutter-plus-noise ratio in target region & $\geq 10$~dB \\

\textbf{Sidelobe Floor (dB)} & Average level of sidelobe floor excluding peak sidelobes & $\leq -30$~dB \\

\end{longtblr}

\subsubsection{Timeliness}

\acs{enrb} and \acs{eorb} products can be generated as soon as L1 data are available, for each \acs{slc} image independently or starting from \acs{cslc} products, with the addition of the Enhancing step.

The Enhancing step can be performed on \acs{slc} or \acs{cslc} products, and does not require any other input data. Additionally, for the timeliness of the identified algorithms, the following considerations are valid:
\begin{itemize}
    \item Ghost Removal: the increment of timeliness is negligible with respect to standard \acs{nrb} or \acs{orb} production.
    \item Sidelobe Suppression: it is of the same order as that required to generate the corresponding standard products. 
\end{itemize}
The composition of the Enhancing step is different for \acs{enrb} and \acs{eorb} products, but considerations about timeliness are the same. The only difference involves Ghost Removal, which is an operation useful only in the maritime environment.

In the framework of this project, the \acs{cslc} procedure includes geometric correctness, even in case of parameter compensation that does not introduce differential errors between reference and secondary images. Taking into account this characteristic, if the \acs{enrb} is generated from the \acs{cslc}, all the calibrations are already performed during the \acs{cslc} processing, and the reference geometry is computed only once for the reference orbit. Therefore, the timeliness does not change depending on the information needed to perform the Geo-to-Radar conversion or to generate the \acs{cslc} product.

The main factor impacting the timeline is instead the availability of external data or reference images, which are required to refine geolocation and mitigate atmospheric effects, particularly ionospheric effects. 

\autoref{fig:enrb_single} shows the flowchart for generating the \acs{enrb}/\acs{eorb} product using a single \acs{slc} image.

\begin{figure}[H]
	\centering
	\resizebox{\textwidth}{!}{
		\begin{tikzpicture}[node distance=0.5cm and 2cm]
			
			% Define the style for the nodes
			% \tikzstyle{node} = [circle, draw, fill=blue!20, text centered, minimum size=1.5cm]
			
			% Draw nodes
			\node[rectin] (init) {\acs{slc}};
                \node[block, right=of init] (enh) {Enhancing};
			\node[block, right=of enh] (g2r) {Geo-to-Radar};
			\node[block, right=of g2r] (rtc) {\acs{rtc}};
			\node[rectaux, below=of init] (dem) {\acs{dem}};
			\node[rectaux, below=of dem] (orbit) {Orbital data};
			\node[rectauxd, below=of orbit] (models) {External models};
			\node[rectind, below=of models] (ref_stack) {Reference SLC stack};
			\node[block, below=of rtc] (geocoding) {Geocoding};
                \node[rectout, below=of geocoding] (nrb) {\acs{enrb}/\acs{eorb}};
			
			% Draw edges with arrows
			\draw[->, thick] (init) -- (enh);
                \draw[->, thick] (enh) -- (g2r);
			\draw[->, thick] (g2r) -- (rtc);
			\draw[->, thick] (rtc) -- (geocoding);
			\draw[->, thick] (dem) -| (g2r);
			\draw[->, thick] (orbit) -| (g2r);
			\draw[->, thick] (models) -| (g2r);
			\draw[->, thick] (ref_stack) -| (g2r);
			\draw[->, thick] (geocoding) -- (nrb);
		\end{tikzpicture}
	}
	\caption{Sample workflow for generation of \acs{enrb} considering each \acs{slc} image separately.}
	\label{fig:enrb_single}
\end{figure}

\autoref{fig:enrb_from_cslc} shows the flowchart for generating the \acs{enrb}/\acs{eorb} product using \acs{cslc} products.

\begin{figure}[H]
	\centering
	\resizebox{\textwidth}{!}{
		\begin{tikzpicture}[node distance=0.5cm and 2cm]
			
			% Define the style for the nodes
			% \tikzstyle{node} = [circle, draw, fill=blue!20, text centered, minimum size=1.5cm]
			
			% Draw nodes
			\node[rectin] (init) {\acs{slc}};
            \node[block, right=of init] (coreg) {Co-registration};
            \node[rectout, right=of coreg] (cslc) {\acs{cslc}};
            \node[block, right=of cslc] (enh) {Enhancing};
            \node[block, below=of enh] (rtc) {\acs{rtc}};    
            \node[block, below=of rtc] (geocoding) {Geocoding};
    
            \node[rectaux, below=of init] (dem) {\acs{dem}};
            \node[rectaux, below=of dem] (orbit) {Orbital data};
            \node[rectauxd, below=of orbit] (models) {External models};
            \node[rectind, below=of models] (ref_stack) {Reference SLC stack};
            \node[rectout, below=of geocoding] (enrb) {\acs{enrb}/\acs{eorb}};
    
            % Draw edges with arrows
            \draw[->, thick] (init) -- (coreg);
            \draw[->, thick] (coreg) -- (cslc);
            \draw[->, thick] (cslc) -- (enh);
            \draw[->, thick] (enh) -- (rtc);
            \draw[->, thick] (rtc) -- (geocoding);
    
            % Draw edges with arrows
            \draw[->, thick] (dem) -| (coreg);
            \draw[->, thick] (orbit) -| (coreg);
            \draw[->, thick] (models) -| (coreg);
            \draw[->, thick] (ref_stack) -| (coreg);
            \draw[->, thick] (geocoding) -- (enrb);
        \end{tikzpicture}
	}
	\caption{Sample workflow for generation of \acs{enrb} starting from \acs{cslc} products.}
	\label{fig:enrb_from_cslc}
\end{figure}

\subsection{Usability, Integrability and Interoperability}
\label{sec:enrb_usability_integrability}

The \acs{enrb} and \acs{eorb} products are improved versions of the \acs{nrb} and \acs{orb} products respectively. For this reason, Sections \autoref{sec:nrb_usability_integrability} for \acs{enrb} and \autoref{sec:orb_usability_integrability} for \acs{eorb} describe already the applications for these products, the usability properties as well as the integrability into different systems and workflows.

Applying techniques, like Ghost Removal and Sidelobe Suppression, significantly improves the quality and interpretability of radar imagery, particularly in high-resolution or complex environments, leading to cleaner and more reliable radar products, that enable better feature extraction, classification, and change detection in SAR applications.

In some of the analysed \acs{nrb} products the speckle filter is not applied. For specific analyses like land cover classification or vegetation monitoring, speckle noise mitigation may be required. This is an additional preprocessing step that one may need to consider.

Finally, radar backscatter can be influenced by surface roughness, moisture content, and other environmental factors. These must be accounted for when interpreting the data, especially in applications where surface characteristics (e.g., soil moisture, vegetation growth) vary over time.

\section{Geo-projected Ground Motion (GGM)}

\subsection{Product Description}
The Geo-projected Ground Motion (\acs{ggm}) product is particularly valuable for ground deformation applications. Its purpose is to derive three-dimensional displacement measurements by combining information from two (or more, when available) acquisition geometries. Unlike conventional SAR-based measurements expressed in the sensor reference frame (azimuth and slant range), the \acs{ggm} expresses motion in a geographic reference system, enabling more intuitive and physically meaningful analysis.

This geo-projected 3D displacement information is especially useful in a wide range of applications, such as glacier motion tracking, earthquake-induced surface deformation, subsidence monitoring, and other geophysical phenomena.

The definition of the \acs{ggm} product involves several critical aspects that must be carefully assessed to determine whether its implementation is justified. One feasible approach is to derive the \acs{ggm} by combining \acs{gpro} \& \acs{gpao} obtained from complementary acquisition geometries. In particular:

\begin{itemize}
    \item \acs{gpro} and \acs{gpao} extracted from phase information. In this case, the product \acs{gpro} coincides with \acs{gunw}. 
    \item \acs{gpro} and \acs{gpao} extracted from amplitude information. 
\end{itemize}
Both of the previous products provide displacement measurements, relative to a reference position, between two acquisition dates in the azimuth/slant range directions by using phase information. By using data from both ascending and descending geometries over the same geographic area, it is possible to compute the three components (north, east, vertical) of the displacement vector in a geographic coordinate frame.

The displacement captured by \acs{ggm} is the cumulative displacement measurements between two relatively close dates. \autoref{tab:ggm_product_characteristics} summarizes the different input configurations and the corresponding displacement components derivable from them.

\begin{longtblr}[
  label = tab:ggm_product_characteristics,
  caption = {Types of Measurements for Geo-projected Ground Motion},
]{
  width = \linewidth,
  colspec = {Q[l,25]Q[l,20]Q[l,20]Q[l,35]},
  hlines,
  hline{1-Z} = {-}{0.05em},
  rowhead=2,
}
\textbf{Input Geometries} & \textbf{GGM Type} & \textbf{Inputs} & \textbf{Measurement Components} \\
\SetCell[r=2]{l} Double geometry (ascending \& descending) & \acs{nev}-\acs{pt}\_GGM & \acs{gpro} \& \acs{gpao} by using amplitude & Along north, east, and vertical directions \\
 & \acs{nev}-IN\_GGM & \acs{gpro} \& \acs{gpao} by using phase & Along north, east, and vertical directions \\
\end{longtblr}

\subsubsection{Input and Auxiliary Data}
The input and auxiliary data  considered for product generation are listed in \autoref{tab:in_aux_ggm_products} (refer to \autoref{tab:ggm_product_characteristics} for \acs{ggm} product type).

\begin{longtblr}[
	label = tab:in_aux_ggm_products,
	entry = none,
	caption = {\acs{ggm} input and auxiliary data},
	]{
		width = \linewidth,
		colspec = {Q[l,20]Q[l,60]Q[c,20]},
		hlines,
		hline{1-Z} = {-}{0.05em},
	}
	\textbf{Name}   & \textbf{Description}   & \textbf{Mandatory (M) / Optional (O)} \\
	\acs{gpro}\_asc & Ground-projected Precise Range Offsets derived from ascending acquisitions by exploring phase or amplitude information & M \\
    \acs{gpro}\_desc & Ground-projected Precise Range Offsets derived from descending acquisitions by exploring phase or amplitude information & M \\
\end{longtblr}

The \acs{ggm} product is generated from input data derived using the same methodology. Therefore, it is not possible to combine pixel offset products obtained through amplitude cross-correlation with those derived from phase information, due to the significantly different sensitivity and accuracy of the input measurements.  

\subsubsection{Main Variables}
\autoref{tab:var_ggm_products} lists the main variables associated with the \acs{ggm} products.

\begin{longtblr}[
	label = tab:var_ggm_products,
	entry = none,
	caption = {\acs{ggm} main variables},
	]{
		width = \linewidth,
		colspec = {Q[l,15]Q[l,52]Q[l,20]Q[l,13]},
		hlines,
		hline{1-Z} = {-}{0.05em},
	}
	\textbf{Name} & \textbf{Description} & \textbf{Dimensions} & \textbf{Data Type} \\
    displacement
    
    \_north & Ground displacement along the north (Y) direction, in millimeters or centimeters. Present in \acs{nev} \acs{ggm} products & x, y & float32 \\
    displacement
    
    \_east & Ground displacement along the east (X) direction, in millimeters or centimeters. Present in \acs{nev} \acs{ggm} products & x, y & float32 \\
    displacement
    
    \_vertical & Ground displacement in the vertical (Z) direction, in millimeters or centimeters. Present in \acs{nev} \acs{ggm} products & x, y & float32 \\
\end{longtblr}

\autoref{tab:var_ggm_coordinates} describes the dimensions/coordinates associated with the \acs{ggm} main variables.

\begin{longtblr}[
    label = tab:var_ggm_coordinates,
    entry = none,
    caption = {\acs{ggm} main variables coordinates/dimensions},
    ]{
        width = \linewidth,
        colspec = {Q[l,20]Q[l,62]Q[l,18]},
        hlines,
        hline{1-Z} = {-}{0.05em},
    }
    \textbf{Name} & \textbf{Description} & \textbf{Data Type} \\
    x & UTM/UPS x (East) coordinates [m] & float64 \\
    y & UTM/UPS y (North) coordinates [m] & float64 \\
\end{longtblr}

\subsubsection{Main Metadata}
The following tables list the main metadata and per-pixel metadata associated with the \acs{ggm} products.

\begin{longtblr}[
	label = tab:meta_ggm_products,
	entry = none,
	caption = {\acs{ggm} main metadata and per-pixel metadata},
	]{
		width = \linewidth,
		colspec = {Q[l,17]Q[l,38]Q[l,20]Q[l,13]Q[c,12]},
		hlines,
		hline{1-Z} = {-}{0.05em},
	}
	\textbf{Name} & \textbf{Description} & \textbf{Dimensions} & \textbf{Data Type} & \textbf{Per-pixel Metadata}\\
    inputs\_IDs & List of the input products used fore generating the \acs{ggm} product & - & List of strings & \\
    date\_start & First or reference date of the interval to which the measurements refer & - & string & \\
    date\_end & Last date of the interval to which the measurements refer & - & string & \\
    reference\_point & ID of the reference point/cell of the measurements & - & string & \\
    quality\_mask & Mask indicating valid/incoherent areas & x, y & byte & \ding{51} \\
    uncertainty
    
    \_north & Uncertainty of the north displacement component. Present in \acs{nev} \acs{ggm} products. & x, y & float32 & \ding{51} \\
    uncertainty
    
    \_east & Uncertainty of the east displacement component. Present in \acs{nev} \acs{ggm} products. & x, y & float32 & \ding{51} \\
    uncertainty
    
    \_vertical & Uncertainty of the vertical displacement component. Present in \acs{nev} \acs{ggm} products. & x, y & float32 & \ding{51} \\
    coherence & Interferometric or amplitude coherence & x, y & float32 & \ding{51} \\
    x\_sample & Pixel spacing in x direction & - & float32 & \\
    y\_sample & Pixel spacing in y direction & - & float32 & \\
\end{longtblr}

The table below describes the dimensions/coordinates associated with the \acs{ggm} main metadata.

\begin{longtblr}[
    label = tab:var_ggm_metadata_coordinates,
    entry = none,
    caption = {\acs{ggm} main metadata coordinates/dimensions},
    ]{
        width = \linewidth,
        colspec = {Q[l,20]Q[l,62]Q[l,18]},
        hlines,
        hline{1-Z} = {-}{0.05em},
    }
    \textbf{Name} & \textbf{Description} & \textbf{Data Type} \\
    x & UTM/UPS x (East) coordinates [m] & float64 \\
    y & UTM/UPS y (North) coordinates [m] & float64 \\
\end{longtblr}

\subsubsection{Reference System Tiling and Slicing}
\autoref{tab:ggm_ref_system_slicing} reports the reference system, the slicing, and grid for \acs{ggm} products.

\begin{longtblr}[
    label=tab:ggm_ref_system_slicing,
    entry = none,
    caption = {Reference system and Slicing for \acs{ggm}},
    ]{
        width=\linewidth,
        colspec = {Q[l,15]Q[l,35]Q[l,35]Q[l,15]},
        %vline{2-5} = {-}{},
        %vline{5-5} = {-}{},
        hlines,
        hline{1-Z} = {-}{0.05em},
        rowhead=1
    }
    \textbf{Product} & \textbf{Reference System} & \textbf{Slicing} & \textbf{Spacing}\\
    \textbf{\acs{ggm}} & Cartographic: UTM/UPS & US-\acs{mgrs} 100 km x 100 km & 50 m \\
\end{longtblr}

\subsection{Product Requirements}
The main problem related to such a product is that more than one acquisition has to be combined and this implies that the optimal (or at least reasonable) couplings among a range of possibilities has to be found. Moreover, in general, it is not always possible to establish what is the “reasonable” coupling without taking into account what the specific use case is (i.e. the best coupling can change when we change the use case the \acs{ggm} has been thought for). One additional negative consequence directly related to the data coupling requirement is that \acs{ggm} production flow is much more complex than the ones designed for other \acs{ard} products. 

Although these considerations lead us to believe that these products are too complex for a routine generation process; we could think of a different approach in which we devolve to the end-user the burden of the data coupling (since is the user who knows the application the GGM will be used for), and we could release as open-source the code necessary for the \acs{ggm} production starting from the \acs{ard}s selected by the user.

\subsubsection{Accuracy}

Before assessing the accuracy performances for \acs{ggm} products, a series of consideration must be done:

The accuracy of Geoprojected Ground Motion (\acs{ggm}) products is primarily dictated by the directional sensitivity and noise characteristics of the input measurements:
\begin{itemize}
    \item \acs{los} (Line of Sight) phase measurements provide high precision and are predominantly sensitive to vertical and east-west deformation components, but exhibit limited sensitivity to north-south motion due to the near-polar orbit geometry. 
    \item In contrast, \acs{lof} (Line of Flight) offsets offer stronger sensitivity to the north-south component but are significantly noisier, typically by an order of magnitude or two, than range offsets, owing to lower coherence and reduced measurement precision.
\end{itemize}
In light of these constraints, the approach evaluated for \acs{ggm} products generation is a full 3D retrieval method that combines azimuth and range offsets from both ascending and descending acquisitions, yielding four independent measurements that jointly constrain the vertical, east, and north components of surface displacement. 

Accuracy assessment is conducted with consideration of the measurement geometries, noise levels, and the propagation of uncertainties through the inversion process. 

The following tables summarize the resulting estimated output accuracies for both \acs{s1} and \acs{rosel} missions. 

The following assumptions (see \autoref{tab:pro_pao_accuracy_phase_s1}) have been used to derive the performances of the \acs{s1} mission:  
\begin{itemize}
\item Range offset accuracy $\sigma_{\Delta^{ph}_{rg}} \simeq 0.005$ m 
\item Azimuth offset accuracy $\sigma_{\Delta^{amp}_{az}} \simeq 0.04$ m. 
\end{itemize}

\begin{longtblr}[
  label = tab:ggm_accuracy_s1,
  entry = none,
  caption = {Accuracy Requirements for \acs{ggm} products (\acs{s1} mission)},
]{
  width = \linewidth,
  colspec = {Q[l,35] Q[l,55] Q[c,20]},
  hlines,
  hline{1-Z} = {-}{0.05em},
  rowhead = 1,
}

\textbf{Metric} & \textbf{Description} & \textbf{Requirement} \\

$\sigma_n$ & North-South component accuracy & $\leq 0.03$ m \\

$\sigma_e$ & East-West component accuracy & $\leq 0.006$ m \\

$\sigma_u$ & Vertical component accuracy & $\leq 0.006$ m \\

\end{longtblr}

The following assumptions (see \autoref{tab:pro_pao_accuracy_phase_s1}) have been used to derive the performances for the \acs{rosel} mission:  
\begin{itemize}
\item Range offset accuracy $\sigma_{\Delta^{ph}_{rg}} \simeq 0.02$ m 
\item Azimuth offset accuracy $\sigma_{\Delta^{amp}_{az}} \simeq 0.12$ m. 
\end{itemize}

\begin{longtblr}[
  label = tab:ggm_accuracy_rosel,
  entry = none,
  caption = {Accuracy Requirements for \acs{ggm} products (\acs{rosel} mission)},
]{
  width = \linewidth,
  colspec = {Q[l,35] Q[l,55] Q[c,20]},
  hlines,
  hline{1-Z} = {-}{0.05em},
  rowhead = 1,
}

\textbf{Metric} & \textbf{Description} & \textbf{Requirement} \\

$\sigma_n$ & North-South component accuracy & $\leq 0.09$ m \\

$\sigma_e$ & East-West component accuracy & $\leq 0.025$ m \\

$\sigma_u$ & Vertical component accuracy & $\leq 0.025$ m \\

\end{longtblr}

\subsubsection{Timeliness}
This is a complex product derived from the combination of multiple \acs{ard} products. The input could derive from interferometric unwrapped phases \acs{gunw}, or precise azimuth and range offsets \acs{gpro} and \acs{gpao}. The combination of the different information does not significantly impact timeliness.

\subsection{Usability, Integrability and Interoperability}
The characteristics of the \acs{ggm} products in terms of usability, integrability and interoperability are in the following table.

\begin{longtblr}[
  caption = {Usability, Integrability, and Interoperability of the \acs{ggm} Product},
  label = {tab:ggm_usability_integrability},
]{
  colspec = {Q[2cm] Q[12cm]},
  rowhead = 1,
  width = \linewidth
}
\hline
\textbf{Aspect} & \textbf{Description} \\
\hline
Usability &
\begin{itemize}
  \item 3D Displacement Estimation (in case of multi-geometry derivation): Provides north-east-vertical(\acs{nev}) ground motion, which is easier to interpret than SAR-native (azimuth/slant-range or \acs{los}) measurements.
  \item Application-Ready: Useful for deformation studies related to earthquakes, landslides, subsidence, glacier flow, and infrastructure monitoring.
  \item User-Friendly Geometry: Expressed in geographic coordinates, facilitating integration into \acs{gis} and downstream geophysical analysis.
  \item Take into account that the product accuracy may depend on assumptions (e.g., negligible north-south motion in \acs{gunw}-only cases), and input data quality (e.g., coherence, ionospheric corrections).
\end{itemize}
\\
Integrability &
\begin{itemize}
  \item It Inherits integrability features from \acs{gpro}, \acs{gpao}, and \acs{gunw}, allowing seamless use in existing deformation workflows.
  \item Integration with \acs{gis} and Modeling Tools: Usable as input in hazard models, crustal and glacier deformation simulation or ground deformation assessments, in general.
  \item Suitable for Multi-Source Analysis: Can be fused with other data for validation or joint deformation analysis.
\end{itemize}
\\
Interoperability &
\begin{itemize}
  \item It Uses standard geographic grids and reference systems (e.g., WGS84, \acs{utm}), supporting integration across \acs{eo} systems.
  \item Cross-Mission Compatibility: It can be built from different SAR missions data.
  \item Format Compatibility: Delivered in common geospatial formats, compatible with \acs{eo} tools.
  \item Supports Interdisciplinary Use: Usable by geologists, civil engineers, glaciologists, and emergency response planners.
\end{itemize}
\\
\hline
\end{longtblr}


\section{Stripmap equivalent CSLC}
\label{sec:strip_eq_cslc}

\subsection{Product Description}
Managing SAR acquisitions in burst mode is more complex than in standard stripmap mode. Burst-mode acquisitions, such as those from TOPSAR or ScanSAR, introduce additional phase modulation due to the scanning nature of the acquisition. As a result, SAR products generated from burst-mode data cannot be directly mosaicked, since the overlapping regions between bursts exhibit phase discontinuities. To produce a coherent image, these effects must be corrected through processes such as resampling, interpolation, or sub-band extraction.

To reduce the complexity for the end user, each burst’s signal can be demodulated and the bursts subsequently mosaicked. This results in a single, continuous image that simulates a Stripmap acquisition.

The resulting Stripmap-equivalent (\acs{stripeq}) \acs{slc} or \acs{cslc} image is a full-resolution SAR product that preserves the resolution of the original input data. It seamlessly merges multiple bursts acquired in wide-swath modes like TOPSAR or ScanSAR into a unified image, maintaining the spatial consistency and phase continuity characteristic of Stripmap mode. The output retains the full information content of the original burst-mode data.

\subsubsection{Input and Auxiliary Data}
The following input and auxiliary data are considered for product generation:

\begin{longtblr}[
	label = tab:in_aux_stripeq_products,
	entry = none,
	caption = {\acs{stripeq} input and auxiliary data},
	]{
		width = \linewidth,
		colspec = {Q[l,20]Q[l,60]Q[c,20]},
		hlines,
		hline{1-Z} = {-}{0.05em},
	}
	\textbf{Name}   & \textbf{Description}   & \textbf{Mandatory (M) / Optional (O)} \\
    \acs{slc} or \acs{cslc} & L1 Single Look Complex image or co-registered L1 Single Look Complex image and annotations & M \\
    Precise Orbits & \acs{slc} precise orbits or \acs{cslc} precise reference orbits. & M \\
\end{longtblr}

\subsubsection{Main Variables}
The following tables list the main variables associated with the \acs{stripeq} products.

\begin{longtblr}[
	label = tab:var_stripeq_products,
	entry = none,
	caption = {\acs{stripeq} main variables},
	]{
		width = \linewidth,
		colspec = {Q[l,15]Q[l,52]Q[l,20]Q[l,13]},
		hlines,
		hline{1-Z} = {-}{0.05em},
	}
	\textbf{Name} & \textbf{Description} & \textbf{Dimensions} & \textbf{Data Type} \\
	\acs{slc}/\acs{cslc} & Single look complex or Co-registered Single look complex data: 
    
    - the amplitude is in \gls{beta0} convention, obtained by the calibration step using the Look-up-Table from the \acs{slc}/\acs{cslc} annotations;
    
    - the phase corresponds to the interpolated and resampled \acs{slc}/\acs{cslc} phase & polarization, acquisition\_date, azimuth\_time, range\_time & complex64 \\
    phase\_model & Topographic phase model, in rad, removed from the original image phase & acquisition\_date, azimuth\_time, range\_time & float32 \\
\end{longtblr}

The table below describes the dimensions/coordinates associated with the \acs{cslc} and \acs{gslc} main variables.

\begin{longtblr}[
	label = tab:var_stripeq_coordinates,
	entry = none,
	caption = {\acs{stripeq} main variables coordinates/dimensions},
	]{
		width = \linewidth,
		colspec = {Q[l,20]Q[l,62]Q[l,18]},
		hlines,
		hline{1-Z} = {-}{0.05em},
	}
	\textbf{Name} & \textbf{Description} & \textbf{Data Type} \\
	polarization & Acquired polarizaions & string \\
    acquisition\_date & Image acquisition dates, formatted as YYYY-MM-DD & string \\
    azimuth\_time & Time along satellite track direction & datetime64 \\
    range\_time & Time across satellite track direction & float64 \\
\end{longtblr}

\subsubsection{Main Metadata}
The following tables list the main metadata and per-pixel metadata associated with the \acs{stripeq} products.

\begin{longtblr}[
	label = tab:meta_stripeq_products,
	entry = none,
	caption = {\acs{stripeq} main metadata and per-pixel metadata},
	]{
		width = \linewidth,
		colspec = {Q[l,17]Q[l,38]Q[l,20]Q[l,13]Q[c,12]},
		hlines,
		hline{1-Z} = {-}{0.05em},
	}
	\textbf{Name} & \textbf{Description} & \textbf{Dimensions} & \textbf{Data Type} & \textbf{Per-pixel Metadata}\\
	input\_slc/cslc & Link to input \acs{slc}/\acs{cslc} & - & string & \\
    orbit\_files & Orbital file path list used for each of the images belonging to the \acs{cslc} stack & acquisition\_date & string & \\
    orbit\_coeffs & Polynomial coefficients describing the orbit associated with each input image belonging to the \acs{cslc} stack & acquisition\_date, axis, degree & float64 & \\
    orbit\_reference\
    
    \_time & Reference time associated with the orbit polynomial coefficients for each input image belonging to the \acs{cslc} stack & acquisition\_date & float64 & \\
    orbit\_type & Type of orbital data used for each of the images belonging to the \acs{cslc} stack: ETAD, POE, POD, RES & acquisition\_date & string & \\
    beta0 & LUT of the beta0 coefficients obtained from each input \acs{slc} annotation belonging to the \acs{cslc} stack & acquisition\_date, azimuth\_time, slant\_range\_time & float32 & \\
    gamma0 & LUT of the gamma0 coefficients obtained from each input \acs{slc} annotation belonging to the \acs{cslc} stack & acquisition\_date, azimuth\_time, slant\_range\_time & float32 & \\
    rcs & LUT of the Radar Cross Section (RCS) coefficients obtained from each input \acs{slc} annotation 
    belonging to the \acs{cslc} stack & acquisition\_date, azimuth\_time, slant\_range\_time & float32 & \\
    noise & LUT of thermal noise correction obtained from each input \acs{slc} annotation belonging to the \acs{cslc} stack & acquisition\_date, azimuth\_time, slant\_range\_time & float32 & \\
    phase\_deramp & The de-ramping phase (image or model) applied to each burst to obtain a continuous phase in the Stripmap Equivalent \acs{slc}/\acs{cslc}. & float32 & \\
    source\_data
    
    \_attributes & Dictionary containing the attributes associated with each image belonging to the \acs{cslc} stack & - & Dictionary & \\
    bursts\_list & A list of the input bursts & - & list of int32 & \\
    x\_sample & Pixel spacing in x direction & - & float32 & \\
    y\_sample & Pixel spacing in y direction & - & float32 & \\
\end{longtblr}

\subsubsection{Reference System Tiling and Slicing}
\label{sec:stripeq_reference_system_tiling}
\autoref{tab:stripeq_ref_system_slicing} reports the reference system, the slicing, and grid for the \acs{stripeq} products.


\begin{longtblr}[
	label=tab:stripeq_ref_system_slicing,
	entry = none,
	caption = {Reference system and Slicing for \acs{stripeq}},
	]{
		width=\linewidth,
		colspec = {Q[l,15]Q[l,40]Q[l,15]Q[l,30]},
		%vline{2-5} = {-}{},
		%vline{5-5} = {-}{},
		hlines,
		hline{1-Z} = {-}{0.05em},
		rowhead=1,
	}
	\textbf{Product} & \textbf{Reference System} & \textbf{Slicing} & \textbf{Spacing}\\
	\acs{stripeq} & Radar: azimuth/slant-range & Burst & Same of the input data \\   
\end{longtblr}

For the \acs{stripeq} product, it is planned to use the same reference radar grid and the same slicing for the ROSE-L and Sentinel-1 missions, to ensure interoperability between the products generated with the data of both missions.

\subsection{Product Requirements}

\subsubsection{Accuracy}

Given that the primary objective of this \acs{ard} product is to streamline the handling of complex burst-mode SAR acquisitions—such as those acquired in \acs{topsar} or ScanSAR modes—while preserving interoperability with standard \acs{stripeq} processing chains, it is critical to ensure that the product remains suitable for conventional amplitude- and phase-based analyses. These include, but are not limited to, amplitude image generation, interferogram computation, and coherence estimation.

Accordingly, the quality indicators have been defined to evaluate the continuity and consistency of derived products obtained from co-registered \acs{stripeq} image pairs. In particular, the assessment focuses on the continuity  of the interferometric phase and coherence maps generated from a pair of co-registered \acs{stripeq} images.

\begin{longtblr}[
	label = tab:stripeq_accuracy,
	entry = none,
	caption = {Accuracy Requirements for \acs{stripeq} product},
]{
	width = \linewidth,
	colspec = {Q[l,30] Q[l,60] Q[c,20]},
	hlines,
	hline{1-Z} = {-}{0.05em},
	rowhead = 1,
}

\textbf{Metric} & \textbf{Description} & \textbf{Requirement} \\

\textbf{Interferometric coherence discontinuity} & Maximum interferometric phase discontinuity between two adjacent bursts (only in the azimuth direction) & $\leq 0.01$ \\

\textbf{Interferometric phase discontinuity (deg)} & Maximum interferometric phase discontinuity between two adjacent bursts (in both range and azimuth directions) in case of absence of motion in the azimuth direction & $\leq 5$~deg \\

\end{longtblr}

\begin{comment}
\begin{itemize}
    \item IRF degradation wrt SLC (tool SCT) to be evaluated on few S1 \& SAOCOM data
    \item PSLR degradation wrt SLC (tool SCT) to be evaluated on few S1 \& SAOCOM data 
    \item ... degradation wrt SLC (tool SCT) to be evaluated on few S1 \& SAOCOM data
    \item geolocation accuracy (tool SCT) in correspondence of CRS
    \item Spectral information preservation (doppler difference wrt input data)
    \item Coregistration accuracy (rg, az) (theoretical formulas, CRS?) as a function of Ancillary data / processing
    \item phase continuity between bursts (theoretical analysis) ???? as a function of Ancillary data / processing
\end{itemize}    
\end{comment}

\subsubsection{Timeliness}
The timeliness is comparable to that of \acs{cslc} generation (see \textsection~\ref{sec:cslc_timeliness}). The de-ramping and de-bursting operations, necessary for generating this product, do not significantly impact timeliness.

\subsection{Usability, Integrability and Interoperability}
Usability and integrability of this product are the same as those collected for \acs{cslc} (\textsection~\ref{sec:usability_integrability_CSLC}), of which this novel product is an extension. It does not require processing related to burst fusion/merging.

